% ---- hand chunk 00-preamble.tex ----
%% =====================================================================
%%  Combinatorial Transseries and Their Inverses
%%  Consolidated companion volume of drafts/series-and-transseries/.
%%  Editorial merge (2026-09-04) of the three independently written
%%  articles filed on 2026-09-04 beside the canonical volume
%%  Transseries_And_Inversion; see the provenance chapter for the absorbed
%%  sources and what each contributed.  This file is the canonical source:
%%  the assembler that produced it read the three source directories,
%%  which were deleted after a residue audit, so edit here.
%%
%%  House style follows the canonical volume: A4, Libertinus with an
%%  lmodern fallback, book class with \part/\chapter.  Label prefixes:
%%  ct: (written for the merge), t1:, t2:, t3: (absorbed from the three
%%  sources, in the order recorded in the provenance chapter).
%% =====================================================================
\documentclass[11pt,openany]{book}
\usepackage[a4paper,margin=26mm,headheight=15pt]{geometry}
\usepackage[T1]{fontenc}
\usepackage[utf8]{inputenc}
\IfFileExists{libertinus.sty}{\usepackage{libertinus}}{\usepackage{lmodern}}
\usepackage{microtype}
\usepackage{amsmath,amssymb,amsthm,mathtools,bm}
% Canonical mathematical notation for active FabiusFunction LaTeX documents.
%
% This file is intentionally a plain TeX input rather than a LaTeX package.
% Every standalone document loads it by an explicit path relative to that
% document's directory; included fragments inherit the definitions from their
% root document.  Keep this file independent of document layout and metadata.
%
% Contract:
%   * one command has one mathematical meaning throughout the active corpus;
%   * names encode semantic roles rather than a document-local choice of letter;
%   * visually similar but differently normalized objects have different print;
%   * every command below is defined normatively in the notation catalogue.
%
% The active corpus excludes docs/archive/.  Historical sources may retain
% their original notation only inside an explicitly labelled quotation.

\makeatletter
\@ifundefined{mathbb}{%
  \PackageError{fabius-notation}{The amssymb package must be loaded first}%
    {Load amsmath and amssymb before inputting fabius-notation.tex.}%
}{}
\@ifundefined{operatorname}{%
  \PackageError{fabius-notation}{The amsmath package must be loaded first}%
    {Load amsmath before inputting fabius-notation.tex.}%
}{}
\makeatother

% Version stamp used by the catalogue and by static audits.
\newcommand{\FabiusNotationVersion}{2026-08-31}

% Number systems.  NaturalNumbers is explicitly the nonnegative convention.
\newcommand{\NaturalNumbers}{\mathbb{N}_{0}}
\newcommand{\PositiveIntegers}{\mathbb{N}_{>0}}
\newcommand{\IntegerNumbers}{\mathbb{Z}}
\newcommand{\RationalNumbers}{\mathbb{Q}}
\newcommand{\RealNumbers}{\mathbb{R}}
\newcommand{\ComplexNumbers}{\mathbb{C}}
\newcommand{\ExtendedRealNumbers}{\overline{\mathbb{R}}}

% The three principal Fabius--Rvachev functions.  Subscripts are deliberate:
% a bare F and a calligraphic F have several unrelated uses in the corpus.
\newcommand{\FabiusBounded}{F_{\mathrm{bd}}}
\newcommand{\FabiusGlobal}{F_{\mathrm{gl}}}
\newcommand{\FabiusQuantile}{Q_{F}}
\newcommand{\FabiusClampedQuantile}{Q_{F,\mathrm{cl}}}
\newcommand{\RvachevUp}{\operatorname{up}}

% Constants, elementary operators, and delimiters.
\newcommand{\EulerE}{\mathrm{e}}
\newcommand{\ImaginaryUnit}{\mathrm{i}}
\newcommand{\Differential}{\mathop{}\!\mathrm{d}}
\newcommand{\AbsoluteValue}[1]{\left\lvert #1\right\rvert}
\newcommand{\Norm}[1]{\left\lVert #1\right\rVert}
\newcommand{\Floor}[1]{\left\lfloor #1\right\rfloor}
\newcommand{\Ceiling}[1]{\left\lceil #1\right\rceil}
\newcommand{\FractionalPart}[1]{\left\{#1\right\}}
\newcommand{\PositivePart}[1]{\left(#1\right)_{+}}
\newcommand{\IndicatorOf}[1]{\mathbf{1}_{#1}}
\newcommand{\SetOf}[1]{\left\{#1\right\}}
\newcommand{\InnerProduct}[1]{\left\langle #1\right\rangle}
\newcommand{\CardinalityOf}[1]{\left\lvert #1\right\rvert}
\newcommand{\DefinitionEquals}{\mathrel{\coloneqq}}
\newcommand{\Sign}{\operatorname{sgn}}
\newcommand{\IdentityOperator}{\operatorname{id}}
\newcommand{\SpanOperator}{\operatorname{span}}
\newcommand{\TraceOperator}{\operatorname{tr}}
\newcommand{\RankOperator}{\operatorname{rank}}
\newcommand{\DiagonalOperator}{\operatorname{diag}}
\newcommand{\ResidueOperator}{\operatorname*{Res}}
\newcommand{\LeastCommonMultiple}{\operatorname{lcm}}
\newcommand{\OrderOperator}{\operatorname{ord}}
\newcommand{\PrincipalLogarithm}{\operatorname{Log}}
\newcommand{\FallingFactorial}[2]{\left(#1\right)^{\underline{#2}}}
\newcommand{\RisingFactorial}[2]{\left(#1\right)^{\overline{#2}}}

% Support, topology, convolution, and metric notation.
\newcommand{\SupportOperator}{\operatorname{supp}}
\newcommand{\SupportOf}[1]{\SupportOperator\!\left(#1\right)}
\newcommand{\TopologicalSupportOf}[1]{\operatorname{tsupp}\!\left(#1\right)}
\newcommand{\Convolution}{\mathbin{*}}
\newcommand{\MetricDistance}{\operatorname{dist}}
\newcommand{\TotalVariationDistance}{d_{\mathrm{TV}}}

% Probability.  Equality and convergence in law are not metric distance.
\newcommand{\Probability}{\mathbb{P}}
\newcommand{\Expectation}{\mathbb{E}}
\newcommand{\Variance}{\operatorname{Var}}
\newcommand{\Covariance}{\operatorname{Cov}}
\newcommand{\LawOperator}{\operatorname{Law}}
\newcommand{\LawOf}[1]{\LawOperator\!\left(#1\right)}
\newcommand{\EqualInLaw}{\mathrel{\overset{\mathrm{law}}{=}}}
\newcommand{\ConvergesInLaw}{\mathrel{\xrightarrow{\mathrm{law}}}}

% Fourier and integral transforms.  The printed labels expose normalization.
% FourierTwoPi uses exp(-2*pi*i*x*xi); FourierAngular uses exp(-i*x*omega).
\newcommand{\FourierTwoPi}[1]{\widehat{#1}_{\,2\pi}}
\newcommand{\FourierAngular}[1]{\widehat{#1}_{\,\mathrm{ang}}}
\newcommand{\LaplaceTransformSymbol}{\mathcal{L}}
\newcommand{\MellinTransformSymbol}{\mathcal{M}}
\newcommand{\LaplaceTransformOf}[1]{\LaplaceTransformSymbol\!\left[#1\right]}
\newcommand{\MellinTransformOf}[1]{\MellinTransformSymbol\!\left[#1\right]}
\newcommand{\MomentGeneratingFunction}[1]{M_{#1}}
\newcommand{\CharacteristicFunction}[1]{\varphi_{#1}}

% The two standard sinc functions must never share one symbol.
\newcommand{\SincRad}{\operatorname{sinc}_{\mathrm{rad}}}
\newcommand{\SincPi}{\operatorname{sinc}_{\pi}}

% Binary arithmetic and Thue--Morse notation.
\newcommand{\BinaryDigitSum}{\operatorname{wt}_{2}}
\newcommand{\ThueMorseSignSymbol}{\varepsilon}
\newcommand{\ThueMorseSign}[1]{\varepsilon_{#1}}
\newcommand{\PadicValuation}[1]{\nu_{#1}}
\newcommand{\TwoAdicValuation}{\nu_{2}}

% q-series notation.  Every base is an explicit argument.
\newcommand{\QPochhammer}[3]{\left(#1;#2\right)_{#3}}
\newcommand{\GaussianBinomial}[3]{%
  \genfrac{[}{]}{0pt}{}{#1}{#2}_{#3}%
}
\newcommand{\QInteger}[2]{\left[#1\right]_{#2}}

% Named special functions.
\newcommand{\Polylogarithm}{\operatorname{Li}}
\newcommand{\LambertW}{W}
\newcommand{\GammaFunction}{\Gamma}
\newcommand{\RiemannZeta}{\zeta}

% Asymptotics.  BigO and LittleO are operator tokens for existing formula
% grammar; prefer the At forms so that the limiting regime is printed.
\newcommand{\BigO}{\mathcal{O}}
\newcommand{\LittleO}{o}
\newcommand{\BigOAt}[2]{\mathcal{O}_{#2}\!\left(#1\right)}
\newcommand{\LittleOAt}[2]{o_{#2}\!\left(#1\right)}

\usepackage{booktabs,longtable,tabularx,array}
\usepackage{xcolor}
\usepackage{enumitem}
\usepackage{fancyhdr}
\usepackage{needspace}
\usepackage{xurl}
\usepackage{listings}
\usepackage{hyperref}
\usepackage[nameinlink,noabbrev,capitalise]{cleveref}

\definecolor{linkblue}{RGB}{24,72,124}
\definecolor{deepgreen}{RGB}{23,102,69}
\definecolor{shadegray}{RGB}{246,247,249}
\definecolor{rulegray}{RGB}{195,201,208}
\hypersetup{
  hypertexnames=false,
  colorlinks=true,linkcolor=linkblue,citecolor=deepgreen,urlcolor=linkblue,
  pdftitle={Combinatorial Transseries and Their Inverses},
  pdfauthor={Vladimir Reshetnikov; ProveIt formal-development synthesis},
  pdfsubject={Consolidated companion volume: all-order forward and inverse expansions of combinatorial counting functions, with exact exponentially small sectors},
  pdfkeywords={transseries, asymptotic inversion, Lambert W, Catalan numbers, Stirling numbers, Motzkin numbers, involutions, alternating permutations, connected graphs, necklaces, Gaussian binomial coefficients, finite fields, Galois numbers, theta functions}
}

\pagestyle{fancy}
\fancyhf{}
\fancyhead[L]{\small Combinatorial transseries and their inverses}
\fancyhead[R]{\small\nouppercase{\leftmark}}
\fancyfoot[C]{\thepage}
\renewcommand{\headrulewidth}{0.4pt}
\setcounter{tocdepth}{1}
\setcounter{secnumdepth}{2}
\setlist{itemsep=0.25em,topsep=0.4em}
\setlength{\parskip}{0.15em}
\setlength{\parindent}{1.1em}
\setlength{\emergencystretch}{2.5em}
\allowdisplaybreaks[2]
\numberwithin{equation}{chapter}

\theoremstyle{plain}
\newtheorem{theorem}{Theorem}[chapter]
\newtheorem{proposition}[theorem]{Proposition}
\newtheorem{lemma}[theorem]{Lemma}
\newtheorem{corollary}[theorem]{Corollary}
\theoremstyle{definition}
\newtheorem{definition}[theorem]{Definition}
\newtheorem{example}[theorem]{Example}
\theoremstyle{remark}
\newtheorem{remark}[theorem]{Remark}
\newtheorem{warning}[theorem]{Warning}
\crefname{warning}{Warning}{Warnings}
\Crefname{warning}{Warning}{Warnings}

\lstset{basicstyle=\ttfamily\footnotesize,columns=fullflexible,keepspaces=true,
  breaklines=true,showstringspaces=false,frame=single,rulecolor=\color{rulegray},
  backgroundcolor=\color{shadegray}}
\newcolumntype{P}[1]{>{\raggedright\arraybackslash}p{#1}}

% ---------------------------------------------------------------------
%  Macros.  One definition per printed meaning; the three sources' letters
%  are reconciled in the notation chapter.  Source-specific spellings that
%  survive in absorbed text are aliases of these.
% ---------------------------------------------------------------------
\newcommand{\ee}{\mathrm e}
\newcommand{\ii}{\mathrm i}
\newcommand{\dd}{\,\mathrm d}                 % differential with thin space (T2)
\newcommand{\ddn}{\mathrm d}                  % bare differential (T1, T3)
\newcommand{\NN}{\mathbb N}\newcommand{\RR}{\mathbb R}\newcommand{\CC}{\mathbb C}\newcommand{\QQ}{\mathbb Q}
\newcommand{\N}{\mathbb N}\newcommand{\R}{\mathbb R}\newcommand{\C}{\mathbb C}\newcommand{\Z}{\mathbb Z}
\newcommand{\F}{\mathbb F}
\newcommand{\OO}{\mathcal O}\newcommand{\bigO}{\mathcal O}
\newcommand{\coefA}[1]{\left[#1\right]}       % T1: [t^n] with one argument
\newcommand{\coefB}[2]{[#1^{#2}]}             % T3: [t^n] with two arguments
\newcommand{\coeff}[2]{\left[#1\right]#2}     % T2
\newcommand{\ExpC}{\mathcal E}\newcommand{\LogC}{\mathcal L}
\newcommand{\cE}{\mathcal E}\newcommand{\cL}{\mathcal L}
\newcommand{\fall}[2]{#1^{\underline{#2}}}
\newcommand{\rising}[2]{(#1)_{#2}}
\newcommand{\qbin}[3]{\genfrac{[}{]}{0pt}{}{#1}{#2}_{#3}}   % three arguments: base explicit
\newcommand{\qbinq}[2]{\genfrac{[}{]}{0pt}{}{#1}{#2}_{q}}   % T3: base q
\newcommand{\qbinQ}[2]{\genfrac{[}{]}{0pt}{}{#1}{#2}_{Q}}
\newcommand{\Sii}[2]{\genfrac{\{}{\}}{0pt}{}{#1}{#2}}
\newcommand{\StirII}[2]{\left\{\begin{matrix}#1\\#2\end{matrix}\right\}}
\newcommand{\StirI}[2]{\left[\begin{matrix}#1\\#2\end{matrix}\right]}
\newcommand{\sone}[2]{s(#1,#2)}
\newcommand{\baseid}[1]{\texttt{\detokenize{#1}}}
\newcommand{\sectorsim}{\mathrel{\sim_{\mathrm{sec}}}}
\newcommand{\pp}[1]{(Q;Q)_{#1}}
\newcommand{\Pinf}{\mathcal P_Q}
\newcommand{\A}{\mathcal A_Q}
\newcommand{\eps}{\varepsilon}
\newcommand{\T}{\Theta}
\newcommand{\one}{\mathbf 1}
\newcommand{\GL}{\operatorname{GL}}\newcommand{\Sp}{\operatorname{Sp}}
\newcommand{\Li}{\operatorname{Li}}\newcommand{\rad}{\operatorname{rad}}
\DeclareMathOperator{\sech}{sech}
\DeclareMathOperator{\Res}{Res}
\DeclareMathOperator{\ord}{ord}
\DeclareMathOperator{\Rea}{Re}
\newenvironment{abstract}{\chapter*{Abstract}\addcontentsline{toc}{chapter}{Abstract}}{}

\begin{document}
\frontmatter
\title{\textbf{Combinatorial Transseries\\and Their Inverses}\\[0.5em]
  \Large Gamma quotients, finite exponential sums, moment sectors, moving saddles,\\
  arithmetic sheets, and $q$-products}
\author{Vladimir Reshetnikov\\\small ProveIt formal-development synthesis}
\date{4 September 2026}
\maketitle

\begin{abstract}
The asymptotic inversion of a combinatorial counting sequence requires more
than reversing its leading term: one must identify the correct asymptotic
scale, choose or deliberately avoid a continuous interpolation, retain
subdominant contributions when they are exact, and control the amplification
of forward errors by inversion.  This volume develops these points, with the
methods of the canonical volume \emph{Transseries: the polynomial--logarithmic
calculus, series reversal at infinity, and the inversion of rapidly growing
functions}, through families that volume does not treat: balanced gamma
quotients (Catalan and Fuss--Catalan numbers, central multinomial
coefficients, rectangular tableaux); finite exponential sums (fixed-column
Stirling and Eulerian numbers) and fixed-cycle Stirling numbers; endpoint
moment sequences (Motzkin and central trinomial, Delannoy, large Schr\"oder);
two-saddle Gaussian moments (involutions); prime- and odd-integer-generated
sectors (alternating permutations); superexponential labelled classes
(connected graphs); divisor sums (necklaces, Lyndon words, irreducible
polynomials); slow growth and finite limits (harmonic functions); and the
exact $q$-products of finite-field enumeration (general linear, symplectic and
unitary group orders, flags, fixed-rank and scaled Gaussian binomials,
$q$-Catalan numbers, and Galois numbers with their root-lattice theta
sectors).

Three inversion engines are proved once and reused: phase-coordinate
inversion of an exactly invertible dominant phase with a closed formula for
every multiplicative-sector coefficient through signed Stirling numbers; a
scaled Lagrange--B\"urmann reversion for corrections that are unbounded but
small relative to the core; and a convergent inverse transseries theorem for
an analytic exponential tail attached to a quadratic, linear, or logarithmic
phase, with a finite nonrecursive coefficient formula and a geometric
remainder.  Every family is supplied with an exact interpolation, an
all-order forward expansion with explicit coefficients, an all-order inverse
construction, a residual certificate, and an integer-threshold rule.
Convergent identities, Poincar\'e expansions with remainders, exactly
normalized exponentially small sectors, formal range inverses, and finite
arithmetic cutoffs are kept separate throughout.  Three verification programs
with their recorded runs accompany the volume.  No complete complex Stokes
classification and no Lean formalization of the new results is claimed.
\end{abstract}

\clearpage
\tableofcontents
\mainmatter

% ---- hand chunk 01-ch1-head.tex ----
\part{Foundations}
\label{ct:part:foundations}

\chapter{What is being inverted}
\label{ct:ch:objects}

The canonical volume of this group \cite{tai} develops the
polynomial--logarithmic calculus, series reversal at infinity, and an
apparatus for inverting a rapidly growing function organized around an
exactly invertible dominant phase and a nonvanishing normalized slope; its
combinatorial case studies are rooted trees, the double and swing
factorials, the partition numbers, Bell and Fubini numbers, and several
special functions.  This volume keeps that architecture and applies it to
families the canonical volume does not treat.  Its interfaces are the
canonical volume's perturbed-inversion theorem, its centred Lambert
reversion, its backward-error certificate, and its staircase theorem for
integer thresholds (the results labelled \baseid{p0:thm:perturbed-inversion},
\baseid{p0:thm:lambert-centered}, \baseid{p0:thm:backward-error},
\baseid{p0:thm:staircase} in its source).  Everything needed below is
nevertheless stated and proved here, so that this volume can be read on its
own.

Before any family is treated, three preliminary matters must be settled:
what an ``inverse'' of a sequence is, why the choice of a continuous
interpolation is a mathematical datum rather than a convenience, and what
the word ``transseries'' will mean.  The three absorbed sources each made
these points in their own words; this chapter makes them once.

\section{Four logically different objects}
\label{ct:sec:four-objects}

% ---- T2 lines 107--122 ----
\subsection{Four logically different objects}

Let $(a_n)$ be eventually positive and strictly increasing. The notation ``inverse of $a_n$'' can refer to four different objects.

\begin{definition}
The \emph{inverse on the range} sends $a_n$ to $n$. The \emph{threshold inverse} is the staircase
\begin{equation}
 N(y)=\min\{n\ge n_0:a_n\ge y\}.
 \label{t2:eq:staircase}
\end{equation}
A \emph{specified continuous inverse} is $F^{-1}$ for an explicitly given, eventually increasing function satisfying $F(n)=a_n$. A \emph{formal asymptotic inverse} is a formal expression whose substitution reverses the chosen forward asymptotic model to every prescribed order.
\end{definition}

The four notions often cooperate, but they are not interchangeable. A formal inverse may approximate the range inverse without specifying values of $F$ between integers. An interpolation may be natural from a gamma quotient, a moment integral, or a finite exponential sum. Nevertheless, ``natural'' is a choice, not a uniqueness theorem.

Throughout, $F^{-1}$ means a compositional inverse, never the reciprocal $1/F$. Unless explicitly stated otherwise, $x\to+\infty$, the parameters such as a tableau height or a Stirling column are fixed, and logarithms are real natural logarithms. All assertions of eventual positivity or monotonicity allow discarding a finite initial interval.
% ---- hand chunk 01-ch1-gauge.tex ----

Throughout, $F^{-1}$ means a compositional inverse of the growth
interpolation, never the reciprocal $1/F$ and never the compositional
inverse of an ordinary generating function.  Unless stated otherwise
$x\to+\infty$, structural parameters such as a tableau height, a Stirling
column, or the base $q$ are fixed, and logarithms are real natural
logarithms.  Assertions of eventual positivity or monotonicity permit
discarding a finite initial interval.

\section{Interpolation is a mathematical datum}

% ---- T2 lines 125--162 ----
\label{t2:subsec:gauge}

\begin{proposition}[Interpolation gauge freedom]
\label{t2:prop:gauge}
Let $F$ be an eventually strictly increasing $C^1$ function. There exist other eventually strictly increasing $C^1$ functions with the same values at every sufficiently large integer whose inverse differs from $F^{-1}$ by a nonvanishing bounded oscillation. There also exist such functions whose inverse differs first at any prescribed inverse-power order, or at an exponentially small order while preserving all inverse-power coefficients.
\end{proposition}

\begin{proof}
For $|\eta|<1/(2\pi)$, define
\[
 h(x)=x+\eta\sin(2\pi x),\qquad \widetilde F=F\circ h.
\]
Then $h(n)=n$ and $h'(x)=1+2\pi\eta\cos(2\pi x)>0$. Thus $\widetilde F(n)=F(n)$ and
\[
 \widetilde F^{-1}=h^{-1}\circ F^{-1}.
\]
The map $h^{-1}(u)-u$ is nonzero and periodic if $\eta\ne0$, so the inverse difference need not even tend to zero.

For the finer assertions replace the sine term by $\eta x^{-r}\sin(2\pi x)$, $r>0$, or by $\eta\ee^{-\rho x}\sin(2\pi x)$, $\rho>0$. In both cases $h'(x)>0$ eventually. Direct substitution in $h(h^{-1}(u))=u$ gives, respectively,
\[
 h^{-1}(u)=u-\eta u^{-r}\sin(2\pi u)+o(u^{-r}),
\]
and
\[
 h^{-1}(u)=u-\eta\ee^{-\rho u}\sin(2\pi u)
           +\bigO(\ee^{-2\rho u}).
\]
The latter change is smaller than every inverse power. All the integer samples are nevertheless unchanged.
\end{proof}

Consequently, interpolation-independence of an algebraic inverse expansion holds only within a class of \emph{asymptotically compatible} interpolations. It does not follow from monotonicity and agreement on the integers. This qualification becomes indispensable when exponentially small corrections are discussed.

If $F$ is an exact increasing interpolation, then, on its eventual range,
\begin{equation}
 N(y)=\left\lceil F^{-1}(y)\right\rceil.
 \label{t2:eq:ceil}
\end{equation}
An approximation $\widehat x$ with $|\widehat x-F^{-1}(y)|\le\varepsilon$ determines this ceiling only when the two endpoints $\widehat x\pm\varepsilon$ have the same ceiling. At a sample value $y=a_n$, an inverse approximation with error less than $1/2$ recovers $n$ by nearest-integer rounding. An unqualified ceiling of an $o(1)$ approximation is not a valid threshold algorithm near a jump.
% ---- T3 lines 128--145 ----
\begin{proposition}[Interpolation can change a flat inverse term]\label{t3:prop:ambiguity}
Suppose $F>0$ is an eventually increasing smooth interpolation, and its logarithmic derivative is eventually bounded below by a positive constant. For fixed $\sigma>0$ and sufficiently small real $\eta$, set
\[
 \widetilde F(x)=F(x)\exp\!\bigl(\eta\ee^{-\sigma x}\sin\pi x\bigr).
\]
Then $\widetilde F(n)=F(n)$ for every integer $n$, and $\widetilde F$ is eventually increasing. If $x=F^{-1}(Y)$, then, locally as $\eta\to0$,
\begin{equation}\label{t3:eq:interpolation-shift}
 \widetilde F^{-1}(Y)-x
 =-\eta\frac{\ee^{-\sigma x}\sin\pi x}{(\log F)'(x)}+O(\eta^2).
\end{equation}
\end{proposition}
\begin{proof}
The sine factor vanishes at integers. The derivative of the added logarithmic term is $O(\ee^{-\sigma x})$, so eventual positivity of the original logarithmic derivative is preserved. Differentiate
$\log F(x(\eta))+\eta\ee^{-\sigma x(\eta)}\sin\pi x(\eta)=\log Y$
with respect to $\eta$ at zero. The implicit function theorem gives the stated expansion.
\end{proof}

In the product examples we therefore \emph{define} the interpolation by a ratio of convergent infinite products. In the theta examples we define one analytic function for each residue class. In the arithmetic examples we distinguish exact integer formulas from explicitly chosen functions on fixed-radical subsequences.
% ---- hand chunk 01-ch1-meaning.tex ----

The two propositions make the same point in additive and in multiplicative
form: the integer samples of a counting sequence leave every exponentially
small inverse term, and even a bounded oscillation of the inverse,
undetermined until an interpolation has been named.  Every analytic inverse
assertion in this volume therefore names the function being inverted --- a
gamma quotient, a moment integral, a finite exponential sum, a ratio of
convergent infinite products, a theta-weighted lattice sum --- and every
statement about a sequence alone is a statement about its range or its
threshold.

\section{What ``transseries'' will mean here}
\label{ct:sec:meaning}

% ---- T2 lines 166--179 ----
We distinguish an inverse-power expansion
\[
 F(x)\sim F_0(x)\sum_{j\ge0}a_jx^{-j}
\]
from a sector expansion such as
\[
 F(x)\sectorsim F_0(x)\sum_{k\ge0}\ee^{-k\omega x}
       \sum_{j\ge0}a_{k,j}x^{-j}.
\]
The latter notation means that the separate sector amplitudes have the stated expansions in a specified exact decomposition or a specified formal model. It does \emph{not} mean that truncating the first sector after ten inverse powers leaves an exponentially small error.

A finite additive semigroup of positive actions is locally finite below each action cutoff. This makes coefficient extraction at a fixed exponential order finite. Prime-generated actions remain locally finite because there are finitely many integers below a fixed bound. Necklace actions, in contrast, accumulate at a finite value; Section~\ref{ct:ch:necklaces} explains the resulting limitation.

Factors such as $\cos(\pi x)$ are bounded oscillatory coefficients. They are not positive monomials of an ordinary real logarithmic--exponential Hardy field. Whenever they occur, we work with the explicitly specified real oscillatory extension, or equivalently with conjugate complex exponential sectors. Calling this distinction out avoids assigning an oscillating term a nonexistent eventual sign.
% ---- T1 lines 121--127 ----
A \emph{block transseries} here consists of explicitly defined functions $A_\nu(x)$, carrying specified exponential or oscillatory factors, together with all-order expansions of their amplitudes. Whenever we write an exact identity such as
\[
 F(x)=A_+(x)+\cos(\pi x)A_-(x),
\]
the blocks are fixed by integrals or products, not merely by their asymptotic coefficients. This is essential: a finite algebraic truncation of $A_+$ usually has an error much larger than $A_-$. Exponentially small inverse sectors are first defined relative to the \emph{exact} dominant block and only then expanded internally.

Oscillatory terms are ordered by their nonoscillatory envelopes. A statement such as $\OO(3^{-2x})$ is meaningful at zeros of $\cos\pi x$; a relative equivalence to $3^{-x}\cos\pi x$ need not be. We use complex harmonics $\ee^{\pm\ii\pi x}$ for coefficient bookkeeping and take real combinations at the end.
% ---- T3 lines 147--166 ----
\subsection{Convergent exponential expansions are transseries too}

Consider
\begin{equation}\label{t3:eq:basic-shape}
 F(x)=\exp\bigl(P(x)+B(\ee^{-\lambda x})\bigr),\qquad
 B(t)=\sum_{m\ge1}b_m t^m,
\end{equation}
where $B$ is analytic at zero and $\lambda>0$. Its relative expansion is
\[
 F(x)=\ee^{P(x)}\sum_{m\ge0}v_m\ee^{-m\lambda x},\qquad v_0=1.
\]
For every fixed $m$, the factor $\ee^{-m\lambda x}$ is smaller than every negative power of $x$. Thus these are exponentially separated transseries sectors even when the series in the small variable $t$ converges. There is no requirement that a transseries be divergent.

The inverse will be written
\begin{equation}\label{t3:eq:inverse-shape-intro}
 F^{-1}(Y)=s+\sum_{m\ge1}d_m(s)\ee^{-m\lambda s},\qquad
 P(s)=\log Y.
\end{equation}
For a quadratic $P$, the $d_m$ are finite polynomials in reciprocal powers of $P'(s)$, not merely formal infinite expansions inside each sector. Keeping the full center $s$ is important: a missing constant in $P$ creates an error of order $1/P'(s)$, which is much larger than any of the exponential corrections.

% ---- hand chunk 01-ch1-rounding.tex ----

\section{Rounding belongs after inversion}
\label{ct:sec:rounding}

For an increasing interpolation valid on all sufficiently large integers,
the threshold \cref{t2:eq:staircase} is the ceiling \cref{t2:eq:ceil} of the
real inverse, and an approximation $\widehat x$ with error at most $\eps$
determines that ceiling only when $\widehat x\pm\eps$ have the same ceiling.
Several families below are interpolated on residue classes rather than on
all integers.

% ---- T3 lines 173--177 ----
For a sheet matching $n\equiv\eps\pmod r$, the corresponding threshold is
\begin{equation}\label{t3:eq:residue-rounding}
 N_{\eps}(Y)=\eps+r\left\lceil\frac{F_{\eps}^{-1}(Y)-\eps}{r}\right\rceil.
\end{equation}
The overall threshold is the minimum of these sheet thresholds, after any finite initial range has been checked. A tiny real error does not automatically certify a ceiling when the true inverse is near an integer. Section~\ref{t3:sec:certification} gives residual and rounding criteria.
% ---- hand chunk 02-ch2-head.tex ----

\chapter{Coefficient calculus and the inversion engines}
\label{ct:ch:engines}

Every family in this volume is treated by the same three-step procedure:
an exact representation of the counting function is chosen; its dominant
logarithmic phase is isolated and inverted exactly in closed form; and the
remaining correction is transported through the inverse by a finite
coefficient formula.  This chapter proves the tools once.  The finite
partition formulas for exponentials and logarithms of power series make
every forward coefficient a finite polynomial; the phase-coordinate
inversion theorem and its multiplicative-sector corollary give every
inverse sector coefficient as a finite differential expression; the scaled
Lagrange--B\"urmann reversion handles corrections that are unbounded but
small relative to the core; the balanced linear--logarithmic core is
inverted to all orders; the convergent inverse transseries theorem for an
analytic exponential tail attached to a quadratic, linear, or logarithmic
phase supplies both a finite nonrecursive coefficient formula and a
geometric remainder; and a residual certificate turns forward residuals
into inverse error bounds.

\section{Finite partition formulas}
\label{ct:sec:partition-formulas}

% ---- T2 lines 186--207 ----
For a formal series $Q(t)=\sum_{j\ge1}q_jt^j$, set
\begin{equation}
 \cE_n(q)=
 \sum_{\substack{m_1,\ldots,m_n\ge0\\\sum j m_j=n}}
       \prod_{j=1}^n\frac{q_j^{m_j}}{m_j!},
 \qquad \cE_0(q)=1.
 \label{t2:eq:E}
\end{equation}
Then $\exp Q(t)=\sum_{n\ge0}\cE_n(q)t^n$. Thus every coefficient is a finite polynomial, even when the power series itself diverges. For example,
\[
 \cE_1=q_1,\quad
 \cE_2=q_2+\frac{q_1^2}{2},\quad
 \cE_3=q_3+q_1q_2+\frac{q_1^3}{6}.
\]
Conversely, if $A(t)=1+\sum_{j\ge1}a_jt^j$, then
\begin{equation}
 \cL_n(a):=[t^n]\log A(t)
 =\sum_{\substack{\sum j m_j=n\\M=\sum m_j}}
   (-1)^{M+1}(M-1)!\prod_{j=1}^n\frac{a_j^{m_j}}{m_j!}.
 \label{t2:eq:L}
\end{equation}
Both formulas follow by expanding the exponential or logarithm and applying the multinomial theorem. They are the ordinary-coefficient form of the Bell-polynomial calculus used in the canonical volume \cite{tai}.
% ---- T1 lines 153--160 ----
We write $s(r,m)$ for the \emph{signed} Stirling numbers of the first kind, normalized by
\begin{equation}
 \fall{z}{r}=\sum_{m=0}^r s(r,m)z^m,
 \qquad
 (\log(1+u))^m=m!\sum_{r\ge m}\frac{s(r,m)}{r!}u^r.
 \label{t1:eq:stirlingsign}
\end{equation}
Thus $s(2,1)=-1$. A sign error here changes the inverse sectors.
% ---- T2 lines 209--241 ----
\begin{lemma}[Gamma-ratio coefficients]
\label{t2:lem:ratio}
For fixed $a,b$, with $x$ positive and large,
\begin{equation}
 \frac{\Gamma(x+a)}{\Gamma(x+b)}
 \sim x^{a-b}\sum_{n\ge0}R_n(a,b)x^{-n},
 \qquad R_n(a,b)=\cE_n(r(a,b)),
 \label{t2:eq:ratio}
\end{equation}
where
\begin{equation}
 r_j(a,b)=\frac{(-1)^{j+1}}{j(j+1)}
          \bigl(B_{j+1}(a)-B_{j+1}(b)\bigr).
 \label{t2:eq:ratio-log}
\end{equation}
Here $B_m(z)$ denotes the Bernoulli polynomial.
\end{lemma}
\begin{proof}
The shifted Stirling expansion gives
\[
 \log\Gamma(x+a)-\log\Gamma(x+b)
 =(a-b)\log x+\sum_{j=1}^{N}r_j(a,b)x^{-j}
       +\bigO(x^{-N-1}).
\]
The expansion is uniform for $a,b$ in a fixed compact set; see the standard gamma asymptotics and Bernoulli-polynomial conventions in \cite{dlmfgamma,dlmfbernoulli}. Exponentiating a finite truncation and using \eqref{t2:eq:E} proves the formula with a remainder of the indicated order. Since $N$ is arbitrary, it proves the all-orders statement. This use of a divergent formal series does not assert convergence.
\end{proof}

For reference, the Bernoulli convention is fixed by
\[
 \frac{t\ee^{zt}}{\ee^t-1}=\sum_{m\ge0}B_m(z)\frac{t^m}{m!},
 \qquad B_1(z)=z-\tfrac12.
\]
This convention is especially important for the shifted half-integer formulas later in the volume.
% ---- hand chunk 02-ch2-phase.tex ----

The two partition formulas are written with the same letters in every
chapter: $\ExpC_n=\cE_n$ for the exponential and $\LogC_n=\cL_n$ for the
logarithm.

\section{Additive perturbations of an exactly invertible phase}
\label{ct:sec:phase}

% ---- T1 lines 163--240 ----
\begin{theorem}[Phase-coordinate inversion]
\label{t1:thm:phase}
Let $\phi$ and $R$ be analytic near $x_0$, with $\phi'(x_0)\ne0$, let $L=\phi(x_0)$, and write $g=\phi^{-1}$ locally. The solution of
\[
 \phi(x)+zR(x)=L,
\]
which equals $x_0$ at $z=0$, is analytic for sufficiently small $z$ and satisfies
\begin{equation}
 x=g(L)+\sum_{m\ge1}\frac{(-z)^m}{m!}
 \frac{\ddn^{m-1}}{\ddn L^{m-1}}
 \left\{g'(L)R(g(L))^m\right\}.
 \label{t1:eq:phaseinv}
\end{equation}
The same identity holds formally whenever the perturbation is in a complete positive-order ideal and the phase slope is invertible.
\end{theorem}
\begin{proof}
Set $v=\phi(x)$. Then $v=L-zR(g(v))$. The Lagrange--B\"urmann formula for $g(v)$ gives
\[
 g(v)=g(L)+\sum_{m\ge1}\frac{(-z)^m}{m!}
 \partial_L^{m-1}\bigl(g'(L)R(g(L))^m\bigr).
\]
Analytic existence follows from the implicit function theorem at $(x,z)=(x_0,0)$. In the formal setting, comparing successive orders determines each new coefficient by a single division by $\phi'(x_0)$; the analytic coefficient identity is a polynomial identity after that division and therefore also proves the formal formula.
\end{proof}

This is the phase-coordinate specialization of the canonical volume's perturbed inversion theorem, not a new version of Lagrange inversion. The useful extension is the following closed formula for \emph{multiplicative} sectors.

\begin{theorem}[All multiplicative-sector coefficients]
\label{t1:thm:multi}
Let $A(x)>0$ on a real tail, let $\phi=\log A$, assume its derivative is nonzero at the chosen base point, and put $g=\phi^{-1}$ locally. For finitely many analytic functions $q_j$, consider
\begin{equation}
 F_{\bm z}(x)=A(x)\left(1+\sum_{j=1}^d z_jq_j(x)\right).
 \label{t1:eq:multifamily}
\end{equation}
Write $L=\log y$, $D=\ddn/\ddn L$, and, for a multi-index $\nu\ne0$, write $r=|\nu|$, $\nu!=\prod_j\nu_j!$, and
\[
 Q_\nu(L)=\prod_j q_j(g(L))^{\nu_j}.
\]
Then the inverse branch at $\bm z=0$ is
\begin{equation}
 \boxed{\quad
 F_{\bm z}^{-1}(y)=g(L)+
 \sum_{\nu\ne0}\frac{\bm z^\nu}{\nu!}
 \sum_{m=1}^{|\nu|}(-1)^m s(|\nu|,m)
 D^{m-1}\!\left(g'(L)Q_\nu(L)\right).
 \quad}
 \label{t1:eq:multiinv}
\end{equation}
It is locally convergent in the sector markers and is also a formal identity. Each coefficient is a finite differential expression.
\end{theorem}
\begin{proof}
Apply \cref{t1:thm:phase} to the perturbation
$R_{\bm z}(x)=\log(1+\sum z_jq_j(x))$, temporarily introducing a second scalar marker for this perturbation. At a fixed degree $r$ in $\bm z$, only terms $m\le r$ contribute. By \eqref{t1:eq:stirlingsign},
\[
 \coefA{\bm z^\nu}R_{\bm z}(x)^m
 =\frac{m!s(r,m)}{\nu!}\prod_jq_j(x)^{\nu_j}.
\]
Substitution in \eqref{t1:eq:phaseinv} gives \eqref{t1:eq:multiinv}. This coefficientwise derivation is legitimate even when one has not established convergence at the final marker values. Analytic convergence near zero follows independently from the implicit function theorem.
\end{proof}

For a single correction $q$, the first two orders are
\begin{equation}
 x=g-g'q+\frac12\left\{g'q^2+D(g'q^2)\right\}+\OO(q^3)
 \label{t1:eq:firstmulti}
\end{equation}
when the derivatives of $q$ have the same envelope order. Here $q$ is evaluated at $g$. In the $x$ coordinate, the first displacement is
\begin{equation}
 -\frac{q(x_0)}{\phi'(x_0)},\qquad x_0=A^{-1}(y).
 \label{t1:eq:firstdisplacement}
\end{equation}

\begin{remark}[Convergence at the physical marker]
A local theorem at $\bm z=0$ does not by itself justify setting every $z_j=1$. In the applications below, the corrections are uniformly small on a fixed complex disk about a large real $x_0$, while $\phi'$ is asymptotic either to a positive constant, to $c\log x_0$, or to $cx_0$. Taylor's formula gives the fixed-point equation
\[
 \delta=-\frac{\log(1+\sum z_jq_j(x_0+\delta))
 +\phi(x_0+\delta)-\phi(x_0)-\phi'(x_0)\delta}{\phi'(x_0)}.
\]
On a sufficiently small fixed disk it is a contraction, uniformly for $|z_j|\le2$, once $x_0$ is large. Its analytic fixed point justifies the marker expansion at $z_j=1$. This also gives a geometric remainder in marker degree, with the correction envelope as small parameter. Infinite Dirichlet sums will be treated by absolute local uniform convergence.
\end{remark}
% ---- hand chunk 02-ch2-transport.tex ----

\begin{remark}[Sector transport of one exact small correction]
\label{ct:rmk:transport}
The single-correction case of \cref{t1:thm:multi} is the form in which the
theorem is used for the endpoint, saddle, and arithmetic families.  Let
$F_0>0$ be smooth on an eventual interval with $(\log F_0)'>0$, let
$\phi_0=\log F_0$ and $g=\phi_0^{-1}$, and let $F=F_0(1+\varepsilon)$ with
$1+\varepsilon>0$ and $\varepsilon$ graded by positive sector weights that
differentiation preserves.  With $h(L)=\log(1+\varepsilon(g(L)))$, the formal
inverse graded by the sectors of $\varepsilon$ is
\begin{equation}
 F^{-1}(\ee^{L})=g(L)+\sum_{m\ge1}\frac{(-1)^m}{m!}
 \Bigl(\frac{\dd}{\dd L}\Bigr)^{m-1}\bigl(g'(L)h(L)^m\bigr),
 \label{ct:eq:transport}
\end{equation}
which is \cref{t1:eq:phaseinv} with $z=1$ and $R=h$: with $z=\phi_0(x)$ the
equation $F(x)=\ee^L$ reads $z+h(z)=L$, and Lagrange--B\"urmann applied to
$g(z)$ rather than to $z$ gives the display.  At each locally finite action
cutoff only finitely many sector products occur, and the coefficients may
contain exact functions of $g(L)$.  The formula shows why derivatives of
oscillatory coefficients must never be omitted.
\end{remark}

% ---- T2 lines 305--316 ----
For example, through second order in an auxiliary sector strength $\eta$, replacing $\varepsilon$ by $\eta\varepsilon$ gives
\begin{align}
 F^{-1}(\ee^L)
 ={}&g-\eta g'\varepsilon(g)\nonumber\\
 &+\frac{\eta^2}{2}\left\{g'\varepsilon(g)^2
       +\frac{\dd}{\dd L}\bigl(g'\varepsilon(g)^2\bigr)\right\}
       +\bigO(\eta^3).
 \label{t2:eq:transport-second}
\end{align}
This displays two sources of nonlinear inverse layers: the logarithm of the forward correction and the change of the evaluation point. Even a missing forward layer can be generated by these operations.

There is an important analytic qualification. An exact $g$ and an exact sector decomposition resolve an exponentially small effect. Replacing $g$ by a fixed finite inverse-power approximation generally does not: the replacement error is usually much larger than that effect. One may either keep $g$ exact, use a rigorously controlled numerical evaluation, or specify a summation and optimal-truncation procedure. We use the first option in the two-saddle and moment examples.
% ---- hand chunk 02-ch2-scaled.tex ----

\section{Scaled local reversion}
\label{ct:sec:scaled}

The phase-coordinate theorem inverts a correction of positive order.  When
the first inverse correction is unbounded --- of order $\sqrt X/\log X$ for
involutions --- but still small relative to the core $X$, the right
displacement variable is relative, not absolute.

% ---- T2 lines 245--272 ----
Suppose an inversion problem has been reduced to
\begin{equation}
 H(x)+R(x)=L,\qquad H(X)=L.
 \label{t2:eq:core-equation}
\end{equation}
The variable $X$ is a \emph{core inverse}. It may already contain a Lambert $W$ function and so absorb infinitely many logarithmic corrections.

\begin{proposition}[Scaled local reversion]
\label{t2:prop:scaled}
Put $x=X(1+z)$ and define the removable quotient
\[
 T_X(u)=\frac{H(X(1+u))-H(X)}{u},\qquad T_X(0)=XH'(X).
\]
Assume $T_X(0)$ is invertible in the coefficient field. In a formal grading in which $R(X(1+u))/T_X(u)$ has positive order, the solution with $z$ of positive order is
\begin{equation}
 z=\sum_{m\ge1}\frac{(-1)^m}{m}
 [u^{m-1}]\left(\frac{R(X(1+u))}{T_X(u)}\right)^m.
 \label{t2:eq:scaledLB}
\end{equation}
For an additive displacement $x=X+z$, the same formula holds with
$T_X(u)=(H(X+u)-H(X))/u$ and $R(X+u)$.
\end{proposition}
\begin{proof}
Equation~\eqref{t2:eq:core-equation} becomes
$z=-R(X(1+z))/T_X(z)$. Insert a bookkeeping parameter $\eta$ multiplying the right-hand side. Ordinary Lagrange--B\"urmann inversion gives the coefficient of $\eta^m$ as the $m$th summand in \eqref{t2:eq:scaledLB}. Positive formal order makes each coefficient in the intended grading a finite sum, so $\eta$ may be set to one formally. Analytically, convergence requires an additional local implicit-function argument; it is not a consequence of this formal calculation alone.
\end{proof}

The point of using a scaled displacement is that the leading inverse correction can be much larger than one. For involutions it is of order $\sqrt X/\log X$; relative to $X$, however, it is small. The wrong choice of an ``infinitesimal'' displacement can obscure an otherwise straightforward reversion.
% ---- hand chunk 02-ch2-linlog.tex ----

\section{Linear--logarithmic cores and all algebraic corrections}
\label{ct:sec:linlog}

The commonest dominant balance in this volume is
\begin{equation}
 \log F(x)\sim ax+\beta\log x+\log C+\sum_{j\ge1}q_jx^{-j},\qquad a>0,
 \label{ct:eq:linlog-model}
\end{equation}
with a differentiated finite-order expansion: balanced gamma quotients
(\cref{ct:ch:gamma}) have this form with the $q_j$ given by Bernoulli
polynomials, and so do the dominant endpoint blocks of the moment sequences
(\cref{ct:ch:moments}), with the $q_j$ obtained by the logarithmic partition
formula from an endpoint amplitude.  In the notation of the gamma-quotient
master expansion \cref{t2:eq:gamma-log} proved in \cref{ct:ch:gamma}, this is
the balanced case $\kappa=0$, $\lambda=a$; the theorem below is stated there
in that notation, but its proof uses only \cref{ct:eq:linlog-model} and a
positive limiting logarithmic slope.

% ---- T2 lines 396--451 ----
Assume $\kappa=0$ and $\lambda=a>0$. With $L=\log(y/C)$, define the large solution $X$ of
\begin{equation}
 aX+\beta\log X=L.
 \label{t2:eq:balanced-core}
\end{equation}
Its exact expression is
\begin{equation}
 X=\begin{cases}
 L/a,&\beta=0,\\[1mm]
 \displaystyle\frac{\beta}{a}
 W_{\star}\!\left(\frac a\beta\ee^{L/\beta}\right),&\beta\ne0,
 \end{cases}
 \label{t2:eq:balanced-W}
\end{equation}
where $W_\star=W_0$ for $\beta>0$ and $W_\star=W_{-1}$ for $\beta<0$. The second branch in the negative-$\beta$ case is essential: $W_0$ would give the small rather than the large solution. The branch notation agrees with \texttt{ProductLog[k,z]} \cite{wlproductlog}.

\begin{theorem}[Explicit centered inverse coefficients]
\label{t2:thm:balanced-inverse}
For the specified gamma-quotient interpolation,
\begin{equation}
 F^{-1}(y)\sim X+\sum_{n\ge1}d_nX^{-n}.
 \label{t2:eq:balanced-inverse}
\end{equation}
Let $Q(t)=\sum_{j\ge1}q_jt^j$ and
\[
 \Psi(t,u)=-\frac1a\left\{\beta\log(1+tu)
                     +Q\!\left(\frac{t}{1+tu}\right)\right\}.
\]
Then the nonrecursive finite formula
\begin{equation}
 d_n=\sum_{m=1}^n\frac1m[t^nu^{m-1}]\Psi(t,u)^m
 \label{t2:eq:balanced-d}
\end{equation}
computes every coefficient. In particular,
\begin{align}
 d_1&=-\frac{q_1}{a},\nonumber\\
 d_2&=\frac{\beta q_1}{a^2}-\frac{q_2}{a},\nonumber\\
 d_3&=-\frac{\beta^2q_1}{a^3}
       +\frac{\beta q_2-q_1^2}{a^2}-\frac{q_3}{a}.
 \label{t2:eq:balanced-first}
\end{align}
Truncation through $d_NX^{-N}$ has error $\bigO(X^{-N-1})$.
\end{theorem}
\begin{proof}
Write $x=X+U$, $t=X^{-1}$. Subtracting \eqref{t2:eq:balanced-core} from $\log F(x)=\log y$ gives formally $U=\Psi(t,U)$. Since $\Psi$ has positive $t$-order, Lagrange--B\"urmann gives \eqref{t2:eq:balanced-d}; terms $m>n$ cannot contribute. Direct coefficient comparison yields \eqref{t2:eq:balanced-first}.

For the analytic assertion, substitute a finite truncation into a sufficiently long finite expansion of $\log F$. Coefficient cancellation leaves a residual $\bigO(X^{-N-1})$. Its derivative is $a+\bigO(X^{-1})\ge a/2$ on a fixed neighborhood for large $X$. Lemma~\ref{t1:thm:certificate} supplies an inverse with the asserted error and proves its local uniqueness. The derivative asymptotics also prove eventual monotonicity globally.
\end{proof}

Expanding $W_{-1}$ itself converts this compact expression into the familiar nested logarithms. For instance, solving \eqref{t2:eq:balanced-core} by substitution gives
\[
 X=\frac1a\left\{L-\beta\log(L/a)
          +\frac{\beta^2\log(L/a)}{L}
          +\bigO\!\left(\frac{(\log L)^2}{L^2}\right)\right\}.
\]
Keeping $X$ exact avoids repeatedly expanding these logarithmic terms.
% ---- T1 lines 283--298 ----
The Lambert core can be eliminated if an expansion in ordinary logarithms is preferred. With $h=\log(L/a)$, the first terms are
\begin{equation}
 F^{-1}(y)=\frac1a\left[
 L-bh+\frac{b^2h-ad_1}{L}
 +\OO\!\left(\frac{(1+|h|)^2}{L^2}\right)\right].
 \label{t1:eq:flatten}
\end{equation}
The general coefficients follow by substituting this ansatz in the phase, or directly from \eqref{t2:eq:balanced-d} and the Lambert expansion. Retaining $X$ is usually simpler and numerically better conditioned.

\subsection{Action collisions and finite computation}
If $q_j(x)=x^{\beta_j}\ee^{-\lambda_jx}U_j(x)$ with finitely many $\lambda_j>0$, an inverse coefficient of multi-degree $\nu$ carries the action
\begin{equation}
 \Lambda_\nu=\sum_j\nu_j\lambda_j.
 \label{t1:eq:action}
\end{equation}
Only finitely many multi-indices lie below any fixed action cutoff, because $|\nu|\le T/\min\lambda_j$. When different multi-indices have the same action, their coefficients must be \emph{added}. Treating them as distinct physical exponentials overcounts the answer. Polynomial and trigonometric amplitude rings are stable under the derivatives in \eqref{t1:eq:multiinv}; no unspecified auxiliary recurrence is needed to generate a coefficient.
% ---- hand chunk 02-ch2-quadratic.tex ----

\section{Convergent inverse transseries for analytic exponential tails}
\label{ct:sec:quadratic}

The $q$-product families of \cref{ct:part:qproducts}, the connected-graph
counts, and the fixed-radical necklace sheets have a different dominant
balance: an exact quadratic, linear, or logarithmic phase $P(x)$ together
with a tail $B(\ee^{-\lambda x})$ that is \emph{analytic} in the small
variable.  For them the inverse transseries genuinely converges, and every
coefficient has a finite nonrecursive formula.  The section proves this
once, in the quadratic-or-linear case first and then for a general analytic
phase.

% ---- T3 lines 262--424 ----
\label{t3:sec:engine}

\subsection{Statement and analytic control}

Let
\begin{equation}\label{t3:eq:Pquadratic}
 P(x)=a x^2+b x+c,
\end{equation}
where either $a>0$, or $a=0$ and $b>0$. The function $B$ in \eqref{t3:eq:basic-shape} is fixed, analytic at zero, and real for sufficiently small real arguments. Let $L=\log Y$ and choose the large real solution of $P(s)=L$:
\begin{equation}\label{t3:eq:center}
 s=\begin{cases}
 \displaystyle\frac{-b+\sqrt{b^2+4a(L-c)}}{2a},&a>0,\\[0.8em]
 \displaystyle\frac{L-c}{b},&a=0.
 \end{cases}
 \qquad D=P'(s)=2as+b.
\end{equation}

\begin{theorem}[Convergent inverse transseries]\label{t3:thm:inverse}
For all sufficiently large real $Y$, the function
$F(x)=\exp(P(x)+B(\ee^{-\lambda x}))$ is increasing near its large real inverse, and
\begin{equation}\label{t3:eq:inverse-full}
 F^{-1}(Y)=s+\sum_{m\ge1}d_m(D)\,t_0^m,
 \qquad t_0=\ee^{-\lambda s}.
\end{equation}
The series converges. In fact, there exist positive constants $R$, $C$, and $s_0$, independent of $s\ge s_0$, such that
\begin{equation}\label{t3:eq:inverse-remainder}
 \left|F^{-1}(Y)-s-\sum_{m=1}^M d_m(D)t_0^m\right|
 \le \frac{C}{D}\frac{(t_0/R)^{M+1}}{1-t_0/R}
 \quad(0<t_0<R).
\end{equation}
Each $d_m$ is an explicitly computable finite polynomial in $D^{-1}$, with coefficients polynomial in $a,\lambda,b_1,\ldots,b_m$. The highest possible power of $D^{-1}$ is $2m-1$.
\end{theorem}

\begin{proof}
Write $x=s+\delta$. The inverse equation is exactly
\begin{equation}\label{t3:eq:implicit-delta}
 D\delta+a\delta^2+B(t\ee^{-\lambda\delta})=0,
\end{equation}
with $t=t_0$ at the end. Choose $D_0>0$ such that $D\ge D_0$ in the range considered. Choose $\rho>0$ so small that $a\rho\le D_0/4$; this restriction is unnecessary when $a=0$. Then choose $R>0$ so small that $B$ is analytic on the closed disk $|u|\le R\ee^{\lambda\rho}$ and
\[
 \sup_{|u|\le R\ee^{\lambda\rho}}|B(u)|<D_0\rho/4.
\]
On $|\delta|=\rho$ and $|t|\le R$, the sum of the last two terms of \eqref{t3:eq:implicit-delta} has modulus less than $D|\delta|$. Rouch\'e's theorem gives exactly one zero in the disk, counted with multiplicity. The zero is therefore simple and depends holomorphically on $t$. It is the branch satisfying $\delta(0)=0$.

Since $B(0)=0$, there is a constant $K$ such that
$|B(t\ee^{-\lambda\delta})|\le K|t|$ on these disks. Equation \eqref{t3:eq:implicit-delta} implies
\[
 (D-a\rho)|\delta|\le K|t|,
 \qquad |\delta|\le \frac{4K|t|}{3D}.
\]
Cauchy's coefficient estimates on a slightly smaller fixed $t$-disk yield
$|d_m|\le C D^{-1}R^{-m}$ after changing $C$ and $R$. Summing the resulting geometric tail proves \eqref{t3:eq:inverse-remainder}. For large $s$, $t_0<R$.

Finally,
\[
 (\log F)'(x)=2ax+b-\lambda\ee^{-\lambda x}B'(\ee^{-\lambda x})>0
\]
for large $x$, so the local branch is the large real inverse. The finite polynomial assertion follows from the explicit formula proved next.
\end{proof}

\subsection{A Lagrange--B\"urmann formula}

Define the local response to a prescribed logarithmic perturbation by
\begin{equation}\label{t3:eq:phi}
 \Phi_D(u)=\begin{cases}
 \displaystyle\frac{-D+\sqrt{D^2-4aB(u)}}{2a},&a>0,\\[0.8em]
 -B(u)/D,&a=0,
 \end{cases}
\end{equation}
where the square root equals $D$ at zero. Then $D\Phi_D+a\Phi_D^2=-B(u)$.

\begin{theorem}[Closed coefficient extraction]\label{t3:thm:coefficient}
For $m\ge1$,
\begin{equation}\label{t3:eq:inverse-extraction}
 \boxed{\displaystyle
 d_m(D)=-\frac1{m\lambda}\coefB{u}{m}
          \exp\bigl(-m\lambda\Phi_D(u)\bigr).}
\end{equation}
Equivalently, put
\begin{equation}\label{t3:eq:Cmk}
 C_{m,k}=\coefB{u}{m}B(u)^k
 =\sum_{\substack{j_1,\ldots,j_k\ge1\\j_1+\cdots+j_k=m}}
       b_{j_1}\cdots b_{j_k}.
\end{equation}
Then the completely finite, nonrecursive formula is
\begin{equation}\label{t3:eq:finite-dm}
 \boxed{\displaystyle
 d_m=-\sum_{k=1}^m\frac{C_{m,k}}{k}
 \sum_{j=0}^{k-1}\binom{k+j-1}{j}
 \frac{a^j(m\lambda)^{k-1-j}}
 {(k-1-j)!\,D^{k+j}}.}
\end{equation}
The $a=0$ case is obtained by retaining only $j=0$.
\end{theorem}

\begin{proof}
Set $u=t\ee^{-\lambda\delta}$. Equation \eqref{t3:eq:implicit-delta} becomes
$\delta=\Phi_D(u)$ and $u=t\exp(-\lambda\Phi_D(u))$.
The elementary Lagrange--B\"urmann identity gives
\[
 [t^m]\Phi_D(u(t))
 =\frac1m[u^{m-1}]\Phi_D'(u)\ee^{-m\lambda\Phi_D(u)}.
\]
Differentiating the exponential and using
$[u^{m-1}]H'(u)=m[u^m]H(u)$ proves \eqref{t3:eq:inverse-extraction}.

To obtain a finite formula, write $v=B(u)$ and
$w(v)=-\Phi_D(u)$. Thus $Dw-aw^2=v$, or
$w=v/(D-aw)$. A second application of Lagrange inversion gives, for $k\ge1$,
\[
 [v^k]\ee^{m\lambda w(v)}
 =\frac{m\lambda}{k}[w^{k-1}]
      \ee^{m\lambda w}(D-aw)^{-k}.
\]
Expand the exponential and the negative binomial power. Their coefficient of $w^{k-1}$ is
\[
 \sum_{j=0}^{k-1}\binom{k+j-1}{j}
 \frac{a^j(m\lambda)^{k-1-j}}{(k-1-j)!D^{k+j}}.
\]
Substitute this expression into \eqref{t3:eq:inverse-extraction}, and expand $B(u)^k$. Only $k\le m$ contributes. This proves \eqref{t3:eq:finite-dm} and also the bound $k+j\le2m-1$ on the reciprocal power of $D$.
\end{proof}

\subsection{The first three sectors and a computation recurrence}

The first coefficients are
\begin{align}
 d_1={}&-\frac{b_1}{D},\label{t3:eq:d1}\\
 d_2={}&-\frac{b_2}{D}-\frac{\lambda b_1^2}{D^2}
                   -\frac{a b_1^2}{D^3},\label{t3:eq:d2}\\
 d_3={}&-\frac{b_3}{D}-\frac{3\lambda b_1b_2}{D^2}
 -\frac{2a b_1b_2+\tfrac32\lambda^2b_1^3}{D^3}
 -\frac{3a\lambda b_1^3}{D^4}
 -\frac{2a^2b_1^3}{D^5}.\label{t3:eq:d3}
\end{align}
The $D^{-3}$ term in $d_2$ comes from curvature of the dominant phase, while its $D^{-2}$ term comes from shifting the exponential argument. Keeping only the response $-B(t_0)/D$ misses both.

For computation, let $\delta_{<m}(t)=\sum_{j=1}^{m-1}d_jt^j$. Coefficient comparison in \eqref{t3:eq:implicit-delta} gives the triangular recurrence
\begin{equation}\label{t3:eq:dm-recurrence}
 d_m=-\frac1D\left(
 a\sum_{i=1}^{m-1}d_i d_{m-i}
 +\sum_{k=1}^m b_k\coefB{t}{m-k}
     \ee^{-k\lambda\delta_{<m}(t)}\right).
\end{equation}
This is convenient for numerical work; \eqref{t3:eq:finite-dm} shows that no recurrence is required in the mathematical specification of a coefficient.

\begin{corollary}[Allowed weights]\label{t3:cor:support}
Let $S=\{m\ge1:b_m\ne0\}$. A coefficient $d_m$ can be nonzero only when $m$ is a finite sum of elements of $S$. In particular, if $B(t)=\widehat B(t^r)$, then $d_m=0$ unless $r\mid m$.
\end{corollary}
\begin{proof}
In \eqref{t3:eq:finite-dm}, every term contributing to $C_{m,k}$ has $m=j_1+\cdots+j_k$ with $j_i\in S$. This proves the assertion. Some allowed weights may still cancel; membership in the semigroup is not a guarantee of nonvanishing.
\end{proof}

\begin{corollary}[Direction of the first correction]\label{t3:cor:sign}
If $b_j=0$ for $j<r$ and $b_r\ne0$, then
\[
 F^{-1}(Y)-s=-\frac{b_r}{D}\ee^{-r\lambda s}
 +O\!\left(\frac{\ee^{-(r+1)\lambda s}}{D}\right).
\]
If all $b_m\ge0$, then all $d_m\le0$ for $a\ge0$.
\end{corollary}
\begin{proof}
The first assertion follows from the finite formula and the remainder bound. In the second assertion all factors in \eqref{t3:eq:finite-dm}, except the initial minus sign, are nonnegative.
\end{proof}
% ---- T3 lines 1316--1361 ----
\subsection{A general inverse formula beyond quadratic phases}

The coefficient extraction mechanism does not depend on $P$ being quadratic. Let $X(v)$ be the local inverse of $P$ at $v=P(s)$, and define
\begin{equation}\label{t3:eq:generalPhi}
 \Phi_s(u)=X(P(s)-B(u))-s.
\end{equation}
Exactly as in Section~\ref{t3:sec:engine}, putting
$u=t\ee^{-\kappa\delta}$ gives $\delta=\Phi_s(u)$. Hence
\begin{equation}\label{t3:eq:general-inverse-extraction}
 d_m(s)=-\frac1{m\kappa}[u^m]\exp\bigl(-m\kappa\Phi_s(u)\bigr).
\end{equation}
For \eqref{t3:eq:logphase}, this is explicit in Lambert $W$:
\begin{equation}\label{t3:eq:W-Phi}
 \Phi_s(u)=-\frac1h W_{-1}\!\left(-\frac hY E_R(u)\right)-s,
\end{equation}
where the branch is the local analytic continuation of its real value at $u=0$.

There is also a finite derivative formula that avoids expanding a special function. Let $b_m=[u^m]B(u)$ and let $C_{m,k}$ be as in \eqref{t3:eq:Cmk}. Define an operator on functions of $s$ by
\begin{equation}\label{t3:eq:shiftedoperator}
 \mathcal D_m=\frac1{P'(s)}\left(\frac{\ddn}{\ddn s}-m\kappa\right).
\end{equation}
Then
\begin{equation}\label{t3:eq:general-derivative}
 \boxed{\displaystyle
 d_m(s)=\sum_{k=1}^m\frac{(-1)^kC_{m,k}}{k!}
             \mathcal D_m^{\,k-1}\left(\frac1{P'(s)}\right).}
\end{equation}
Only finitely many differentiations are involved at a prescribed weight.

\begin{proof}[Proof of the derivative formula]
Temporarily insert a scalar parameter $\eta$ in
$P(x)+\eta B(z\ee^{-\kappa x})=P(s)$, and put $v=P(x)$. Ordinary Lagrange inversion for
$v+\eta\widetilde B(v)=P(s)$, followed by the analytic function $x=X(v)$, gives the coefficient of $\eta^k$ as
\[
 \frac{(-1)^k}{k!}
 \left(\frac{\ddn}{\ddn v}\right)^{k-1}
 \left(X'(v)\widetilde B(v)^k\right)_{v=P(s)}.
\]
Extract the coefficient of $z^m$. It is
\[
 \frac{(-1)^k C_{m,k}}{k!}
 \left(\frac1{P'(s)}\frac{\ddn}{\ddn s}\right)^{k-1}
 \left(\frac{\ee^{-m\kappa s}}{P'(s)}\right).
\]
Conjugating the differential operator by $\ee^{-m\kappa s}$ gives \eqref{t3:eq:shiftedoperator}. For fixed $m$, terms with $k>m$ cannot contribute. Set $\eta=1$ and remove the common factor $\ee^{-m\kappa s}$ to obtain \eqref{t3:eq:general-derivative}.
\end{proof}
% ---- hand chunk 02-ch2-residual.tex ----

\section{Residuals certify inverse errors}
\label{ct:sec:residual}

Every inverse expansion in this volume is asymptotic or convergent in a
stated sense; what turns a finite truncation into a trustworthy number is a
residual bound divided by a slope bound.

% ---- T1 lines 1256--1275 ----
\begin{theorem}[Local inverse certificate]
\label{t1:thm:certificate}
Let $F>0$ on $I=[x_*-h,x_*+h]$, let $\Phi=\log F$ be continuously differentiable there, and suppose $\Phi'(x)\ge m>0$ on $I$. If
\begin{equation}
 |\Phi(x_*)-\log y|<mh,
 \label{t1:eq:residualbracket}
\end{equation}
then $F(x)=y$ has a unique root $x\in I$, and
\begin{equation}
 |x-x_*|\le\frac{|\Phi(x_*)-\log y|}{m}.
 \label{t1:eq:inversebound}
\end{equation}
If only an approximation $\widetilde\Phi(x_*)$ with error at most $\epsilon$ is known, replace the numerator in both tests by
$|\widetilde\Phi(x_*)-\log y|+\epsilon$.
\end{theorem}
\begin{proof}
Integrating the derivative lower bound to either endpoint shows that the two endpoint residuals have opposite signs under \eqref{t1:eq:residualbracket}. Continuity gives a root, strict monotonicity gives uniqueness, and the mean value theorem gives \eqref{t1:eq:inversebound}. The version with an approximate phase follows from the triangle inequality.
\end{proof}

This elementary certificate is stronger than checking that the last displayed term is small. A small last term does not bound an uncomputed divergent tail, and a small forward relative error is not automatically a small inverse error unless the logarithmic slope is controlled.
% ---- T2 lines 334--344 ----
For rapidly growing positive functions it is usually best to apply this lemma to $f=\log F$. A relative forward error of size $r$ becomes an inverse error of order $r/(\log F)'$. The relevant slopes in this volume are approximately
\[
 \begin{array}{c|c}
 \text{forward growth}&\text{logarithmic slope}\\\hline
 C\ee^{ax}x^\beta&a\\
 (x/\ee)^{x/2}\ee^{\sqrt x}&\tfrac12\log x\\
 \Gamma(x+1)c^x&\log x+\log c\\
 2^{x(x-1)/2}&(x-\tfrac12)\log2
 \end{array}
\]
For harmonic functions, by contrast, the derivative of the function itself is about $1/x$, so forward errors are amplified by a factor of order $x$.
% ---- T1 lines 1277--1283 ----
For balanced gamma ratios one may use the Binet/Stirling bound \eqref{t1:eq:binetbound}. For the exact moment blocks, finite beta expansions with Taylor remainder bounds or rigorous integration give a phase enclosure. For zigzag Dirichlet tails, the simple bound
\begin{equation}
 \sum_{\substack{d>D\\d\ \mathrm{odd}}}d^{-s}
 \le D^{-s}+\frac{D^{1-s}}{s-1},\qquad s>1,
 \label{t1:eq:dirichlettail}
\end{equation}
comes from comparison with a decreasing integral; it is intentionally not sharp. Derivative tails can similarly be bounded by integrals of $(\log t)^j t^{-s}$. For the Gaussian-binomial logarithm, Cauchy bounds on a circle $|t|=r<\sqrt q$ or direct absolute sums in \eqref{t1:eq:qbinH} bound the remaining convergent tail. These analytic bounds, together with a derivative lower bound on an explicit interval, can be turned into rigorous inverse enclosures. The numerical tests below do not implement interval arithmetic and are not being presented as such certificates.
% ---- hand chunk 03-part2-head.tex ----

\part{Gamma quotients and the negative Lambert branch}
\label{ct:part:gamma}

\chapter{The gamma-quotient engine and the balanced families}
\label{ct:ch:gamma}

Many combinatorial sequences with purely exponential growth times an
algebraic factor are quotients of gamma functions at linear arguments:
Catalan and Fuss--Catalan numbers, central multinomial coefficients, and
the numbers of standard Young tableaux of a rectangle of fixed height.  A
single expansion, driven by the Bernoulli polynomials of the shifted
Stirling series, gives all their forward coefficients, and the balanced
core theorem of \cref{ct:sec:linlog} gives all their inverse coefficients.
The recurring feature is the \emph{negative} Lambert branch: an algebraic
prefactor $x^{\beta}$ with $\beta<0$ forces $W_{-1}$.

% ---- T2 lines 349--392 ----
\section{A master expansion}

Consider the positive function
\begin{equation}
 F(x)=K\ee^{\mu x}x^\nu
       \prod_{i=1}^r\Gamma(a_ix+b_i)^{\epsilon_i},
 \qquad K>0,\quad a_i>0,
 \label{t2:eq:gamma-general}
\end{equation}
where all parameters are fixed and real. Real powers are defined using the positive gamma values at large positive arguments.

\begin{theorem}[Gamma-quotient coefficient engine]
\label{t2:thm:gamma}
The logarithm of \eqref{t2:eq:gamma-general} has the all-orders expansion
\begin{equation}
 \log F(x)\sim\kappa x\log x+\lambda x+\beta\log x+\log C
                      +\sum_{j\ge1}q_jx^{-j},
 \label{t2:eq:gamma-log}
\end{equation}
where
\begin{align}
 \kappa&=\sum_i\epsilon_i a_i,
 &\lambda&=\mu+\sum_i\epsilon_i a_i(\log a_i-1),\nonumber\\
 \beta&=\nu+\sum_i\epsilon_i(b_i-\tfrac12),
 &\log C&=\log K+\frac{\log(2\pi)}2\sum_i\epsilon_i
           +\sum_i\epsilon_i(b_i-\tfrac12)\log a_i,
 \label{t2:eq:gamma-constants}\\
 q_j&=\frac{(-1)^{j+1}}{j(j+1)}
       \sum_i\frac{\epsilon_iB_{j+1}(b_i)}{a_i^j}.
 \label{t2:eq:gamma-q}
\end{align}
The normalized amplitude coefficients are $\cE_n(q)$.
\end{theorem}
\begin{proof}
Insert the shifted Stirling expansion
\[
 \log\Gamma(ax+b)
 \sim (ax+b-\tfrac12)\log(ax)-ax+\tfrac12\log(2\pi)
 +\sum_{j\ge1}\frac{(-1)^{j+1}B_{j+1}(b)}{j(j+1)(ax)^j}
\]
into the logarithm of \eqref{t2:eq:gamma-general}. Collect $x\log x$, $x$, $\log x$, constants, and inverse powers. There are finitely many gamma factors, so the fixed-order remainders add without changing their order. Formula~\eqref{t2:eq:E} then exponentiates the result. Differentiated versions follow from the corresponding differentiated gamma asymptotics \cite{dlmfgamma}.
\end{proof}

The case $\kappa=0$ is called \emph{balanced}. It includes many combinatorial sequences with purely exponential leading growth multiplied by an algebraic factor. The gamma origin also specifies a real interpolation; it does not specify a unique interpolation among all smooth functions.
% ---- hand chunk 03-balanced.tex ----

\section{All-orders inversion of the balanced case}
\label{ct:sec:balanced}

When $\kappa=0$ and $\lambda=a>0$, the expansion \cref{t2:eq:gamma-log} is
the linear--logarithmic model \cref{ct:eq:linlog-model}, and
\cref{t2:thm:balanced-inverse} gives the inverse to all orders: the core
\cref{t2:eq:balanced-W} is the large solution of $aX+\beta\log X=\log(y/C)$,
with $W_{-1}$ when $\beta<0$, and the coefficients $d_n$ of
\cref{t2:eq:balanced-d} are finite expressions in $a$, $\beta$, and the
$q_j$.  The first source stated the balanced case directly, as a theorem on
products $K\prod_i\Gamma(a_ix+b_i)^{\epsilon_i}$ with $\sum_i\epsilon_ia_i=0$;
its parameters $a=\sum_i\epsilon_ia_i\log a_i$,
$b=\sum_i\epsilon_i(b_i-\tfrac12)$,
$C=K(2\pi)^{\frac12\sum_i\epsilon_i}\prod_ia_i^{\epsilon_i(b_i-1/2)}$ are the
constants \cref{t2:eq:gamma-constants} with $\mu=\nu=0$ and $\kappa=0$, and
its logarithmic coefficients are \cref{t2:eq:gamma-q}.  The differentiated
Stirling expansion gives $(\log F)'=a+\beta/x+\bigO(x^{-2})$, so a balanced
gamma quotient with $a>0$ is positive and strictly increasing on a tail,
which is the hypothesis of the inverse theorem.

% ---- T1 lines 431--449 ----
\subsection{What is and is not fixed beyond all orders}
\label{t1:subsec:binet}
The gamma interpolation in \eqref{t2:eq:gamma-general} is an exact function, not an arbitrary function with a Stirling expansion. A convenient exact completion is Binet's formula \cite{dlmfbinet}:
\begin{equation}
 \log\Gamma(z)=(z-\tfrac12)\log z-z+\tfrac12\log(2\pi)
 +\int_0^\infty\ee^{-zt}
 \left(\frac12-\frac1t+\frac1{\ee^t-1}\right)\frac{\ddn t}{t},
 \quad z>0.
 \label{t1:eq:binet}
\end{equation}
It fixes the gamma block including its flat remainder. The Taylor expansion of the kernel at zero gives the Stirling coefficients. Its nonzero complex singularities occur at $2\pi\ii\mathbb Z\setminus\{0\}$, which explains factorial divergence, but does \emph{not} authorize assigning an arbitrary real-axis sector $c\ee^{-2\pi x}$.

For computational error control, keep the elementary logarithms in \eqref{t1:eq:binet} unexpanded and truncate the usual Stirling tail at $m=M-1$. With $z_i=a_ix+b_i>0$, the absolute error of the weighted logarithmic product is bounded by
\begin{equation}
 \mathcal B_M(x)=\sum_i|\epsilon_i|
 \frac{|B_{2M}|}{2M(2M-1)z_i^{2M-1}}.
 \label{t1:eq:binetbound}
\end{equation}
This uses the first-omitted-term bound on the positive real axis \cite{dlmfgamma}. It concerns this unexpanded-log evaluation; additional re-expansion errors must be included if one instead evaluates a polynomial in $1/x$.
% ---- T2 lines 453--495 ----
\section{Catalan numbers}
\label{t2:subsec:catalan}

The Catalan numbers count, among many equivalent objects, rooted full binary plane trees by internal vertices. Their standard formula and gamma continuation are \cite{oeiscatalan,wlcatalan}
\begin{equation}
 C_n=\frac1{n+1}\binom{2n}{n},\qquad
 C(x)=\frac{\Gamma(2x+1)}{\Gamma(x+1)\Gamma(x+2)}.
 \label{t2:eq:catalan}
\end{equation}
They illustrate the distinction between Catalan numbers as coefficients of a quadratic inverse and the inverse of the \emph{counting function} $C(x)$; it is the latter that is studied here.

Theorem~\ref{t2:thm:gamma} gives $a=\log4$, $\beta=-3/2$, $C=\pi^{-1/2}$, and
\begin{align}
 \log C(x)\sim{}&x\log4-\frac32\log x-\frac12\log\pi
 -\frac9{8x}+\frac1{2x^2}-\frac{21}{64x^3}
 +\frac1{4x^4}-\frac{129}{640x^5}+\cdots.
 \label{t2:eq:catalan-log}
\end{align}
Exponentiating,
\begin{equation}
 C(x)\sim\frac{4^x}{\sqrt\pi x^{3/2}}
 \left(1-\frac9{8x}+\frac{145}{128x^2}
        -\frac{1155}{1024x^3}+\frac{36939}{32768x^4}+\cdots\right).
 \label{t2:eq:catalan-amplitude}
\end{equation}
All coefficients are supplied by \eqref{t2:eq:gamma-q} and \eqref{t2:eq:E}, not just the displayed ones.

Define
\begin{equation}
 X=-\frac3{2a}W_{-1}\!\left[
       -\frac{2a}{3}\left(\frac1{\sqrt\pi y}\right)^{2/3}\right],
 \qquad a=\log4.
 \label{t2:eq:catalan-core}
\end{equation}
Then
\begin{align}
 C^{-1}(y)={}&X+\frac9{8aX}
 +\frac1{X^2}\left(\frac{27}{16a^2}-\frac1{2a}\right)\nonumber\\
 &+\frac1{X^3}\left(\frac{81}{32a^3}-\frac{129}{64a^2}
                    +\frac{21}{64a}\right)+\bigO(X^{-4}).
 \label{t2:eq:catalan-inverse}
\end{align}
There is no unknown constant or auxiliary recurrent sequence in this description. Gamma asymptotics and the finite coefficient sum determine every further term.
% ---- T2 lines 497--526 ----
\section{Fuss--Catalan numbers}

For a fixed integer $p\ge2$, full $p$-ary plane trees give the sequence
\begin{equation}
 C_{p,n}=\frac1{(p-1)n+1}\binom{pn}{n},\qquad
 C_p(x)=\frac{\Gamma(px+1)}{\Gamma(x+1)\Gamma((p-1)x+2)}.
 \label{t2:eq:fuss}
\end{equation}
The ternary case is OEIS A001764 \cite{oeisfuss}. The elementary generating-function equation $T=1+zT^p$ and Lagrange inversion also give \eqref{t2:eq:fuss}: for $U=T-1$, $U=z(1+U)^p$, so
$[z^n]U=\binom{pn}{n-1}/n$.

Here
\begin{equation}
 \rho_p=\frac{p^p}{(p-1)^{p-1}},\quad
 a_p=\log\rho_p,\quad \beta=-\frac32,\quad
 C_p^{(0)}=\sqrt{\frac{p}{2\pi(p-1)^3}},
 \label{t2:eq:fuss-constants}
\end{equation}
and
\begin{equation}
 q_j^{(p)}=\frac{(-1)^{j+1}}{j(j+1)}
 \left\{B_{j+1}(1)(p^{-j}-1)
             -\frac{B_{j+1}(2)}{(p-1)^j}\right\}.
 \label{t2:eq:fuss-q}
\end{equation}
For example,
\[
 q_1^{(p)}=\frac1{12}\left(\frac1p-1-\frac{13}{p-1}\right).
\]
Thus $C_p(x)\sim C_p^{(0)}\rho_p^x x^{-3/2}\sum_m\cE_m(q^{(p)})x^{-m}$, and its inverse is obtained from \eqref{t2:eq:balanced-W} and \eqref{t2:eq:balanced-d}. For $p=3$, $\rho_3=27/4$, $C_3^{(0)}=\sqrt3/(4\sqrt\pi)$, and $q_1=-43/72$. This entire fixed-$p$ family shares the same logarithmic exponent, while its growth rate and inverse correction polynomials vary with $p$.
% ---- hand chunk 03-fuss-extra.tex ----

\begin{remark}[The Fuss--Catalan inverse in closed form]
\label{ct:rmk:fuss-inverse}
For every $p\ge2$ the second logarithmic coefficient is
$q_2^{(p)}=1/(2(p-1)^2)$, since $B_3(1)=0$ and $B_3(2)=3$ in
\cref{t2:eq:fuss-q}.  With $L=\log(y/C_p^{(0)})$ and $a_p=\log\rho_p$, the
core and the first correction are
\begin{equation}
 X=-\frac{3}{2a_p}\,W_{-1}\!\Bigl(-\frac{2a_p}{3}\,\ee^{-2L/3}\Bigr),\qquad
 C_p^{-1}(y)=X-\frac{q_1^{(p)}}{a_pX}+\bigO(X^{-2}),
 \label{ct:eq:fuss-inverse}
\end{equation}
and \cref{t2:eq:fuss-q} with \cref{t2:eq:balanced-d} provides every further
term as a finite sum.  At $p=2$ this is \cref{t2:eq:catalan-inverse}.
\end{remark}

% ---- T2 lines 528--567 ----
\section{Central multinomial coefficients and rectangular tableaux}

For a fixed integer $d\ge2$, consider
\begin{equation}
 M_d(x)=\frac{\Gamma(dx+1)}{\Gamma(x+1)^d}.
 \label{t2:eq:multinomial}
\end{equation}
At $x=n$, this counts assignments of $dn$ labeled objects to $d$ labeled boxes containing $n$ objects each, by the elementary multinomial formula. Its constants are
\begin{equation}
 a=d\log d,\quad \beta=\frac{1-d}{2},\quad
 C=\sqrt d\,(2\pi)^{(1-d)/2}.
 \label{t2:eq:multi-constants}
\end{equation}
Only odd inverse powers occur in its logarithmic correction:
\begin{equation}
 q_{2m-1}=\frac{B_{2m}}{2m(2m-1)}
                 \bigl(d^{1-2m}-d\bigr),\qquad q_{2m}=0.
 \label{t2:eq:multi-q}
\end{equation}
Even inverse powers reappear after exponentiation and inversion. The absence of an order in the logarithm is therefore not the absence of that order in the final inverse.

The number of standard Young tableaux of a rectangle with $d$ rows and $n$ columns is
\begin{equation}
 T_d(n)=(dn)!\frac{\prod_{j=0}^{d-1}j!}{\prod_{j=0}^{d-1}(n+j)!}.
 \label{t2:eq:tableaux}
\end{equation}
One route to this formula is the hook-length formula: the hook lengths in row $i$ are $(n+d-i),(n+d-i-1),\ldots,(d-i+1)$, whose product is $(n+d-i)!/(d-i)!$. Multiplying over rows gives \eqref{t2:eq:tableaux}. The three-row case is recorded as A005789 \cite{oeistableaux}. We take the gamma version of \eqref{t2:eq:tableaux} as the real extension.

Now
\begin{align}
 a&=d\log d,&
 \beta&=-\frac{d^2-1}{2},&
 C&=\sqrt d\,(2\pi)^{(1-d)/2}\prod_{j=0}^{d-1}j!,
 \label{t2:eq:tableaux-constants}\\
 q_r&=\frac{(-1)^{r+1}}{r(r+1)}
 \left\{\frac{B_{r+1}(1)}{d^r}
             -\sum_{j=0}^{d-1}B_{r+1}(j+1)\right\}.
 \label{t2:eq:tableaux-q}
\end{align}
The inverse again uses $W_{-1}$. The case $d=2$ reduces exactly to Catalan numbers, providing a consistency check between two combinatorial descriptions. No assertion here is uniform as $d\to\infty$: varying the rectangle height changes the asymptotic problem.
% ---- hand chunk 03-tableaux-extra.tex ----

The rectangular family separates the exponential rate from the power-law
prefactor: it has the same rate $d\log d$ as the central multinomial family,
but its exponent $\beta$ is quadratic rather than linear in $d$.
Consequently the coefficient of $\log\log y$ in the flattened inverse
(the nested-logarithm expansion following \cref{t2:eq:balanced-first})
changes from $(d-1)/(2d\log d)$ to $(d^2-1)/(2d\log d)$.

% ---- T2 lines 569--579 ----
\section{What the gamma engine does not automatically supply}

When $\kappa>0$, the convenient unbalanced core is
\begin{equation}
 X=\frac{L/\kappa}{W_0\!\left((L/\kappa)\ee^{\lambda/\kappa}\right)},
 \qquad L=\log(y/C),
 \label{t2:eq:unbalanced-core}
\end{equation}
which solves $\kappa X\log X+\lambda X=L$. The remaining terms are handled by Proposition~\ref{t2:prop:scaled}. This observation connects the balanced families to the factorial-growth examples below.

However, the Bernoulli expansions of gamma functions are generally divergent. An all-orders expansion is not automatically a complete resurgent transseries with known Stokes multipliers. In this volume the gamma-quotient results are rigorous all-orders algebraic sectors and their inverse sectors. We do not attach arbitrary exponentially small terms to them or infer a summation convention from the integer samples.
% ---- hand chunk 04-part3-head.tex ----

\part{Finite exponential sums and fixed cycle counts}
\label{ct:part:finite}

\chapter{Stirling and Eulerian columns}
\label{ct:ch:columns}

The two kinds of Stirling numbers exhibit remarkably different inverse
scales when the second index is fixed.  Set partitions into $k$ blocks and
permutations with $m$ descents produce finite sums of ordinary exponentials
whose inverses are \emph{convergent} expansions in generalized powers of
$1/y$; permutations with $k$ cycles produce factorial growth multiplied by a
polynomial in $\log n$, whose inverse has a Lambert core, an order-one
shift, and a hierarchy of $\log\log X/\log X$ layers before any $X^{-1}$
correction.  The standard combinatorial meanings and sign conventions are
those of \cite{dlmfstirling,wlstirling2,wlstirling1}.

% ---- T2 lines 586--631 ----
\section{An exactly solvable multiexponential class}

\begin{theorem}[Convergent multivariate inverse]
\label{t2:thm:finite-exp}
Let
\begin{equation}
 F(x)=A\ee^{ax}\left(1+\sum_{j=1}^r c_j\ee^{-\omega_jx}\right),
 \quad A>0,\quad a>0,\quad \omega_j>0,
 \label{t2:eq:finite-exp}
\end{equation}
with fixed real coefficients. Put
\[
 X=\frac{\log(y/A)}a,\qquad
 \beta_j=\frac{\omega_j}{a},\qquad
 t_j=(y/A)^{-\beta_j}.
\]
For a nonzero multi-index $\bm m$, let $M=\sum_jm_j$ and
$B=\sum_j\beta_jm_j$. Then, for sufficiently large $y$,
\begin{equation}
 F^{-1}(y)=X-\frac1a
 \sum_{\bm m\ne0}
 \frac{(M-1)!}{\prod_jm_j!}
 \binom{B-1}{M-1}\prod_j(c_jt_j)^{m_j}.
 \label{t2:eq:finite-exp-inverse}
\end{equation}
This is an actually convergent multivariate Taylor series near $\bm t=0$, not merely a Poincar\'e expansion. The generalized binomial coefficient is the polynomial
$\binom{z}{m}=z(z-1)\cdots(z-m+1)/m!$.
\end{theorem}
\begin{proof}
Let $\delta=a(x-X)$ and $w=\ee^{-\delta}$. The equation $F(x)=y$ is equivalent to
\begin{equation}
 w=1+\sum_{j=1}^r c_jt_jw^{\beta_j}.
 \label{t2:eq:w-finite}
\end{equation}
The right-hand powers are analytic in $w$ near $1$. At $(w,\bm t)=(1,0)$ the derivative of the left side minus the right side with respect to $w$ is $1$. The analytic implicit-function theorem therefore gives a unique analytic $w(\bm t)$ near that point.

Write $u=w-1$, insert an auxiliary parameter $\eta$ multiplying the sum, and apply Lagrange inversion to $\log(1+u)$:
\[
 \log(1+u)=\sum_{M\ge1}\frac{\eta^M}{M}
 [z^{M-1}]\frac1{1+z}
       \left(\sum_jc_jt_j(1+z)^{\beta_j}\right)^M.
\]
The multinomial coefficient for a given $\bm m$ is $M!/\prod m_j!$, and the remaining extraction is $\binom{B-1}{M-1}$. Since $x=X-a^{-1}\log w$, this proves \eqref{t2:eq:finite-exp-inverse}. The analytic implicit-function construction proves convergence. Finally $F'(x)/(A\ee^{ax})\to a>0$, so the branch is the eventual real inverse.
\end{proof}

When two multi-indices have the same total action $B$, their contributions must be added. A total-degree truncation is useful computationally, but ordering the actual asymptotic terms requires ordering the weights $B$, not just the integers $M$. The action semigroup is locally finite because all $\beta_j$ are positive and there are finitely many generators.
% ---- hand chunk 04-stirling-falling.tex ----

\begin{remark}[The falling-factorial form of the coefficients]
\label{ct:rmk:falling}
Since $(M-1)!\binom{B-1}{M-1}=\prod_{h=1}^{M-1}(B-h)=\fall{B}{M}/B$, the
coefficient of $\prod_j(c_jt_j)^{m_j}$ in \cref{t2:eq:finite-exp-inverse}
is $-\fall{B}{M}/(aB\,\bm m!)$, with $\bm m!=\prod_jm_j!$.  This is also what
\cref{t1:thm:multi} produces directly: with $A(x)=A\ee^{ax}$,
$q_j(x)=c_j\ee^{-\omega_jx}$, $g=X$, $g'=1/a$ and $D=\dd/\dd L$, one has
$D^{m-1}(g'Q_{\bm m})=a^{-1}(-B)^{m-1}Q_{\bm m}$, and the signed-Stirling
identity \cref{t1:eq:stirlingsign} gives
$\sum_{m=1}^{M}(-1)^m\sone{M}{m}(-B)^{m-1}=-\fall{B}{M}/B=-\prod_{h=1}^{M-1}(B-h)$.
Whenever $B\in\{1,\ldots,M-1\}$ the coefficient vanishes; such cancellations
are structural, not numerical coincidences.  Marker degree $M$ and
numerical size are different orderings: a degree-one term with a large
action can be smaller than many higher-degree terms with smaller actions,
so the formula supports either a total-degree truncation or an
action-sorted truncation, provided the convention is stated.
\end{remark}

% ---- T2 lines 633--677 ----
\section{Stirling numbers of the second kind and surjections}

For a fixed integer $k\ge2$, inclusion--exclusion gives
\begin{equation}
 S_k(x)=\frac1{k!}\sum_{j=1}^k(-1)^{k-j}\binom{k}{j}j^x,
 \qquad x>0.
 \label{t2:eq:stirling2-cont}
\end{equation}
At positive integers $n$, this equals $\StirII nk$. Indeed, among the $k^n$ maps from an $n$-element labeled set to $k$ labeled boxes, inclusion--exclusion removes maps omitting one or more boxes. Dividing the number of surjections by $k!$ forgets the labels of the nonempty boxes. Equation~\eqref{t2:eq:stirling2-cont} is also our exact continuous interpolation.

Theorem~\ref{t2:thm:finite-exp} applies with
\begin{equation}
 A=\frac1{k!},\qquad a=\log k,\qquad
 c_j=(-1)^{k-j}\binom{k}{j},\qquad
 \omega_j=\log\frac{k}{j}\quad(1\le j<k).
 \label{t2:eq:stirling2-params}
\end{equation}
Thus every inverse coefficient is given by the finite generalized-binomial expression in \eqref{t2:eq:finite-exp-inverse}. Multiplying the forward function by $k!$ gives the corresponding inverse for the surjection count, simply by replacing the target $y$ by $y/k!$.

For $k=2$ the formula is elementary:
\[
 S_2(x)=2^{x-1}-1,\qquad
 S_2^{-1}(y)=1+\frac{\log(y+1)}{\log2}.
\]
Its convergent expansion in $1/y$ provides a check on the general theorem.

For $k=3$, define
\[
 \beta=\frac{\log(3/2)}{\log3},\qquad
 X=\frac{\log(6y)}{\log3},\qquad
 t=(6y)^{-\beta},\quad s=(6y)^{-1}.
\]
Since $S_3(x)=(3^x-3\cdot2^x+3)/6$, the inverse starts as
\begin{align}
 (\log3)(S_3^{-1}(y)-X)
 ={}&3t-3s+\frac92(1-2\beta)t^2+9\beta ts-\frac92s^2
 \nonumber\\
 &+\bigO((|t|+|s|)^3).
 \label{t2:eq:stirling3-inverse}
\end{align}
This display is arranged by multivariate degree. In dominance order, $t^2$ precedes $s$, since $2\beta<1$. The pure cubic term, useful for a dominance-ordered expansion, is
\[
 \frac92(3\beta-1)(3\beta-2)t^3.
\]
The complete formula automatically places mixed terms and accounts for any collisions. Unlike a factorial asymptotic series, this inverse expansion genuinely converges for sufficiently large target values.
% ---- T1 lines 530--575 ----
\section{Fixed numbers of descents: Eulerian columns}
Let $A(n,m)$ count permutations of $n$ elements with exactly $m$ descents, with zero-based $m$. The Eulerian triangle is A008292 \cite{oeiseulerian}. For fixed $m\ge1$ define
\begin{equation}
 E_m(x)=\sum_{j=0}^m(-1)^j\binom{x+1}{j}(m+1-j)^x.
 \label{t1:eq:eulerianinterp}
\end{equation}
For integer $n\ge1$ this is the standard finite inclusion--exclusion formula for $A(n,m)$. The binomial coefficient in \eqref{t1:eq:eulerianinterp} means the degree-$j$ polynomial $(x+1)x\cdots(x+2-j)/j!$, so there is no ambiguity at real $x$.

Put $a=\log(m+1)$, $X=\log y/a$, and, for $1\le j\le m$, put
\[
 P_j(x)=(-1)^j\binom{x+1}{j},\qquad
 \lambda_j=\log\frac{m+1}{m+1-j}.
\]
Then
\begin{equation}
 E_m(x)=\ee^{ax}\left(1+\sum_{j=1}^mP_j(x)\ee^{-\lambda_jx}\right).
 \label{t1:eq:eulerianforward}
\end{equation}
This is again exact, and it proves eventual positivity and monotonicity. It also displays a new feature: the subdominant exponential amplitudes are nonconstant polynomials.

\begin{proposition}[All Eulerian inverse sectors]\label{t1:prop:eulerian}
For $\nu\in\NN^m\setminus\{0\}$, let $r=|\nu|$, $\Lambda_\nu=\sum_j\nu_j\lambda_j$, and $P_\nu(X)=\prod_jP_j(X)^{\nu_j}$. The contribution of multi-degree $\nu$ to $E_m^{-1}(y)-X$ is
\begin{equation}
 \delta_\nu=
 \frac{\ee^{-\Lambda_\nu X}}{a\nu!}
 \sum_{\ell=1}^r(-1)^\ell s(r,\ell)a^{1-\ell}
 (\partial_X-\Lambda_\nu)^{\ell-1}P_\nu(X).
 \label{t1:eq:eulerianinverse}
\end{equation}
The sum of all these terms converges for sufficiently large $y$. Each sector amplitude is a polynomial in $X$, of degree at most $\sum_jj\nu_j$.
\end{proposition}
\begin{proof}
Use \cref{t1:thm:multi} with $g=X$ and $D=a^{-1}\partial_X$. Conjugating differentiation by $\ee^{-\Lambda_\nu X}$ gives
\[
 D^{\ell-1}(g'Q_\nu)=a^{-\ell}\ee^{-\Lambda_\nu X}
 (\partial_X-\Lambda_\nu)^{\ell-1}P_\nu(X).
\]
This proves the coefficient formula and the degree bound. On a fixed disk about large $X$, the functions $P_j(X+\delta)\ee^{-\lambda_j(X+\delta)}$ and their derivatives are uniformly exponentially small. The contraction argument after \cref{t1:thm:multi} therefore proves convergence at the physical marker values.
\end{proof}

For example, $E_1(x)=2^x-x-1$, hence
\[
 E_1^{-1}(y)=X+\frac{X+1}{\log2}\,2^{-X}
 +\OO(X^2 2^{-2X}),\qquad X=\log_2 y.
\]
Formula \eqref{t1:eq:eulerianinverse} supplies the complete amplitude at every order, not merely this first correction. After $2^{-X}=y^{-1}$ is substituted, the inverse is a convergent series of polynomials in $\log y$ multiplied by powers of $y^{-1}$. For general $m$, the exponents need not be integers.
% ---- T2 lines 679--802 ----
\section{Unsigned Stirling numbers of the first kind}

Write $c(n,k)=\StirI nk$ for the number of permutations of $n$ labels with $k$ cycles. The polynomial identity
\[
 z(z+1)\cdots(z+n-1)=\sum_{k=0}^n c(n,k)z^k
\]
follows by adding the new label either as a singleton cycle or in one of the existing cycle positions. It gives a particularly useful interpolation for each fixed $k\ge1$:
\begin{equation}
 C_k(x)=[z^k]\frac{\Gamma(x+z)}{\Gamma(z)},\qquad x>0.
 \label{t2:eq:stirling1-cont}
\end{equation}
The coefficient is taken in the analytic germ at $z=0$; $1/\Gamma(z)$ has a zero there. This is an explicit function, not a recurrence.

The Wolfram Language function \texttt{StirlingS1[n,k]} uses the \emph{signed} convention, so at integers
\[
 c(n,k)=(-1)^{n-k}\,\texttt{StirlingS1[n,k]}.
\]
The gamma expression \eqref{t2:eq:stirling1-cont} fixes the positive unsigned continuation; analytically continuing $(-1)^{x-k}$ would produce a different object.

Using the Taylor expansions of $\log\Gamma$ at $x$ and at $1$ gives the exact finite-coefficient representation
\begin{align}
 C_k(x)=\Gamma(x)[z^{k-1}]\exp\Bigg(& (\psi(x)+\gamma)z\nonumber\\
 &+\sum_{j\ge2}\left\{\frac{\psi^{(j-1)}(x)}{j!}
                  -\frac{(-1)^j\zeta(j)}j\right\}z^j\Bigg).
 \label{t2:eq:stirling1-polygamma}
\end{align}
Only $j\le k-1$ can affect this coefficient. Here $\psi$ is the logarithmic derivative of gamma and $\gamma$ is Euler's constant.

\begin{theorem}[All-orders fixed-column expansion]
\label{t2:thm:stirling1}
Put $\ell=\log x$ and
\begin{align}
 r_j(z)&=\frac{(-1)^{j+1}}{j(j+1)}
              \bigl(B_{j+1}(z)-B_{j+1}(0)\bigr),\nonumber\\
 P_{k,m}(\ell)&=[z^{k-1}]\frac{\ee^{\ell z}}{\Gamma(1+z)}
                      \cE_m(r(z)).
 \label{t2:eq:stirling1-P}
\end{align}
For fixed $k$,
\begin{equation}
 C_k(x)\sim\Gamma(x)\sum_{m\ge0}P_{k,m}(\log x)x^{-m}.
 \label{t2:eq:stirling1-forward}
\end{equation}
The leading polynomial $P_{k,0}$ has degree $k-1$, and every $P_{k,m}$ with $m\ge1$ has degree at most $k-2$. In particular,
\begin{align}
 P_{1,0}&=1,&P_{2,0}&=\ell+\gamma,\nonumber\\
 P_{3,0}&=\frac{(\ell+\gamma)^2-\zeta(2)}2,&
 P_{4,0}&=\frac{(\ell+\gamma)^3-3\zeta(2)(\ell+\gamma)+2\zeta(3)}6.
 \label{t2:eq:stirling1-leading}
\end{align}
\end{theorem}
\begin{proof}
Rewrite \eqref{t2:eq:stirling1-cont} as
\[
 C_k(x)=\Gamma(x)[z^{k-1}]
        \frac1{\Gamma(1+z)}\frac{\Gamma(x+z)}{\Gamma(x)}.
\]
Lemma~\ref{t2:lem:ratio} expands the gamma ratio uniformly on a fixed small circle about $z=0$. Cauchy's coefficient formula therefore permits extraction term by term, with a fixed-order remainder bounded by an inverse power times a fixed power of $\log x$. The expression \eqref{t2:eq:stirling1-P} follows.

For $m\ge1$, the polynomial $\cE_m(r(z))$ vanishes at $z=0$, because each $r_j(0)=0$. Its product with $\ee^{\ell z}$ can consequently contribute at most $\ell^{k-2}$ to $[z^{k-1}]$. For $m=0$ the leading term is $\ell^{k-1}/(k-1)!$. Finally,
\[
 \frac1{\Gamma(1+z)}
 =\exp\left(\gamma z+\sum_{j\ge2}\frac{(-1)^{j+1}\zeta(j)}jz^j\right)
\]
gives the displayed leading polynomials by finite extraction.
\end{proof}

\section{The nested-logarithmic inverse}

Define
\begin{equation}
 L=\log\left(\frac{y(k-1)!}{\sqrt{2\pi}}\right),\qquad
 X=\frac{L}{W_0(L/\ee)},\qquad \ell=\log X,
 \label{t2:eq:stirling1-core}
\end{equation}
and the polynomial in $1/\ell$
\begin{equation}
 Q_k(1/\ell)=\frac{(k-1)!P_{k,0}(\ell)}{\ell^{k-1}}.
 \label{t2:eq:stirling1-Q}
\end{equation}
It tends to one. Taking logarithms in Theorem~\ref{t2:thm:stirling1}, and also using Stirling's formula for $\Gamma(x)$, produces
\begin{align}
 \log\left(\frac{C_k(x)(k-1)!}{\sqrt{2\pi}}\right)
 \sim{}&x\log x-x-\tfrac12\log x
       +\log\bigl((k-1)!P_{k,0}(\log x)\bigr)
 \nonumber\\
 &+\sum_{m\ge1}U_{k,m}(\log x)x^{-m}.
 \label{t2:eq:stirling1-log}
\end{align}
Each $U_{k,m}$ is explicit: add the gamma logarithmic coefficient at order $m$ to
\[
 [t^m]\log\left(1+\sum_{j\ge1}
       \frac{P_{k,j}(\ell)}{P_{k,0}(\ell)}t^j\right).
\]
Thus the coefficients are rational functions of $\ell$, which themselves possess expansions in $1/\ell$.

\begin{corollary}[First inverse layers]
\label{t2:cor:stirling1-inverse}
For the interpolation \eqref{t2:eq:stirling1-cont},
\begin{equation}
 C_k^{-1}(y)=X+\frac12
       -\frac{(k-1)\log\ell}{\ell}
       -\frac{\log Q_k(1/\ell)}{\ell}
       +\bigO(X^{-1}).
 \label{t2:eq:stirling1-inverse}
\end{equation}
In particular,
\begin{equation}
 C_k^{-1}(y)=X+\frac12-
       \frac{(k-1)\log\log X}{\log X}
       -\frac{(k-1)\gamma}{(\log X)^2}
       +\bigO((\log X)^{-3}).
 \label{t2:eq:stirling1-inverse-expanded}
\end{equation}
All further $X^{-1}$ and logarithmic layers are obtained by applying Proposition~\ref{t2:prop:scaled} to \eqref{t2:eq:stirling1-log}.
\end{corollary}
\begin{proof}
The core $H(x)=x\log x-x$ has $H'(X)=\ell$. Its residual at $X$ is
$-\ell/2+(k-1)\log\ell+\log Q_k(1/\ell)+\bigO(X^{-1})$.
Subtracting this residual divided by $\ell$ gives \eqref{t2:eq:stirling1-inverse}. The proposed displacement is bounded. Taylor's theorem then leaves a residual $\bigO(X^{-1})$, whose inverse error is no larger than $\bigO(X^{-1})$ because the slope grows like $\ell$. Also
$Q_k(1/\ell)=1+(k-1)\gamma/\ell+\bigO(\ell^{-2})$, proving the second display. Differentiated asymptotics show positivity and eventual monotonicity, so the continuous inverse exists.
\end{proof}

A useful practical distinction emerges. For fixed $k$, second-kind inverses are logarithmic cores plus convergent generalized powers of $1/y$. First-kind inverses have a Lambert core, an order-one shift, and a hierarchy involving $\log\log X/\log X$ before any $X^{-1}$ correction is reached. The two ``Stirling'' inverse problems belong to different scales.
% ---- hand chunk 05-part4-head.tex ----

\part{Moment sequences and endpoint sectors}
\label{ct:part:moments}

\chapter{Endpoint extraction: Motzkin, trinomial, Delannoy, and Schr\"oder numbers}
\label{ct:ch:moments}

A sequence that is the moment sequence of a positive measure has a canonical
positive interpolation, the moment integral at a real exponent, and its
asymptotics are governed by the largest support point.  When the support
also has a negative part, splitting the integral at zero produces a second
\emph{exact} block, and the real interpolation carries a cosine factor whose
choice is deliberate.  Endpoint extraction turns both blocks into all-order
algebraic expansions with finite coefficients; the balanced core theorem
inverts the dominant block; and the multiplicative-sector theorem transports
the exact small block into an oscillatory inverse lattice.

% ---- T2 lines 807--838 ----
\section{An endpoint coefficient lemma}

A recurring integral has, near its largest positive support point $b$, the form
\[
 K\int t^{x+\alpha-1}(b-t)^{\nu-1}A(b-t)\dd t,
 \qquad \nu>0,
\]
where $A$ is analytic at zero. The change $t=b(1-v)$ reduces the endpoint to a beta-type integral. Taylor expansion of $A(bv)$ and Lemma~\ref{t2:lem:ratio} then produce every inverse-power coefficient.

\begin{lemma}[Endpoint extraction]
\label{t2:lem:endpoint}
Suppose the part of an integral away from $v=0$ is exponentially smaller than its $v=0$ endpoint contribution, and its local integrand is
\[
 Kb^x(1-v)^{x+\alpha-1}v^{\nu-1}(1-\kappa v)^\theta.
\]
After normalizing by $Kb^x\Gamma(\nu)x^{-\nu}$, its inverse-power coefficient at order $m$ is
\begin{equation}
 a_m=\sum_{j=0}^m\binom\theta j(-\kappa)^j
          \rising\nu j\,R_{m-j}(\alpha,\alpha+\nu+j).
 \label{t2:eq:endpoint-coeff}
\end{equation}
\end{lemma}
\begin{proof}
On a small fixed interval about zero, expand $(1-\kappa v)^\theta$ through order $N$. The remainder is bounded by a constant times $v^{N+1}$. Its integral is $\bigO(x^{-\nu-N-1})$, as follows either by comparison with $\ee^{-xv}$ or by a beta integral. Extending each polynomial term to $v=1$ changes it only exponentially, since the extension begins a fixed distance from zero. The $j$th beta integral is
\[
 B(x+\alpha,\nu+j)=\Gamma(\nu+j)
 \frac{\Gamma(x+\alpha)}{\Gamma(x+\alpha+\nu+j)}.
\]
Apply Lemma~\ref{t2:lem:ratio}, divide by $\Gamma(\nu)x^{-\nu}$, and collect $j+(m-j)=m$. This proves both the coefficient formula and the fixed-order remainder. It is the endpoint form of Watson's lemma \cite{dlmflaplace}.
\end{proof}

The proof is local. It does not assert convergence of an infinite endpoint Taylor expansion throughout the whole integration interval. In the Motzkin example we have a stronger globally convergent beta representation; the Delannoy and Schr\"oder applications only need this local lemma.
% ---- hand chunk 05-motzkin-head.tex ----

\section{Motzkin and central trinomial numbers: two exact endpoint blocks}
\label{ct:sec:motzkin}

The Motzkin and central trinomial sequences share one signed-support moment
model, with the exponent $\alpha=\pm\tfrac12$ distinguishing them; the
treatment below keeps $\alpha$ general so that both are handled at once.

% ---- T1 lines 580--716 ----
\subsection{A common moment model}
The Motzkin and central trinomial sequences are A001006 and A002426 \cite{oeismotzkin,oeistrinomial}. Their elementary finite sums are
\begin{equation}
 M_n=\sum_{j=0}^{\lfloor n/2\rfloor}\binom n{2j}\frac1{j+1}\binom{2j}j,
 \qquad
 T_n=\sum_{j=0}^{\lfloor n/2\rfloor}\binom n{2j}\binom{2j}j.
 \label{t1:eq:endpointcounts}
\end{equation}
Both fit the signed-support moment identity
\begin{equation}
 a_n=C_\alpha\int_{-1}^3 t^n\bigl((3-t)(t+1)\bigr)^\alpha\,\ddn t,
 \quad
 (\alpha,C_\alpha)=
 \begin{cases}
 (\tfrac12,\tfrac1{2\pi}),&a_n=M_n,\\
 (-\tfrac12,\tfrac1\pi),&a_n=T_n.
 \end{cases}
 \label{t1:eq:endpointmoment}
\end{equation}
To verify it, set $t=1+s$. For $\alpha=1/2$, the even moments of the normalized semicircle density on $[-2,2]$ are $(j+1)^{-1}\binom{2j}j$; for $\alpha=-1/2$, the even moments of the normalized arcsine density are $\binom{2j}j$. Odd moments vanish. Both even-moment assertions follow by the substitution $s=2\cos\theta$ and a beta integral. Expanding $(1+s)^n$ gives \eqref{t1:eq:endpointcounts}.

For real $x>-1$, define the positive blocks
\begin{align}
 A_+(x)&=C_\alpha\int_0^3t^x\bigl((3-t)(t+1)\bigr)^\alpha\,\ddn t,\nonumber\\
 A_-(x)&=C_\alpha\int_0^1t^x\bigl((3+t)(1-t)\bigr)^\alpha\,\ddn t,
 \label{t1:eq:endpointblocks}
\end{align}
and specify the real interpolation
\begin{equation}
 F_\alpha(x)=A_+(x)+\cos(\pi x)A_-(x).
 \label{t1:eq:endpointinterp}
\end{equation}
This agrees with \eqref{t1:eq:endpointmoment} at every nonnegative integer. The cosine is a deliberate interpolation choice, rather than an assertion that $t^x$ is a real function on the negative half-axis.

\subsection{Exact convergent beta expansions}
Let $p=\alpha+1$. Direct substitutions $t=3(1-u)$ and $t=1-u$ give
\begin{align}
 A_+(x)&=C_\alpha3^{x+p}4^\alpha
 \sum_{j\ge0}\binom\alpha j(-3/4)^j B(x+1,p+j),\nonumber\\
 A_-(x)&=C_\alpha4^\alpha
 \sum_{j\ge0}\binom\alpha j(-1/4)^j B(x+1,p+j).
 \label{t1:eq:betablocks}
\end{align}
Both series are absolutely convergent for $x>-1$. Indeed the binomial series for $(1-cu)^\alpha$ converges absolutely and uniformly on $0\le u\le1$ when $c=3/4$ or $1/4$, and the remaining weight is integrable. Equivalently, each sum is
\begin{equation}
 B(x+1,p)\,{}_2F_1(-\alpha,p;x+p+1;c).
 \label{t1:eq:hyperblock}
\end{equation}
These exact formulas fix each block, including its remainder beyond every finite algebraic order. They are also convenient numerical representations.

\subsection{All-order endpoint coefficients}
For fixed shifts $a,b$, define
\begin{equation}
 G_n(a,b)=\ExpC_n\!\left(
 \left\{\frac{(-1)^{m+1}\bigl(B_{m+1}(a)-B_{m+1}(b)\bigr)}{m(m+1)}\right\}_{m\ge1}
 \right),\qquad G_0(a,b)=1.
 \label{t1:eq:gammaratiocoeff}
\end{equation}
The shifted Stirling expansion proves
$\Gamma(x+a)/\Gamma(x+b)\sim x^{a-b}\sum_{n\ge0}G_n(a,b)x^{-n}$.
For $c\in\{3/4,1/4\}$ set
\begin{equation}
 U_n(\alpha,c)=\sum_{j=0}^n\binom\alpha j(-c)^j(p)_j
 G_{n-j}(1,p+j+1),
 \label{t1:eq:endpointcoeff}
\end{equation}
where $(p)_j=p(p+1)\cdots(p+j-1)$, and define
\begin{equation}
 K_+=C_\alpha3^p4^\alpha\Gamma(p),\qquad
 K_-=C_\alpha4^\alpha\Gamma(p).
 \label{t1:eq:endpointconstants}
\end{equation}
Then the complete algebraic amplitudes are
\begin{align}
 A_+(x)&\sim K_+3^x x^{-p}\sum_{n\ge0}U_n(\alpha,3/4)x^{-n},\nonumber\\
 A_-(x)&\sim K_-x^{-p}\sum_{n\ge0}U_n(\alpha,1/4)x^{-n}.
 \label{t1:eq:endpointasympt}
\end{align}

\begin{proof}[Remainder justification]
For any fixed $N$, Taylor's theorem on $[0,1]$ writes
\[
 (1-cu)^\alpha=\sum_{j<N}\binom\alpha j(-cu)^j+u^N R_N(u),
\]
with $R_N$ bounded because $1-cu\ge1-c>0$. Its integral remainder is bounded by a constant times $B(x+1,p+N)=\OO(x^{-p-N})$. Each of the finitely many retained beta functions has the gamma-ratio expansion used in \eqref{t1:eq:gammaratiocoeff}; collecting powers yields \eqref{t1:eq:endpointcoeff}. Differentiation in $x$ inserts $\log(1-u)$ after extracting $3^x$ from $A_+$. The same endpoint estimate justifies differentiated expansions. The estimates are uniform for bounded imaginary displacements of $x$, since the absolute value of $(1-u)^x$ is $(1-u)^{\Re x}$. Thus the asymptotics used in inverse error estimates are valid, not merely formal manipulations of \eqref{t1:eq:betablocks}.
\end{proof}

For Motzkin numbers, $(K_+,K_-)=(3\sqrt3/(2\sqrt\pi),1/(2\sqrt\pi))$; for central trinomial coefficients, $(K_+,K_-)=(\sqrt3/(2\sqrt\pi),1/(2\sqrt\pi))$. The first four coefficients of the unit amplitudes are
\begin{center}
\begin{tabular}{@{}lrrrr@{}}
\toprule
Block & $U_0$ & $U_1$ & $U_2$ & $U_3$\\
\midrule
Motzkin $+$ & $1$ & $-39/16$ & $2665/512$ & $-87885/8192$\\
Motzkin $-$ & $1$ & $-33/16$ & $1945/512$ & $-57435/8192$\\
Trinomial $+$ & $1$ & $-3/16$ & $1/512$ & $135/8192$\\
Trinomial $-$ & $1$ & $-5/16$ & $49/512$ & $145/8192$\\
\bottomrule
\end{tabular}
\end{center}
Since $A_+'(x)/A_+(x)=\log3-p/x+\OO(x^{-2})$ and $A_-/A_+=\OO(3^{-x})$, the function $F_\alpha$ is positive and strictly increasing for all sufficiently large real $x$.

\subsection{The inverse: an algebraic core and an oscillatory lattice}
First invert $A_+$ exactly and denote that inverse by $g(y)$. Its algebraic expansion is obtained from \eqref{t2:eq:balanced-W}--\eqref{t2:eq:balanced-d} using
\begin{equation}
 a=\log3,\qquad b=-p,\qquad C=K_+,\qquad
 d_n=\LogC_n\bigl(U_1(\alpha,3/4),U_2(\alpha,3/4),\ldots\bigr).
 \label{t1:eq:endpointinverseparams}
\end{equation}
In particular, $c_1=39/(16\log3)$ for Motzkin and $c_1=3/(16\log3)$ for central trinomial coefficients. The common Lambert core is the large solution of
\[
 (\log3)X-p\log X=\log(y/K_+).
\]

Now use the exact relative perturbation
\begin{equation}
 q(x)=\cos(\pi x)\frac{A_-(x)}{A_+(x)}.
 \label{t1:eq:endpointq}
\end{equation}
In \cref{t1:thm:multi}, take $A=A_+$ and this single $q$. This is a complete, convergent expansion in the marker multiplying the second exact block, at sufficiently large $y$. In particular, with $x_0=g(y)$,
\begin{equation}
 F_\alpha^{-1}(y)=x_0-
 \frac{3^{-x_0-p}}{\log3}\cos(\pi x_0)
 +\OO(3^{-x_0}/x_0)+\OO(3^{-2x_0}).
 \label{t1:eq:endpointinversefirst}
\end{equation}
The error is an additive envelope estimate; it does not divide by a cosine that might vanish.

At marker degree $r$, the exponential envelope is $3^{-rx_0}$. Differentiating $\cos^r(\pi x_0)$ produces only frequencies
\[
 r\pi,(r-2)\pi,\ldots,-r\pi.
\]
The quotient of the two unit series in \eqref{t1:eq:endpointasympt}, followed by \eqref{t1:eq:multiinv}, gives every algebraic amplitude of every such harmonic by finite coefficient operations. This describes the full nonlinear interaction lattice for the chosen exact two-block decomposition.

\begin{warning}[An algebraic approximation is not an exact block]
Replacing $x_0=A_+^{-1}(y)$ by a fixed finite algebraic truncation is adequate for algebraic accuracy. Its error generally exceeds $3^{-x_0}$, however. Formula \eqref{t1:eq:endpointinversefirst} is an exponentially accurate statement only when the dominant block is evaluated exactly or with a remainder controlled to the required exponential scale. The exact beta or hypergeometric representations make that distinction operational.
\end{warning}
% ---- hand chunk 05-motzkin-extra.tex ----

\subsection{Higher Motzkin coefficients and the inverse in balanced-core notation}
\label{ct:sec:motzkin-extra}

For the Motzkin case $\alpha=\tfrac12$, write $M_\pm=A_\pm$ and
$a^\pm_m=U_m(\tfrac12,\kappa_\pm)$ with $\kappa_+=\tfrac34$, $\kappa_-=\tfrac14$.
In the endpoint-lemma notation \cref{t2:eq:endpoint-coeff} these are
$a^\pm_m=\sum_{j=0}^m\binom{1/2}{j}(-\kappa_\pm)^j\rising{3/2}{j}R_{m-j}(1,j+\tfrac52)$,
which is \cref{t1:eq:endpointcoeff} with $p=\tfrac32$ and $G_n=R_n$.  One
further coefficient and the normalized ratio of the two blocks are

% ---- T2 lines 907--922 ----
The first coefficients are
\begin{center}
\begin{tabular}{c@{\qquad}rrrrr}
\toprule
 & $m=0$&$m=1$&$m=2$&$m=3$&$m=4$\\\midrule
$a_m^+$&$1$&$-39/16$&$2665/512$&$-87885/8192$&$11422299/524288$\\[2pt]
$a_m^-$&$1$&$-33/16$&$1945/512$&$-57435/8192$&$6977019/524288$\\\bottomrule
\end{tabular}
\end{center}
In particular,
\begin{equation}
 \frac{M_-(x)}{M_+(x)}
 \sim\frac{3^{-x}}{3\sqrt3}\left(1+\frac3{8x}+\cdots\right).
 \label{t2:eq:motzkin-ratio}
\end{equation}
This is a genuine, exactly normalized second sector, not a guessed correction inferred from a recurrence.
% ---- T2 lines 926--953 ----
For the dominant sector, Theorem~\ref{t2:thm:balanced-inverse} applies with
\[
 a=\log3,\quad\beta=-\frac32,\quad C=\frac{3\sqrt3}{2\sqrt\pi},
 \qquad q_j=\cL_j(a^+).
\]
Although $M_+$ is not itself a finite gamma quotient, the proof of that theorem uses only its differentiated all-orders expansion and a positive limiting logarithmic slope. The moment integral supplies these hypotheses. For instance,
\[
 q_1=-\frac{39}{16},\qquad q_2=\frac{143}{64},\qquad
 M_+^{-1}(y)=X+\frac{39}{16aX}
  +\frac1{X^2}\left(\frac{117}{32a^2}-\frac{143}{64a}\right)+\cdots,
\]
where $aX-\tfrac32\log X=\log(y/C)$.

For the small sector, retain the \emph{exact} dominant inverse $x_+=M_+^{-1}(y)$. Taylor expansion of \eqref{t1:eq:endpointinterp} and a derivative bound give
\begin{equation}
 M^{-1}(y)=x_+-\frac{\cos(\pi x_+)M_-(x_+)}{M_+'(x_+)}
                +\bigO(3^{-2x_+}).
 \label{t2:eq:motzkin-inverse-exact}
\end{equation}
Equivalently, with a uniform additive error even at zeros of the cosine,
\begin{equation}
 M^{-1}(y)=x_+-\frac{\cos(\pi x_+)}{3\sqrt3\log3}\,3^{-x_+}
 +\bigO\!\left(\frac{3^{-x_+}}{x_+}\right)
 +\bigO(3^{-2x_+}).
 \label{t2:eq:motzkin-inverse-leading}
\end{equation}
All higher exponential layers and their trigonometric coefficients follow from Theorem~\ref{ct:rmk:transport} with
$\varepsilon(x)=\cos(\pi x)M_-(x)/M_+(x)$. The exact integrals and the convergent beta representations make the meaning of these layers unambiguous. Differentiated endpoint asymptotics imply $M_+'(x)/M_+(x)\to\log3$, while the other term and its derivative are exponentially smaller; hence $M$ is eventually increasing.
% ---- T2 lines 955--1039 ----
\section{Central Delannoy numbers}

Set
\begin{equation}
 a_0=3-2\sqrt2,\qquad b_0=3+2\sqrt2,\qquad
 d=b_0-a_0=4\sqrt2,\qquad \kappa=\frac{b_0}{d}.
 \label{t2:eq:DS-constants}
\end{equation}
The central Delannoy numbers, A001850, have generating function
$(1-6z+z^2)^{-1/2}$ \cite{oeisdelannoy}. A positive real interpolation is
\begin{align}
 D(x)&=\frac1\pi\int_0^\pi(3+2\sqrt2\cos\theta)^x\dd\theta\nonumber\\
 &=\frac1\pi\int_{a_0}^{b_0}
           \frac{t^x}{\sqrt{(b_0-t)(t-a_0)}}\dd t.
 \label{t2:eq:delannoy-moment}
\end{align}
It is also $P_x(3)$ in the usual Legendre-function notation \cite{dlmflegendre}. To check the integer sequence directly, summing the geometric series in the first integral and substituting $u=\tan(\theta/2)$ gives
\[
 \frac1\pi\int_0^\pi\frac{\dd\theta}{1-z(3+2\sqrt2\cos\theta)}
 =\frac1{\sqrt{(1-a_0z)(1-b_0z)}}.
\]

The endpoint $t=b_0$ yields
\begin{equation}
 D(x)\sim C_D b_0^x x^{-1/2}\sum_{m\ge0}d_mx^{-m},
 \qquad C_D=\sqrt{\frac{b_0}{\pi d}},
 \label{t2:eq:delannoy-asymp}
\end{equation}
where
\begin{equation}
 d_m=\sum_{j=0}^m\binom{-1/2}{j}(-\kappa)^j
             \rising{1/2}{j}R_{m-j}(1,j+\tfrac32).
 \label{t2:eq:delannoy-coeff}
\end{equation}
For example,
\[
 d_0=1,\qquad d_1=-\frac14+\frac{3\sqrt2}{32},\qquad
 d_2=\frac{113}{1024}-\frac{9\sqrt2}{128}.
\]
The inverse has the balanced core with $a=\log b_0$, $\beta=-1/2$, $C=C_D$, and coefficients $q_j=\cL_j(d)$. Its first displacement is
\[
 D^{-1}(y)=X+\frac{1/4-3\sqrt2/32}{(\log b_0)X}+\bigO(X^{-2}).
\]
All orders follow from the two finite coefficient formulas \eqref{t2:eq:delannoy-coeff} and \eqref{t2:eq:balanced-d}.

\section{Large Schr\"oder numbers}

The large Schr\"oder numbers $1,2,6,22,90,\ldots$, A006318, have generating function
$(1-z-\sqrt{1-6z+z^2})/(2z)$ \cite{oeisschroder}. Define
\begin{equation}
 S(x)=\frac1{2\pi}\int_{a_0}^{b_0}
                t^{x-1}\sqrt{(b_0-t)(t-a_0)}\dd t.
 \label{t2:eq:schroder-moment}
\end{equation}
This integral exists for every real $x$ because $a_0>0$. Its mass at $x=0$ is one: the elementary integral of the square-root density divided by $t$ is
$\pi((a_0+b_0)/2-\sqrt{a_0b_0})=2\pi$. Applying the same elementary Stieltjes-transform calculation as above to
$1/[t(1-zt)]=1/t+z/(1-zt)$ gives the stated generating function. Thus $S(n)$ is the desired integer sequence.

Endpoint extraction now yields
\begin{equation}
 S(x)\sim C_S b_0^xx^{-3/2}\sum_{m\ge0}s_mx^{-m},
 \qquad C_S=\frac{\sqrt{b_0d}}{4\sqrt\pi},
 \label{t2:eq:schroder-asymp}
\end{equation}
with
\begin{equation}
 s_m=\sum_{j=0}^m\binom{1/2}{j}(-\kappa)^j
                  \rising{3/2}{j}R_{m-j}(0,j+\tfrac32).
 \label{t2:eq:schroder-coeff}
\end{equation}
The shift zero in the gamma ratio is significant: the moment integrand is $t^{x-1}$, not $t^x$. The first coefficients are
\[
 s_0=1,\qquad s_1=-\frac34-\frac{9\sqrt2}{32},\qquad
 s_2=\frac{665}{1024}+\frac{45\sqrt2}{128}.
\]
Again the inverse is given to all orders by $a=\log b_0$, $\beta=-3/2$, $C=C_S$, and $q_j=\cL_j(s)$. In particular,
\[
 S^{-1}(y)=X+\frac{3/4+9\sqrt2/32}{(\log b_0)X}+\bigO(X^{-2}).
\]

\section{A boundary on the endpoint interpretation}

Both $D(x)$ and $S(x)$ have exact positive-support moment definitions. These definitions fix their real inverses and provide differentiated all-orders expansions. They do \emph{not}, without further analysis, justify adding a nonzero lower-endpoint exponential sector merely because the generating function has a second algebraic singularity or a recurrence has a second characteristic root. A solution of the recurrence can have a zero coefficient in that formal mode, and endpoint decompositions can depend on a contour or a summation convention.

Accordingly, the Delannoy and Schr\"oder results above claim the full algebraic sector and its inverse, not a computed hyperasymptotic decomposition. Motzkin numbers are different: splitting an actual real moment integral at zero produced two exact functions before any expansion was made. That stronger input is what licensed the explicit exponentially small inverse term.
% ---- T1 lines 1181--1227 ----
\section{General signed-support moment sequences}
The Motzkin and trinomial construction extends beyond their particular densities. Suppose $a>b>0$, $a>1$, and $w$ is a nonnegative integrable density on $[-b,a]$. Assume that, near the two outer endpoints,
\begin{align}
 w(a(1-u))&=u^{p_+-1}\sum_{j\ge0}h_{+,j}u^j,\nonumber\\
 w(-b(1-u))&=u^{p_--1}\sum_{j\ge0}h_{-,j}u^j,
 \label{t1:eq:generalendpoints}
\end{align}
where $p_\pm>0$, the local amplitude series converge, and $h_{\pm,0}>0$. Define the exact moment blocks
\[
 A_+(x)=\int_0^a t^xw(t)\,\ddn t,\qquad
 A_-(x)=\int_0^b t^xw(-t)\,\ddn t,
\]
and $F(x)=A_+(x)+\cos(\pi x)A_-(x)$. Its integer values are the moments of $w$ on $[-b,a]$.

\begin{proposition}[Endpoint transfer with explicit coefficients]\label{t1:prop:endpointtransfer}
Put $K_+=a h_{+,0}\Gamma(p_+)$ and $K_-=b h_{-,0}\Gamma(p_-)$. Then
\begin{align}
 A_+(x)&\sim K_+a^xx^{-p_+}\sum_{n\ge0}u_{+,n}x^{-n},\nonumber\\
 A_-(x)&\sim K_-b^xx^{-p_-}\sum_{n\ge0}u_{-,n}x^{-n},
 \label{t1:eq:generalmomentblocks}
\end{align}
with $u_{\pm,0}=1$ and
\begin{equation}
 u_{\pm,n}=\sum_{j=0}^n
 \frac{h_{\pm,j}}{h_{\pm,0}}(p_\pm)_j
 G_{n-j}(1,p_\pm+j+1).
 \label{t1:eq:generalmomentcoeff}
\end{equation}
The interpolation is eventually increasing, and its relative small block is
\begin{equation}
 \frac{A_-(x)}{A_+(x)}\sim
 \frac{K_-}{K_+}\left(\frac ba\right)^x
 x^{p_+-p_-}\frac{\sum u_{-,n}x^{-n}}{\sum u_{+,n}x^{-n}}.
 \label{t1:eq:generalmomentr}
\end{equation}
Its inverse is obtained by \eqref{t2:eq:balanced-d} for $A_+$ and \eqref{t1:eq:multiinv} for the exact multiplicative correction. The first displacement, with $x_0=A_+^{-1}(y)$, is
\begin{equation}
 -\frac{K_-/K_+}{\log a}
 \left(\frac ba\right)^{x_0}x_0^{p_+-p_-}\cos(\pi x_0)
\end{equation}
up to a relative algebraic amplitude error interpreted additively at cosine zeros, and terms of second exponential order.
\end{proposition}
\begin{proof}
Use $t=a(1-u)$ and $t=b(1-u)$ near the respective endpoints. A finite local Taylor expansion of the density amplitude reduces every retained term to a beta integral, with an error of the next algebraic order as in \cref{ct:ch:moments}. The part of either integral a fixed distance from its outer endpoint is exponentially smaller than its own leading block and therefore smaller than each fixed algebraic remainder. Applying the gamma-ratio expansion yields \eqref{t1:eq:generalmomentcoeff}. The dominant logarithmic derivative tends to $\log a>0$, and the relative negative block and its derivative tend to zero exponentially because $b/a<1$. The inverse assertions follow from the two inversion theorems.
\end{proof}

The exact integrals are part of the statement. The endpoint amplitude coefficients alone need not determine all exponentially small contributions within either block, particularly when the density has additional interior structure. The proposition is a reusable block decomposition and all-order endpoint theorem, not a classification of every complex saddle or every interior singularity.
% ---- hand chunk 06-part5-head.tex ----

\part{Moving saddles}
\label{ct:part:saddles}

\chapter{Involutions: two real saddles and a stretched-exponential inverse}
\label{ct:ch:involutions}

The involution numbers are Gaussian moments $\mathbb E(1+Z)^n$, and splitting
the Gaussian integral at zero produces two exact positive functions, each
with a unique saddle that moves like $\sqrt x$ while its width stays of order
one.  Two finite formulas give every saddle coefficient; the inverse of the
dominant block has an \emph{unbounded} first correction, of order
$\sqrt X/\log X$, before its half-power layers; and the second block
produces a stretched-exponential inverse lattice with envelope
$\ee^{-2\sqrt{x}}$.  The chapter closes with a general rule for predicting
the scale of such a drift.

% ---- T2 lines 1044--1070 ----
\section{An exact Gaussian decomposition}

An involution consists of fixed points and disjoint transpositions. Choosing $j$ transpositions gives
\begin{equation}
 I_n=\sum_{j=0}^{\lfloor n/2\rfloor}
        \frac{n!}{2^jj!(n-2j)!},\qquad
 \sum_{n\ge0}I_n\frac{z^n}{n!}=\ee^{z+z^2/2}.
 \label{t2:eq:involution-enumeration}
\end{equation}
These are the involution or telephone numbers, A000085 \cite{oeisinvolution}. If $Z$ is a standard normal random variable, its exponential moment gives $I_n=\mathbb E(1+Z)^n$. Equivalently, define for $x>-1$
\begin{equation}
 A_\sigma(x)=\frac1{\sqrt{2\pi}}\int_0^\infty
          t^x\exp\!\left(-\frac{(t-\sigma)^2}{2}\right)\dd t,
 \qquad \sigma\in\{+1,-1\}.
 \label{t2:eq:involution-A}
\end{equation}
Splitting the Gaussian moment at zero gives
\begin{equation}
 I_n=A_+(n)+(-1)^nA_-(n).
 \label{t2:eq:involution-exact-lattice}
\end{equation}
We select the real interpolation
\begin{equation}
 I(x)=A_+(x)+\cos(\pi x)A_-(x).
 \label{t2:eq:involution-cont}
\end{equation}
Both $A_\sigma$ are positive exact functions. They are therefore suitable sector normalizations: no subdominant constant is being selected by truncating a divergent dominant series.
% ---- T2 lines 1072--1168 ----
\section{The moving saddle}

For fixed $x$, the logarithm of the integrand in \eqref{t2:eq:involution-A} is
\[
 h_\sigma(t)=x\log t-\frac{(t-\sigma)^2}{2}.
\]
Its unique maximum occurs at
\begin{equation}
 r_\sigma=\frac{\sigma+\sqrt{1+4x}}2,
 \qquad r_\sigma^2-\sigma r_\sigma=x,
 \label{t2:eq:involution-saddle}
\end{equation}
and its curvature is
\begin{equation}
 b_\sigma=-h_\sigma''(r_\sigma)
          =2-\frac{\sigma}{r_\sigma}.
 \label{t2:eq:involution-curvature}
\end{equation}
In particular, the saddle moves like $\sqrt x$, but its Gaussian width remains of order one.

Expand at $r=r_\sigma$ using $u=t-r$. For $k\ge3$,
\[
 \frac{h_\sigma^{(k)}(r)}{k!}
       =\frac{(-1)^{k-1}x}{kr^k}.
\]
Thus the phase has the form
\[
 h_\sigma(r)-\frac{b_\sigma u^2}{2}
       +\sum_{k\ge3}\frac{(-1)^{k-1}x}{kr^k}u^k.
\]
Multiplying these Taylor terms and integrating Gaussian moments is enough to compute every coefficient. The following theorem turns that procedure into a single finite coefficient formula.

\begin{theorem}[All coefficients of both involution saddles]
\label{t2:thm:involution-coeff}
Put $t=x^{-1/2}$ and, in the formulas below, use $t$ as a formal variable. Define
\begin{equation}
 R_\sigma(t)=\frac{\sigma t+\sqrt{4+t^2}}2,
 \qquad b_\sigma(t)=2-\frac{\sigma t}{R_\sigma(t)}.
 \label{t2:eq:involution-R}
\end{equation}
For a finitely supported vector $(m_3,m_4,\ldots)$ of nonnegative integers, set
$K=\sum_{k\ge3}km_k$ and $M=\sum_{k\ge3}m_k$. Define the formal Gaussian correction
\begin{equation}
 \mathcal S_\sigma(t)=
 \sum_{\substack{(m_k)\ge0\\K\ \mathrm{even}}}
 \frac{(-1)^M(K-1)!!}{\prod_{k\ge3}m_k!k^{m_k}}
 \frac{t^{K-2M}}{R_\sigma(t)^K b_\sigma(t)^{K/2}},
 \label{t2:eq:involution-S}
\end{equation}
where the empty vector contributes $1$ and $(-1)!!=1$. Finally put
\begin{align}
 B_\sigma(t)={}&
 \exp\!\left(\frac{\log R_\sigma(t)}{t^2}
       +\frac{\sigma R_\sigma(t)}{2t}-\frac\sigma t-\frac14\right)
 \sqrt{\frac2{b_\sigma(t)}}\,\mathcal S_\sigma(t).
 \label{t2:eq:involution-B}
\end{align}
The apparent negative powers in the exponent cancel. Every coefficient $[t^N]B_\sigma(t)$ is a finite sum: only vectors with $K-2M\le N$ are needed. If $B(t)=B_+(t)$, then $B_-(t)=B(-t)$, and
\begin{equation}
 A_\sigma(x)\sim
 \frac{\ee^{-1/4}}{\sqrt2}\left(\frac{x}{\ee}\right)^{x/2}
          \ee^{\sigma\sqrt x}B(\sigma x^{-1/2}).
 \label{t2:eq:involution-saddle-expansion}
\end{equation}
\end{theorem}

\begin{proof}
The normalized Gaussian moment of $u^K$ is zero for odd $K$ and equals $(K-1)!!\,b_\sigma^{-K/2}$ for even $K$. For a vector $(m_k)$, the product of Taylor coefficients contributes
\[
 \prod_{k\ge3}\frac1{m_k!}
       \left(\frac{(-1)^{k-1}x}{kr^k}\right)^{m_k}.
\]
When $K$ is even, its sign is $(-1)^{K-M}=(-1)^M$. Since $x=t^{-2}$ and $r=R_\sigma/t$, this gives exactly \eqref{t2:eq:involution-S}. Each nonzero $m_k$ contributes positive weight $k-2$, so extraction at a fixed power is finite.

The saddle identity $r^2-\sigma r=x$ simplifies the phase to
\[
 h_\sigma(r)=x\log r-\frac x2+\frac{\sigma r}{2}-\frac12.
\]
Dividing $\ee^{h_\sigma(r)}/\sqrt{b_\sigma}$ by the leading factor in \eqref{t2:eq:involution-saddle-expansion} yields the exponential and square root in \eqref{t2:eq:involution-B}. The expansions $R_\sigma=1+\sigma t/2+t^2/8+\cdots$ and $\log R_\sigma=\sigma\operatorname{arsinh}(t/2)$ show cancellation of the negative powers. Replacing $(\sigma,t)$ by $(+,\sigma t)$ proves the sign symmetry; $K-2M$ is even whenever $K$ is even.

For a rigorous fixed-order remainder, note that
$h_\sigma''(v)=-x/v^2-1\le-1$ for $v>0$. The phase is globally strictly concave, and the contribution outside $|v-r_\sigma|\le x^\epsilon$ is bounded relative to its maximum by a Gaussian tail $\bigO(\ee^{-c x^{2\epsilon}})$. On the inner interval, Taylor expansion to a sufficiently high finite order is uniform; choose $\epsilon>0$ small relative to the desired order. Its remainder is a polynomially bounded factor times a higher power of $x^{-1/2}$, integrable against a slightly weaker Gaussian. Extending the local Gaussian integral to the real line changes it by another negligible tail. This proves the expansion to every fixed order. The same argument after finitely many differentiations supplies the differentiated expansions used in inversion. This is a moving-saddle version of Laplace's method \cite{dlmflaplace}.
\end{proof}

The first terms are
\begin{align}
 B(t)={}&1+\frac7{24}t-\frac{119}{1152}t^2
       -\frac{7933}{414720}t^3
       +\frac{1967381}{39813120}t^4
       -\frac{57200419}{1337720832}t^5+\cdots,
 \label{t2:eq:involution-B-first}\\
 \log B(t)={}&\frac7{24}t-\frac7{48}t^2
      +\frac{37}{1920}t^3+\frac5{128}t^4
      -\frac{5879}{107520}t^5+\cdots.
 \label{t2:eq:involution-logB}
\end{align}
The amplitude coefficients in \eqref{t2:eq:involution-B-first} agree with the independent high-order expansion of Ekhad, Kauers, and Zeilberger \cite{ekz}. They are included here as checks of the general saddle formula, not as a claim of priority for those numerical coefficients.
% ---- hand chunk 06-inv-second.tex ----

\section{A second finite formula: Gaussian moments of the shifted saddle}
\label{ct:sec:inv-second}

The partition sum \cref{t2:eq:involution-S} organizes the saddle expansion
by the Taylor coefficients of the phase.  Centring the integration variable
at $\sqrt x$ instead of at the exact saddle $r_\sigma$ gives a second finite
formula for the same amplitude, in terms of the moments of a Gaussian with
mean $\tfrac12$ and variance $\tfrac12$; the two formulas were verified to
agree through $t^6$ in exact rational arithmetic
(\cref{ct:ch:numerics,ct:app:audit}).  Throughout this section the
amplitude is written $B(t)=\sum_{m\ge0}b_mt^m$ as in
\cref{t2:eq:involution-B}, and $H(t)=\log B(t)=\sum_{m\ge1}h_mt^m$.

% ---- T1 lines 748--809 ----
Put $\varepsilon=x^{-1/2}$ and define
\begin{equation}
 \mu_N=\sum_{j=0}^{\lfloor N/2\rfloor}
 \frac{N!(1/2)^{N-2j}}{4^j j!(N-2j)!}.
 \label{t1:eq:ivmu}
\end{equation}
This is the $N$th moment of a Gaussian variable with mean $1/2$ and variance $1/2$. Let $b_0=1$, and, for $m\ge1$, set
\begin{equation}
 b_m=\sum_{\substack{k_1,\ldots,k_m\ge0\\\sum_{r=1}^m r k_r=m}}
 \mu_{\sum_{r=1}^m(r+2)k_r}
 \prod_{r=1}^m\frac{1}{k_r!}
 \left(\frac{(-1)^{r+1}}{r+2}\right)^{k_r}.
 \label{t1:eq:ivcoeff}
\end{equation}

\begin{theorem}[All-order involution blocks]
\label{t1:thm:ivforward}
For each fixed integer $M\ge0$, as $x\to+\infty$,
\begin{equation}
 A_\sigma(x)=\frac1{\sqrt2}
 \exp\left(\frac x2(\log x-1)+\sigma\sqrt x-\frac14\right)
 \left(\sum_{m=0}^M b_m\sigma^m x^{-m/2}
 +\OO(x^{-(M+1)/2})\right).
 \label{t1:eq:ivforward}
\end{equation}
The expansion may be differentiated a fixed number of times, with the corresponding differentiated asymptotic orders. In particular, the relative small block has the all-order amplitude
\begin{equation}
 \frac{A_-(x)}{A_+(x)}\sim
 \ee^{-2\sqrt x}\frac{B(-x^{-1/2})}{B(x^{-1/2})},
 \qquad B(t)=\sum_{m\ge0}b_mt^m.
 \label{t1:eq:ivratio}
\end{equation}
The series $A$ in this statement is an asymptotic series, not a claimed convergent representation of the exact integrals.
\end{theorem}
\begin{proof}
Set $t=\sqrt x+v$ in \eqref{t2:eq:involution-A}. The exponent becomes
\begin{align*}
 x\log t-\frac{(t-\sigma)^2}{2}
 ={}&\frac x2(\log x-1)+\sigma\sqrt x-\frac12-v^2+\sigma v\\
 &+\sum_{r\ge1}\frac{(-1)^{r+1}}{r+2}\varepsilon^r v^{r+2}.
\end{align*}
Completing the square gives $-v^2+\sigma v=-(v-\sigma/2)^2+1/4$. The Gaussian integral of this leading local factor is $\sqrt\pi$, accounting for $1/\sqrt2$ and $\ee^{-1/4}$ in \eqref{t1:eq:ivforward}. Expanding the last exponential and collecting $\varepsilon^m$ gives a finite sum indexed by $\sum rk_r=m$. The power of $v$ is $N=\sum(r+2)k_r$. Gaussian moments with mean $\sigma/2$ are $\sigma^N\mu_N$, and $N\equiv m\pmod2$. This yields \eqref{t1:eq:ivcoeff} and the sign $\sigma^m$.

Here is the finite-order justification of the local expansion. For fixed $M$, restrict first to $|v|\le x^\eta$, where $\eta>0$ is sufficiently small, depending on $M$. Taylor's theorem for $\log(1+v/\sqrt x)$ has a remainder bounded by a constant times the next power, and the exponent outside its quadratic part is $o(v^2)+\OO(1)$. Expanding to an order larger than $M$ and integrating against a Gaussian majorant bounds the integrated remainder by $\OO(\varepsilon^{M+1})$. On the rest of the positive integration interval, the phase is strictly concave, with second derivative $-x/t^2-1\le-1$; its saddle differs from $\sqrt x$ by a bounded amount. Hence the omitted tails are bounded by a Gaussian tail $\OO(\ee^{-c x^{2\eta}})$ relative to the saddle maximum. This is smaller than any prescribed algebraic power. The lower limit $v=-\sqrt x$ can therefore be replaced by $-\infty$ at finite asymptotic order. Factors $\log t$ introduced by differentiation are treated by the same estimates and local expansion. For a bounded imaginary displacement, center the real integration variable at $\sqrt{\Re x}$ and expand the extra factor $\exp(\ii\Im(x)\log t)$; the same Gaussian bounds give estimates uniform on fixed-width complex neighborhoods of the real tail. Division of the two unit asymptotic series proves \eqref{t1:eq:ivratio}.
\end{proof}

The first coefficients are
\begin{align}
 B(t)={}&1+\frac7{24}t-\frac{119}{1152}t^2
 -\frac{7933}{414720}t^3+\frac{1967381}{39813120}t^4\nonumber\\
 &-\frac{57200419}{1337720832}t^5
 +\frac{6340449533}{687970713600}t^6+\OO(t^7).
 \label{t1:eq:ivfirstcoeff}
\end{align}
For later use put $H(t)=\log B(t)=\sum_{m\ge1}h_mt^m$, interpreted formally. Then $h_m=\LogC_m(b_1,b_2,\ldots)$, and
\begin{equation}
 h_1=\frac7{24},\quad h_2=-\frac7{48},\quad
 h_3=\frac{37}{1920},\quad h_4=\frac5{128},\quad
 h_5=-\frac{5879}{107520},\quad h_6=\frac{43}{1440}.
 \label{t1:eq:ivlogcoeff}
\end{equation}
These rational coefficients and their recurrence audit are reproduced in the computational supplement and \cref{t1:app:audit}.
% ---- T2 lines 1170--1177 ----
Combining the exact decomposition and the two expansions gives
\begin{equation}
 I(x)\sectorsim
 \frac{\ee^{-1/4}}{\sqrt2}\left(\frac{x}{\ee}\right)^{x/2}\ee^{\sqrt x}
 \left\{B(x^{-1/2})+\cos(\pi x)\ee^{-2\sqrt x}B(-x^{-1/2})\right\}.
 \label{t2:eq:involution-transseries}
\end{equation}
The second sector is stretched-exponential relative to the first. At integers its sign alternates. A finite truncation of $B$ has an error much larger than that second sector; the exact functions $A_\pm$ retain its meaning despite this fact.
% ---- hand chunk 06-inv-inverse.tex ----

\section{The inverse correction is not bounded}
\label{ct:sec:inv-inverse}

% ---- T1 lines 812--879 ----
For the positive block $A_+$, introduce
\begin{equation}
 L=2\log y+\frac12+\log2,\qquad
 X=\frac{L}{W_0(L/\ee)},\qquad \lambda=\log X.
 \label{t1:eq:ivcore}
\end{equation}
The carrier satisfies $X(\log X-1)=L$. The positive-block inverse obeys, to every fixed asymptotic order,
\begin{equation}
 x(\log x-1)+2\sqrt x+2H(x^{-1/2})=L.
 \label{t1:eq:ivinveq}
\end{equation}
The first residual at $X$ is $2\sqrt X$. Dividing by the slope $\log X$ shows already that the leading correction is $-2\sqrt X/\log X$, which tends to $-\infty$. A fixed-point ansatz $x=X+\OO(X^{-1})$ would therefore fail at its first step.

Set $x=X+\sqrt X\,V$, $\varepsilon=X^{-1/2}$, and define
$G(u)=(1+u)\log(1+u)-u$. After subtracting the carrier equation and dividing by $\sqrt X$, \eqref{t1:eq:ivinveq} becomes
\begin{equation}
 \lambda V+\frac{G(\varepsilon V)}{\varepsilon}
 +2\sqrt{1+\varepsilon V}
 +2\varepsilon H\!\left(\varepsilon(1+\varepsilon V)^{-1/2}\right)=0.
 \label{t1:eq:ivVeq}
\end{equation}
The apparent division by $\varepsilon$ is harmless because $G(u)$ starts with $u^2/2$.

\begin{theorem}[All-order inverse coefficients for the dominant involution block]\label{t1:thm:ivinverse}
There is a unique formal solution
\begin{equation}
 V(\varepsilon,\lambda)=\sum_{j\ge0}v_j(\lambda)\varepsilon^j,
 \qquad v_j\in\QQ(\lambda),\qquad v_0=-2/\lambda.
 \label{t1:eq:ivV}
\end{equation}
It gives the asymptotic expansion of $A_+^{-1}(y)$ through
$x=X+\sqrt X V(X^{-1/2},\log X)$. To specify every coefficient without recurrence, put $v_0=-2/\lambda$ and
\begin{align}
 \Xi(\varepsilon,z)=-\frac1\lambda\Bigg\{&
 \frac{G(\varepsilon(v_0+z))}{\varepsilon}
 +2\sqrt{1+\varepsilon(v_0+z)}-2\nonumber\\
 &+2\varepsilon H\!\left(\varepsilon(1+\varepsilon(v_0+z))^{-1/2}\right)\Bigg\}.
 \label{t1:eq:ivXi}
\end{align}
Then, for $j\ge1$,
\begin{equation}
 v_j(\lambda)=\sum_{m=1}^j\frac1m
 \coefA{\varepsilon^j z^{m-1}}\Xi(\varepsilon,z)^m.
 \label{t1:eq:ivVclosed}
\end{equation}
\end{theorem}
\begin{proof}
Since $\lambda v_0+2=0$, \eqref{t1:eq:ivVeq} is equivalent to $z=\Xi(\varepsilon,z)$ for $z=V-v_0$. The series $\Xi$ belongs to $\varepsilon\QQ(\lambda)[[\varepsilon,z]]$. Formal Lagrange inversion proves uniqueness and \eqref{t1:eq:ivVclosed}. Only $h_1,\ldots,h_{j-1}$ can enter $v_j$, so the formula is finite when combined with \eqref{t1:eq:ivcoeff} and \eqref{t2:eq:L}.

For asymptotic validity, use a finite number of terms of \cref{t1:thm:ivforward} in the logarithmic equation and substitute the truncated $V$. Coefficient cancellation leaves a residual of the next half-power order, with rational functions of $\lambda$. For $\lambda=\log X$ large these are bounded by fixed powers of $\lambda$. The derivative of the unscaled logarithmic phase is $\log X+\OO(X^{-1/2})$, bounded below by a positive multiple of $\log X$. The mean value theorem converts the phase residual into the next inverse error. Equivalently, on $|V|\le C/\lambda$, the derivative of the left side of \eqref{t1:eq:ivVeq} in $V$ is $\lambda+\OO(\varepsilon)$, giving a uniform normalized-slope estimate. This proves the expansion in successive powers of $\varepsilon$ with rational-logarithmic coefficients.
\end{proof}

The first corrections are
\begin{align}
 v_0&=-\frac2\lambda,\qquad
 v_1=\frac{2(\lambda-1)}{\lambda^3},\nonumber\\
 v_2&=-\frac{7\lambda^4+12\lambda^2-56\lambda+48}{12\lambda^5},\label{t1:eq:ivvfirst}\\
 v_3&=\frac{7\lambda^6-28\lambda^4-96\lambda^2+320\lambda-240}{24\lambda^7}.
 \nonumber
\end{align}
Thus the first two displacements from $X$ are
\begin{equation}
 A_+^{-1}(y)=X-\frac{2\sqrt X}{\log X}
 +\frac{2(\log X-1)}{(\log X)^3}
 +\OO\!\left(\frac1{\sqrt X\log X}\right).
 \label{t1:eq:ivinvfirst}
\end{equation}
The last error follows from the displayed $v_2$ and the next-order residual estimate. More generally, truncating \eqref{t1:eq:ivV} through $v_N$ gives, for fixed $N\ge1$, the conservative absolute error $\OO(X^{-N/2})$.
% ---- hand chunk 06-inv-z.tex ----

\begin{remark}[The same expansion in half-powers of $X$]
\label{ct:rmk:zj}
Expanding \cref{t1:eq:ivV} in $X$ gives
$x_+\sim X+\sum_{j\ge0}z_j(\ell)X^{(1-j)/2}$ with $\ell=\log X$ and
$z_j=v_j$; the second source derived the same coefficients in this form,
$z_0=-2/\ell$, $z_1=2/\ell^2-2/\ell^3$,
$z_2=-7/(12\ell)-1/\ell^3+14/(3\ell^4)-4/\ell^5$, which are
\cref{t1:eq:ivvfirst} after expansion, by the scaled reversion
\cref{t2:prop:scaled} with $x=X(1+tz)$, $t=X^{-1/2}$, and the finite
extraction $z_j=\sum_{m=1}^{j}\frac1m[t^ju^{m-1}]\Phi(t,u)^m$ of the
displaced equation.  Using $z_0,z_1,z_2$ leaves an error $\bigO(X^{-1})$;
the symbolic audit of \cref{ct:ch:numerics} substitutes them into the
displaced equation and verifies exact cancellation through $t^2$.
\end{remark}

% ---- T1 lines 881--901 ----
\section{The exponentially small inverse lattice}
The exact dominant-block inverse is now denoted $x_0=A_+^{-1}(y)$. Since
\[
 (\log A_+)'(x)=\frac12\log x+\frac1{2\sqrt x}+\OO(x^{-3/2}),
\]
and $A_-/A_+=\OO(\ee^{-2\sqrt x})$, the interpolation \eqref{t2:eq:involution-cont} is eventually positive and strictly increasing. Apply \cref{t1:thm:multi} with
\begin{equation}
 q(x)=\cos(\pi x)\frac{A_-(x)}{A_+(x)}.
 \label{t1:eq:ivq}
\end{equation}
It gives the exact analytic marker expansion about $x_0$; its coefficients have the half-power amplitudes obtained from \eqref{t1:eq:ivratio}. The first sector is
\begin{align}
 I^{-1}(y)={}&x_0-
 \frac{2\ee^{-2\sqrt{x_0}}}{\log x_0}\cos(\pi x_0)
 +\OO\!\left(\frac{\ee^{-2\sqrt{x_0}}}{\sqrt{x_0}\log x_0}\right)\nonumber\\
 &+\OO\!\left(\frac{\ee^{-4\sqrt{x_0}}}{\log x_0}\right).
 \label{t1:eq:ivexpinverse}
\end{align}
At marker degree $r$, all terms have the envelope $\ee^{-2r\sqrt{x_0}}$ and harmonic frequencies $r\pi,(r-2)\pi,\ldots,-r\pi$. Differentiation also generates powers of $x_0^{-1/2}$ and $(\log x_0)^{-1}$. Formula \eqref{t1:eq:multiinv}, with the exact $q$ in \eqref{t1:eq:ivq}, is an all-sector formula including these interactions.

This is a different transseries scale from the endpoint examples: the small parameter is $\ee^{-2\sqrt{x_0}}$, not $\ee^{-c x_0}$. It is also different from the Gaussian-binomial example below, whose \emph{entire inverse} grows only like $\sqrt{\log y}$. The placement of a square root in the forward phase or in the inverse carrier cannot be interchanged.
% ---- T2 lines 1257--1264 ----
The distinction between $X$ and $x_+$ matters even inside the action. Since
$x_+-X\sim-2\sqrt X/\log X$,
\[
 \ee^{-2\sqrt{x_+}}
 =\ee^{-2\sqrt X}\exp\left(\frac2{\log X}
             +\bigO\!\left(\frac1{\sqrt X(\log X)^2}\right)\right).
\]
Thus replacing $x_+$ by $X$ without expanding the induced logarithmic correction discards an entire part of the sector amplitude. Replacing $\cos(\pi x_+)$ by $\cos(\pi X)$ is even more serious, because $x_+-X$ is unbounded. One must retain or expand the full phase, not just the exponential decay rate.
% ---- T1 lines 1229--1250 ----
\section{Predicting the size of an inverse drift}
A second useful principle concerns the scale of the first inverse correction. Consider a phase
\begin{equation}
 \Phi_0(x)=c x^\rho(\log x-\eta),\qquad c>0,\quad\rho>0.
 \label{t1:eq:powerlogphase}
\end{equation}
The large solution of $\Phi_0(X)=L$ is
\begin{equation}
 X=\left\{\frac{\rho L}{cW_0((\rho/c)\ee^{-\rho\eta}L)}\right\}^{1/\rho}.
 \label{t1:eq:powerlogcore}
\end{equation}
Indeed $w=\rho(\log X-\eta)$ satisfies $w\ee^w=(\rho/c)\ee^{-\rho\eta}L$.
Suppose the additional phase is
$R(x)\sim d x^\sigma(\log x)^\tau$ with $0\le\sigma<\rho$, and its derivatives have the corresponding orders. Linearizing gives
\begin{equation}
 \delta_1=-\frac{R(X)}{\Phi_0'(X)}
 \sim-\frac d{c\rho}X^{\sigma-\rho+1}(\log X)^{\tau-1}.
 \label{t1:eq:driftscale}
\end{equation}
The condition $R=o(\Phi_0)$ only guarantees $\delta_1=o(X)$; it does not guarantee $\delta_1=o(1)$. Formula \eqref{t1:eq:driftscale} identifies the correct shift scale before any long coefficient calculation. The involution case has $(c,\rho,\eta,d,\sigma,\tau)=(1/2,1,1,1,1/2,0)$ and recovers $-2\sqrt X/\log X$.

Once the scale is chosen, write $x=X+s(X)V$ with $s(X)$ matching \eqref{t1:eq:driftscale} or its algebraic part, normalize by the phase slope, and apply the fixed-point coefficient extraction used in \eqref{t1:eq:ivVclosed}. The scale must also be small enough relative to $X$ for the normalized Taylor expansion to be asymptotically ordered. This gives a practical extension rule to further combinatorial saddle-point problems with lower-order power--logarithmic phases.
% ---- hand chunk 07-part6-head.tex ----

\part{Arithmetic sectors}
\label{ct:part:arithmetic}

\chapter{Alternating permutations: a gamma block and odd-integer sectors}
\label{ct:ch:zigzag}

The Euler zigzag numbers are exactly a gamma block times a Dirichlet series
over the odd integers, with the parity of the index selecting the character.
The gamma block is inverted by a centred Lambert carrier whose algebraic
corrections have only odd orders; the Dirichlet tail is transported into
inverse sectors indexed by odd integers, whose coefficients are finite
ordered-factorization sums, and equivalently by odd prime powers through
the Euler product.  The two parity sheets are inverted separately and the
integer threshold is reconstructed from both.

% ---- T1 lines 906--1055 ----
\section{Exact parity-resolved continuations}
Let $Z_n$ denote the Euler zigzag numbers, A000111, normalized by
\begin{equation}
 \sum_{n\ge0}Z_n\frac{z^n}{n!}=\sec z+\tan z
 =\tan(\pi/4+z/2).
 \label{t1:eq:zigzagegf}
\end{equation}
They count permutations with a prescribed alternating orientation; using both orientations would double the values for $n\ge2$ \cite{oeiszigzag}. Set
\begin{equation}
 a=\pi/2,\qquad \mathcal G(x)=2\Gamma(x+1)a^{-x-1},\qquad s=x+1.
 \label{t1:eq:zigzaggamma}
\end{equation}
Define, for $x>0$,
\begin{align}
 Z_0(x)&=\mathcal G(x)\beta(x+1),\nonumber\\
 Z_1(x)&=\mathcal G(x)(1-2^{-x-1})\zeta(x+1),
 \label{t1:eq:zigzaginterp}
\end{align}
where $\beta(s)=\sum_{j\ge0}(-1)^j(2j+1)^{-s}$ is Dirichlet beta \cite{wlbeta}. The subscripts $0$ and $1$ refer to even and odd integer arguments, respectively.

\begin{proposition}[Exact odd-integer forward blocks]\label{t1:prop:zigzagblocks}
For $n\ge1$,
$Z_n=Z_0(n)$ when $n$ is even and $Z_n=Z_1(n)$ when $n$ is odd. More generally,
\begin{equation}
 Z_\rho(x)=\mathcal G(x)
 \left(1+\sum_{\substack{d\ge3\\d\ \mathrm{odd}}}
 \chi_\rho(d)d^{-(x+1)}\right),
 \quad
 \chi_0(d)=(-1)^{(d-1)/2},\quad \chi_1(d)=1.
 \label{t1:eq:zigzagdirichlet}
\end{equation}
The Dirichlet series is absolutely convergent for $x>0$, locally uniformly in a complex neighborhood of every sufficiently large real $x$.
\end{proposition}
\begin{proof}
The poles of \eqref{t1:eq:zigzagegf} are $\rho_k=a+2\pi k=a(1+4k)$, with residue $-2$. The regularized partial fraction expansion, obtained by logarithmic differentiation of the sine product or equivalently the cotangent expansion, is
\[
 \tan(\pi/4+z/2)=1+
 2\sum_{k\in\mathbb Z}\left(\frac1{\rho_k-z}-\frac1{\rho_k}\right).
\]
The summands are $\OO(k^{-2})$ on compact sets avoiding poles. For $n\ge1$, extracting $z^n$ yields
\[
 Z_n=2n!\sum_{k\in\mathbb Z}\rho_k^{-n-1}.
\]
If $n$ is even, the positive odd denominators of residue class $1$ modulo $4$ have positive sign and those of class $3$ have negative sign, giving beta. If $n$ is odd, both signs are positive and all positive odd denominators occur, giving $(1-2^{-s})\zeta(s)$. Absolute local uniform convergence follows by comparison with $\sum d^{-\Re(s)}$.
\end{proof}

Both $Z_\rho$ are positive and strictly increasing for sufficiently large real $x$. The relative Dirichlet correction and its derivatives tend to zero exponentially, whereas
$(\log\mathcal G)'(x)=\psi(x+1)-\log a\sim\log(x/a)>0$.
Here $\psi=\Gamma'/\Gamma$.

\section{A centered Lambert carrier}
The half-shift $z=x+1/2$ eliminates the logarithmic prefactor from the gamma phase:
\begin{equation}
 \log\mathcal G(x)\sim z(\log z-1-\log a)+\log4
 +\sum_{\ell\ge1}\frac{\beta_{2\ell-1}}{z^{2\ell-1}},
 \quad
 \beta_{2\ell-1}=\frac{B_{2\ell}(1/2)}{2\ell(2\ell-1)}.
 \label{t1:eq:zigzagphase}
\end{equation}
This is a direct application of the shifted Stirling formula; odd Bernoulli polynomials at $1/2$ vanish. In particular $\beta_1=-1/24$ and $\beta_3=7/2880$.

Put
\begin{equation}
 L=\log(y/4),\qquad
 Z=\frac{L}{W_0(L/(\ee a))},\qquad
 \lambda=\log(Z/a).
 \label{t1:eq:zigzagcore}
\end{equation}
The large solution is intended. For every finite algebraic order, the exact inverse $\mathcal G^{-1}(y)$ has the form
\begin{equation}
 \mathcal G^{-1}(y)\sim Z-\frac12+
 \sum_{n\ge1}c_n(\lambda)Z^{-n}.
 \label{t1:eq:zigzagginverse}
\end{equation}
An all-order formula follows by putting $t=1/Z$ and
\begin{equation}
 \Psi_{\mathrm z}(t,u)=-\frac1\lambda\left\{
 \frac{(1+tu)\log(1+tu)-tu}{t}
 +\sum_{\ell\ge1}\beta_{2\ell-1}
 \left(\frac{t}{1+tu}\right)^{2\ell-1}\right\}.
 \label{t1:eq:zigzagpsi}
\end{equation}
Then
\begin{equation}
 c_n(\lambda)=\sum_{m=1}^n\frac1m
 \coefA{t^n u^{m-1}}\Psi_{\mathrm z}(t,u)^m.
 \label{t1:eq:zigzagcoeff}
\end{equation}
Only odd $n$ occur. Indeed if $u(t)$ solves the fixed-point equation, then $-u(-t)$ does too, and uniqueness in $t\QQ(\lambda)[[t]]$ makes $u$ odd. The first terms are
\begin{equation}
 \mathcal G^{-1}(y)=Z-\frac12+\frac1{24\lambda Z}
 -\frac{14\lambda^2+10\lambda+5}{5760\lambda^3 Z^3}
 +\OO(Z^{-5}),
 \label{t1:eq:zigzagfirst}
\end{equation}
where the error is conservative for large $\lambda=\log(Z/a)$. As with the other gamma examples, \eqref{t1:eq:zigzagginverse} is an algebraic expansion of a specified exact gamma block, not an independently chosen exponentially accurate substitute for it.

\section{Complete inverse action collisions as ordered factorizations}
For this subsection, write $g=g(L)=\mathcal G^{-1}(\ee^L)$, $D=\ddn/\ddn L$, and $g'=1/(\psi(g+1)-\log a)$. Changing $L$ by the constant $\log4$ used in \eqref{t1:eq:zigzagcore} does not change the differentiation operator.

For an odd integer $N\ge3$ and $r\ge1$, define the finite ordered-factorization coefficient
\begin{equation}
 b_{\rho,r}(N)=
 \sum_{\substack{d_1\cdots d_r=N\\d_j\ge3\ \mathrm{odd}}}
 \prod_{j=1}^r\chi_\rho(d_j).
 \label{t1:eq:orderedfactor}
\end{equation}
The sum is empty for $r>\lfloor\log_3N\rfloor$. In fact $\chi_0$ is multiplicative on odd positive integers, so $b_{0,r}(N)=\chi_0(N)b_{1,r}(N)$, but it is useful to retain the general form.

\begin{theorem}[Odd-integer inverse sectors]\label{t1:thm:zigzaginverse}
For sufficiently large $y$, the parity-resolved inverse is given by the convergent expansion
\begin{align}
 Z_\rho^{-1}(y)=g+\sum_{\substack{N\ge3\\N\ \mathrm{odd}}}N^{-(g+1)}
 \sum_{r=1}^{\lfloor\log_3N\rfloor}\frac{b_{\rho,r}(N)}{r!}
 \sum_{m=1}^r(-1)^m s(r,m)
 \left(D-(\log N)g'\right)^{m-1}g'.
 \label{t1:eq:zigzaginverse}
\end{align}
The differential operator is a variable-coefficient operator: successive applications include the derivatives of $g'$. For each $N$, the displayed coefficient is finite.
\end{theorem}
\begin{proof}
Let $q(x)=\sum_{d\ge3,\,d\text{ odd}}\chi_\rho(d)d^{-(x+1)}$. Absolute convergence gives
\[
 q(g)^r=\sum_{N\ \mathrm{odd}}b_{\rho,r}(N)N^{-(g+1)}.
\]
For any differentiable $f$,
\[
 D\bigl(N^{-(g+1)}f\bigr)
 =N^{-(g+1)}\bigl(D-(\log N)g'\bigr)f.
\]
Apply the single-marker form of \eqref{t1:eq:multiinv}, conjugate each derivative as above, and collect equal $N$. This proves the coefficient identity.

For the analytic justification, choose a fixed complex disk about large real $g$. The sums of $d^{-\Re(g+1)+h}(\log d)^j$, for fixed disk width $h$ and every fixed $j$, are finite and exponentially small in $g$. They majorize the correction and its derivatives there. The fixed-point map is therefore analytic and contractive uniformly for the marker in a disk of radius larger than one. Cauchy estimates on a smaller disk provide an absolutely summable majorant for its marker coefficients and the underlying Dirichlet sums. Thus the derivative operations and the regrouping into products $N$ are valid, and the series equals the desired inverse.
\end{proof}

The first terms illustrate a visible parity distinction:
\begin{equation}
 Z_0^{-1}(y)=g+\frac{3^{-(g+1)}}{\psi(g+1)-\log a}
 +\OO\!\left(\frac{5^{-(g+1)}}{\log g}\right),
 \label{t1:eq:zigzagevenfirst}
\end{equation}
whereas the sign of the displayed correction is negative for $Z_1^{-1}$. The sector $N=9$ receives both the original forward term $d=9$ and a nonlinear product of two $d=3$ terms. It cannot be obtained by independently inverting just the individual Dirichlet summands.

The exact discrete threshold can be reconstructed without imposing a single mixed-parity continuation:
\begin{equation}
 N_Z(y)=\min_{\rho\in\{0,1\}}
 \left\{\rho+2\left\lceil\frac{Z_\rho^{-1}(y)-\rho}{2}\right\rceil\right\}
 \label{t1:eq:zigzagthreshold}
\end{equation}
for sufficiently large $y$. This formula follows by finding the least admissible even index and the least admissible odd index separately.
% ---- hand chunk 07-zigzag-extra.tex ----

\section{A full real interpolation and the uncentred core}
\label{ct:sec:zigzag-full}

The parity sheets $Z_0,Z_1$ of \cref{t1:eq:zigzaginterp} interpolate the even
and odd subsequences separately.  In the second source's notation
$E_{\mathrm e}=Z_0$, $E_{\mathrm o}=Z_1$ and $\mathcal G(x)$ is
\cref{t1:eq:zigzaggamma}; it also fixes a single real interpolation of the
whole sequence, and inverts the gamma block without the half-shift.

% ---- T2 lines 1306--1320 ----
Two useful smooth envelopes for $x>0$ are therefore
\begin{equation}
 E_{\mathrm e}(x)=\mathcal G(x)\beta(x+1),\qquad
 E_{\mathrm o}(x)=\mathcal G(x)(1-2^{-x-1})\zeta(x+1).
 \label{t2:eq:zigzag-envelopes}
\end{equation}
They interpolate the even and odd subsequences respectively, not the full sequence. A full real interpolation is
\begin{equation}
 E(x)=\mathcal G(x)\left\{
 \sum_{\substack{d\ge1\\d\equiv1\ (4)}}d^{-x-1}
 -\cos(\pi x)\sum_{\substack{d\ge1\\d\equiv3\ (4)}}d^{-x-1}
 \right\},\qquad x>0.
 \label{t2:eq:zigzag-cont}
\end{equation}
The series converge absolutely. Since the bracket is $1+\bigO(3^{-x})$ with similarly small derivatives, each chosen function is eventually positive and increasing.
% ---- T2 lines 1322--1390 ----
\section{The common factorial core and its inverse}

Set
\[
 \bar a=\log\frac2\pi,\qquad C_A=\frac{4\sqrt{2\pi}}\pi.
\]
Stirling's expansion gives
\begin{equation}
 \log \mathcal G(x)\sim x(\log x-1+\bar a)+\frac12\log x+\log C_A
       +\sum_{m\ge1}\frac{B_{2m}}{2m(2m-1)x^{2m-1}}.
 \label{t2:eq:zigzag-A-log}
\end{equation}
For $L=\log(y/C_A)$, define
\begin{equation}
 X=\frac{L}{W_0(2L/(\ee\pi))},\qquad
 \ell=\log X,\qquad \varsigma=\ell+\bar a.
 \label{t2:eq:zigzag-core}
\end{equation}
The positive increasing branch of $A$ then has the inverse expansion
\begin{align}
 \mathcal G^{-1}(y)={}&X-\frac\ell{2\varsigma}
 +\frac1X\left(-\frac1{12\varsigma}+\frac\ell{4\varsigma^2}
                         -\frac{\ell^2}{8\varsigma^3}\right)
 +\bigO(X^{-2}).
 \label{t2:eq:zigzag-A-inverse}
\end{align}
To prove it, substitute $x=X+d_0+d_1/X$ into \eqref{t2:eq:zigzag-A-log}. The order-one equation gives $d_0=-\ell/(2s)$. At order $X^{-1}$, the terms are $\varsigma d_1+1/12+d_0/2+d_0^2/2$, giving the displayed $d_1$. Lemma~\ref{t1:thm:certificate} and the slope $\varsigma+o(1)$ justify the error. Arbitrarily high coefficients follow from the unbalanced gamma core and Proposition~\ref{t2:prop:scaled}.

\section{Prime powers organize the logarithmic sectors}

For $\Re s>1$, Euler products give
\begin{align}
 \log\beta(s)&=\sum_{\substack{p\ \mathrm{prime}\\p\ne2}}
               \sum_{m\ge1}\frac{\chi_4(p)^m}{m p^{ms}},\nonumber\\
 \log((1-2^{-s})\zeta(s))&=
             \sum_{\substack{p\ \mathrm{prime}\\p\ne2}}
               \sum_{m\ge1}\frac1{m p^{ms}},
 \label{t2:eq:zigzag-primes}
\end{align}
where $\chi_4(p)=1$ for $p\equiv1\pmod4$ and $-1$ for $p\equiv3\pmod4$. These identities can be proved directly by multiplying the absolutely convergent geometric factors over primes, using unique factorization, and expanding $-\log(1-z)=\sum_{m\ge1}z^m/m$.

Thus the primitive logarithmic sectors are indexed by odd prime powers, with actions $m\log p$. Exponentiating, differentiating, and reverting generates products indexed by odd composite integers. At any fixed action cutoff only finitely many such integers occur. This arithmetic structure is qualitatively different from the two-saddle involution case, but it is still locally finite.

Let $x_0=\mathcal G^{-1}(y)$ be retained exactly and let
\[
 p_0=(\log\mathcal G)'(x_0)=\psi(x_0+1)+\log(2/\pi).
\]
The first parity-dependent inverse layers are
\begin{align}
 E_{\mathrm e}^{-1}(y)&=x_0+\frac{3^{-x_0-1}}{p_0}
                  +\bigO\!\left(\frac{5^{-x_0}}{\log x_0}\right),\nonumber\\
 E_{\mathrm o}^{-1}(y)&=x_0-\frac{3^{-x_0-1}}{p_0}
                  +\bigO\!\left(\frac{5^{-x_0}}{\log x_0}\right).
 \label{t2:eq:zigzag-inverse-parity}
\end{align}
For the full interpolation the first correction is instead
\begin{equation}
 E^{-1}(y)=x_0+\frac{\cos(\pi x_0)3^{-x_0-1}}{p_0}
            +\bigO\!\left(\frac{5^{-x_0}}{\log x_0}\right).
 \label{t2:eq:zigzag-full-inverse}
\end{equation}
Here $5^{-x_0}$ dominates both the next integer sector and the square $9^{-x_0}$ of the first correction. These formulas follow from Theorem~\ref{ct:rmk:transport}, which also supplies every higher inverse coefficient. For the two envelopes, one may use the logarithmic sector sums \eqref{t2:eq:zigzag-primes} directly as $h$ in that theorem.

The sector scale looks unusual when expressed entirely in $y$. Since $x_0\sim\log y/\log\log y$,
\begin{equation}
 d^{-x_0}=\exp\left[-(\log d)\frac{\log y}{\log\log y}(1+o(1))\right].
 \label{t2:eq:zigzag-y-scale}
\end{equation}
For fixed $d>1$, this is smaller than every inverse power of $\log y$, but larger than every fixed positive power $y^{-\epsilon}$. It lies between familiar logarithmic and algebraic target scales. Keeping the core inverse in the action makes this hierarchy much easier to manipulate correctly.
% ---- hand chunk 07-graphs-head.tex ----

\chapter{Connected labelled graphs: every exponential layer}
\label{ct:ch:graphs}

The count of connected labelled graphs is the formal logarithm of a series
with zero radius of convergence, so no analytic singularity transfer is
available.  A giant-component argument nevertheless proves every fixed
exponential layer of $c_n/2^{\binom n2}$ with an explicit polynomial
coefficient, including a layer that vanishes identically and is regenerated
by inversion; the inverse has a quadratic core with no Lambert function and
is rigorous on the sequence's range without any interpolation.

% ---- T2 lines 1395--1558 ----
\section{A formal generating function with zero analytic radius}

Let
\[
 a_n=2^{\binom n2}
\]
be the number of simple labeled graphs on $n$ vertices, and let $c_n$ be the number of connected such graphs. These are A006125 and A001187 \cite{oeisallgraphs,oeisconnected}. The exponential formula for connected components reads
\begin{equation}
 A(z)=\sum_{n\ge0}a_n\frac{z^n}{n!},\qquad
 \log A(z)=\sum_{n\ge1}c_n\frac{z^n}{n!}.
 \label{t2:eq:graph-egf}
\end{equation}
These are formal power series. The first has radius of convergence zero, so analytic singularity transfer at a positive radius would be inapplicable. Nevertheless, fixed coefficient extraction and formal logarithms are entirely legitimate.

For completeness, a graph has a unique partition of its vertex labels into connected components. The exponential expansion of the formal series for connected graphs sums over such set partitions, with the factor $1/m!$ removing the ordering of $m$ components. This proves \eqref{t2:eq:graph-egf} without making any analytic assertion.

Define the reciprocal coefficients
\begin{equation}
 b_k=k![z^k]\frac1{A(z)}.
 \label{t2:eq:graph-b}
\end{equation}
They have the explicit finite partition formula
\begin{equation}
 b_k=k!\sum_{\substack{\sum_{j=1}^k jm_j=k\\M=\sum_{j=1}^k m_j}}
       (-1)^M M!\prod_{j=1}^k
            \frac1{m_j!}\left(\frac{2^{\binom j2}}{j!}\right)^{m_j},
 \qquad b_0=1.
 \label{t2:eq:graph-b-finite}
\end{equation}
This follows by expanding $1/A=\sum_{M\ge0}(-1)^M(A-1)^M$. The first coefficients are
\begin{equation}
 b_0,b_1,b_2,b_3,b_4,b_5,b_6
   =1,-1,0,-2,-24,-544,-22320.
 \label{t2:eq:graph-b-first}
\end{equation}

\section{A fixed-order exponential transseries theorem}

\begin{theorem}[Connected-graph expansion]
\label{t2:thm:graphs}
For every fixed integer $K\ge0$,
\begin{equation}
 \frac{c_n}{2^{\binom n2}}
 =\sum_{k=0}^{K}\binom nk b_k
       2^{k(k+1)/2}\,2^{-kn}
  +\bigO_K\!\left(n^{K+1}2^{-(K+1)n}\right).
 \label{t2:eq:graph-expansion}
\end{equation}
Thus every exponential coefficient is the polynomial
\begin{equation}
 P_k(x)=\binom xk b_k2^{k(k+1)/2}.
 \label{t2:eq:graph-P}
\end{equation}
In particular,
\begin{equation}
 \frac{c_n}{2^{\binom n2}}
 =1-2n2^{-n}-128\binom n3 2^{-3n}
       -24576\binom n4 2^{-4n}
       +\bigO(n^5 2^{-5n}).
 \label{t2:eq:graph-first}
\end{equation}
The entire $2^{-2n}$ sector vanishes.
\end{theorem}

\begin{proof}
Expand the formal logarithm in \eqref{t2:eq:graph-egf} as
\[
 \log A(z)=\sum_{m\ge1}\frac{(-1)^{m-1}}m(A(z)-1)^m.
\]
A contribution to $[z^n]$ is an ordered composition of $n$ into $m$ positive block sizes. If one block has size $n-k>n/2$, it is unique. Its $m$ possible positions cancel the factor $1/m$. The other blocks have total size $k$ and their alternating sum is exactly $[z^k](1/A(z))$. Consequently the total contribution of all terms with a giant block of size $n-k$ is
\begin{equation}
 \binom nk b_k a_{n-k},\qquad 0\le k<n/2.
 \label{t2:eq:graph-giant}
\end{equation}
Dividing by $a_n$ and using
\[
 \binom{n-k}{2}-\binom n2=-kn+\frac{k(k+1)}2
\]
gives the coefficients in \eqref{t2:eq:graph-expansion}.

It remains to bound the omitted terms. For $k\ge1$, the absolute sum defining $b_k$ is bounded by
\[
 |b_k|\le k!\,2^{k-1}a_k.
\]
Indeed, there are $2^{k-1}$ ordered compositions of $k$; for each one the labeling factor is at most $k!$, and the product of internal graph counts is at most $a_k$. Hence the absolute relative contribution of \eqref{t2:eq:graph-giant} is at most
\begin{equation}
 (2n)^k2^{-k(n-k)}.
 \label{t2:eq:graph-tail-bound}
\end{equation}
For $K+1\le k\le n/4$, the ratio of successive bounds is
$2n\,2^{-n+2k+1}$, which is less than $1/2$ for large $n$. The sum is therefore bounded by a constant times its first term,
$\bigO_K(n^{K+1}2^{-(K+1)n})$. For $n/4<k<n/2$, the exponent $k(n-k)$ is at least $3n^2/16$ up to an inessential integer-endpoint error, so these terms are $2^{-c n^2+\bigO(n\log n)}$ and are negligible at every fixed exponential order.

Finally consider terms having no block larger than $n/2$. If their sizes are $j_1,\ldots,j_m$, then
\[
 \sum_i j_i^2\le\frac n2\sum_i j_i=\frac{n^2}2.
\]
The product of their graph counts divided by $a_n$ is at most $2^{-n^2/4}$. Bounding all labeling factors by $n!$ and the number of ordered compositions by $2^{n-1}$ bounds their total by
$2^{-n^2/4+\bigO(n\log n)}$. Combining the three bounds proves \eqref{t2:eq:graph-expansion}. Substituting \eqref{t2:eq:graph-b-first} proves \eqref{t2:eq:graph-first}.
\end{proof}

This is an all-orders theorem in exponentially small lattice sectors. It is not a convergent-series assertion for a fixed real $x$. In particular, we have not defined a canonical connected-graph interpolation by summing the formal expression $\sum_kP_k(x)2^{-kx}$.

\section{A quadratic core and explicit inverse coefficients}

Let $a=\log2$. The logarithm of the all-graph count has the exact quadratic core
\begin{equation}
 H(x)=\frac a2x(x-1),\qquad
 X=\frac{1+\sqrt{1+8\log y/a}}2,
 \qquad p=H'(X)=a(X-\tfrac12).
 \label{t2:eq:graph-core}
\end{equation}
There is no Lambert function at this stage. Define $t=2^{-X}$ and seek a formal range inverse
\begin{equation}
 x\sim X+\sum_{N\ge1}d_N(X)t^N.
 \label{t2:eq:graph-inverse}
\end{equation}
All coefficients have a finite extraction formula. Put
\begin{equation}
 R_X(t,u)=\log\left(1+\sum_{k\ge1}
                 P_k(X+u)t^k\ee^{-kau}\right).
 \label{t2:eq:graph-R}
\end{equation}
Then Proposition~\ref{t2:prop:scaled}, in its additive form, gives
\begin{equation}
 d_N(X)=\sum_{m=1}^{N}\frac{(-1)^m}{m}
 [t^Nu^{m-1}]\left(\frac{R_X(t,u)}{p+au/2}\right)^m.
 \label{t2:eq:graph-d}
\end{equation}
The denominator is simply $(H(X+u)-H(X))/u$. At exponential order $N$, only $P_1,\ldots,P_N$ are needed, each given explicitly by \eqref{t2:eq:graph-b-finite} and \eqref{t2:eq:graph-P}.

The first two coefficients are
\begin{align}
 d_1(X)&=\frac{2X}{p},\nonumber\\
 d_2(X)&=\frac{2X^2-2(aX-1)d_1(X)-(a/2)d_1(X)^2}{p}.
 \label{t2:eq:graph-d-first}
\end{align}
To see the second formula directly, use
\[
 R_X(t,u)=-2Xt+2(aX-1)ut-2X^2t^2+\cdots
\]
in $pu+(a/2)u^2+R_X(t,u)=0$. Although $P_2=0$, the quadratic logarithm and the shift of the first sector generate a nonzero $d_2$. This is a concrete example of inverse sector creation.

\begin{corollary}[A rigorous inverse on the graph sequence's range]
\label{t2:cor:graph-range}
At $y=c_n$, the truncation through order $N$ of \eqref{t2:eq:graph-inverse} satisfies
\begin{equation}
 X+\sum_{j=1}^Nd_j(X)2^{-jX}
 =n+\bigO_N\!\left(n^N2^{-(N+1)n}\right).
 \label{t2:eq:graph-range-error}
\end{equation}
No interpolation of $c_n$ is required for this statement.
\end{corollary}
\begin{proof}
For fixed $N$, use the smooth finite forward model
\[
 F_N(x)=\ee^{H(x)}\left(1+\sum_{k=1}^NP_k(x)2^{-kx}\right).
\]
It is positive and increasing for all sufficiently large $x$, with logarithmic derivative $ax+\bigO(1)$. Theorem~\ref{t2:thm:graphs} gives a logarithmic discrepancy between $F_N(n)$ and $c_n$ of order $n^{N+1}2^{-(N+1)n}$. Lemma~\ref{t1:thm:certificate} therefore places the exact inverse $F_N^{-1}(c_n)$ within $\bigO_N(n^N2^{-(N+1)n})$ of $n$.

The formal coefficient cancellation in \eqref{t2:eq:graph-d} gives the same error for its finite inverse truncation. One can bound the omitted terms directly: $P_k(x)=\bigO_k(x^k)$, $p\asymp X$, and \eqref{t2:eq:graph-d} gives $d_j(X)=\bigO_j(X^{j-1})$. Taylor remainders for this finite model then have size $\bigO_N(X^{N+1}2^{-(N+1)X})$ in logarithmic phase, or $\bigO_N(X^N2^{-(N+1)X})$ after division by the slope. Since $X=n+\bigO(2^{-n})$, the same bound holds with $n$. Combining the estimates proves the result.
\end{proof}

Thus the inverse expansion is operationally meaningful even without a continuous connected-graph counting function. At sufficiently large exact counts, nearest-integer rounding recovers the vertex number. A threshold for an arbitrary target requires either bracketing consecutive exact counts or an interpolation together with the separation test from Section~\ref{ct:ch:objects}.
% ---- hand chunk 07-necklace-head.tex ----

\chapter{Necklaces, Lyndon words, and irreducible polynomials}
\label{ct:ch:necklaces}

Divisor sums have exponential corrections whose \emph{coefficients} are
arithmetic indicators and whose \emph{actions} accumulate at a finite
value.  Neither feature fits a finite exponential grid.  This chapter first
shows what fails and what survives (finite-cutoff expansions and
finite-model range inverses with a trigonometric divisibility filter), and
then constructs, on the subsequences of fixed radical, exact analytic sheets
whose inverse transseries converge and whose leading inverse is $W_{-1}$;
a uniform two-candidate rule resolves the integer threshold without
choosing any interpolation at all.

% ---- T2 lines 1563--1675 ----
\section{Exact divisor sums}

Fix an alphabet size $q\ge2$. The number of length-$n$ necklaces under cyclic rotation is
\begin{equation}
 N_q(n)=\frac1n\sum_{d\mid n}\varphi(d)q^{n/d},
 \label{t2:eq:necklace}
\end{equation}
and the number of primitive necklaces, equivalently Lyndon words of length $n$ over the ordered alphabet, is
\begin{equation}
 L_q(n)=\frac1n\sum_{d\mid n}\mu(d)q^{n/d}.
 \label{t2:eq:lyndon}
\end{equation}
The binary cases are A000031 and A001037 \cite{oeisnecklace,oeislyndon}.

For \eqref{t2:eq:necklace}, Burnside averaging over the $n$ rotations counts $q^{\gcd(n,r)}$ words fixed by rotation through $r$ places. Grouping by $d=n/\gcd(n,r)$ gives $\varphi(d)$ rotations of the corresponding type. For \eqref{t2:eq:lyndon}, let $P_q(n)$ be the number of aperiodic words. Every word has a unique least period dividing $n$, so $q^n=\sum_{d\mid n}P_q(d)$. M\"obius inversion gives $P_q(n)=\sum_{d\mid n}\mu(d)q^{n/d}$, and every aperiodic orbit has exactly $n$ words. This proves both formulas.

Factoring out $q^n/n$, write
\begin{equation}
 N_q(n)=\frac{q^n}{n}\left(1+
 \sum_{\substack{d\mid n\\d\ge2}}\varphi(d)\ee^{-\omega_dn}\right),
 \qquad \omega_d=(\log q)\left(1-\frac1d\right).
 \label{t2:eq:necklace-sectors}
\end{equation}
For Lyndon words, replace $\varphi(d)$ by $\mu(d)$. Each divisor contributes a different exponential action.

\begin{proposition}[Uniform finite action cutoffs]
\label{t2:prop:necklace-cutoff}
For every fixed integer $D\ge2$,
\begin{align}
 \frac{nN_q(n)}{q^n}
 &=1+\sum_{d=2}^{D}\varphi(d)\,\mathbf1_{d\mid n}
             q^{-n(1-1/d)}
    +\bigO\!\left(nq^{-n(1-1/(D+1))}\right),\nonumber\\
 \frac{nL_q(n)}{q^n}
 &=1+\sum_{d=2}^{D}\mu(d)\,\mathbf1_{d\mid n}
             q^{-n(1-1/d)}
    +\bigO\!\left(nq^{-n(1-1/(D+1))}\right).
 \label{t2:eq:necklace-cutoff}
\end{align}
\end{proposition}
\begin{proof}
Every omitted divisor satisfies $d\ge D+1$, hence
$1-1/d\ge1-1/(D+1)$. In the necklace sum use
$\sum_{d\mid n}\varphi(d)=n$. In the primitive sum use
$\sum_{d\mid n}|\mu(d)|\le\sum_{d\mid n}1\le n$.
The stated bounds follow immediately.
\end{proof}

\section{Why constant-coefficient exponentials are insufficient}

Let $\epsilon_n=nN_q(n)/q^n-1$. The first candidate action is $(\log q)/2$. But Proposition~\ref{t2:prop:necklace-cutoff} gives
\begin{equation}
 q^{n/2}\epsilon_n\longrightarrow
 \begin{cases}1,&n\to\infty\ \text{through even integers},\\
               0,&n\to\infty\ \text{through odd integers}.
 \end{cases}
 \label{t2:eq:necklace-parity-obstruction}
\end{equation}
For primitive necklaces the even limit is $-1$. Thus no single constant coefficient can describe the first exponential correction on all integers. Arithmetic indicators are essential, not negligible nuisances.

Moreover,
\begin{equation}
 \omega_2<\omega_3<\cdots<\log q,
 \qquad \omega_d\longrightarrow\log q.
 \label{t2:eq:necklace-action-accumulation}
\end{equation}
There are infinitely many primitive actions below the finite action $\log q$. In particular, the nonlinear square of the first sector has action $2\omega_2=\log q$, but infinitely many elementary divisor sectors precede it. A finite grid of positive generators cannot have this property: its support below every fixed bound is finite.

This is an obstruction to the \emph{locally finite finite-grid calculus}, not a proof that no generalized transseries can exist. A well-ordered support can have order type $\omega$ and accumulate at a cut; larger Hahn-type or arithmetic-coefficient constructions may accommodate it. What fails is the convenient claim that finitely many sector coefficients suffice to reach every bounded action.

\section{Explicit finite-cutoff interpolations and inverses}

The dominant counting model is $f(x)=q^x/x$. Put $a=\log q$. Its large inverse is
\begin{equation}
 X=-\frac1a W_{-1}\!\left(-\frac ay\right),
 \qquad aX-\log X=\log y.
 \label{t2:eq:necklace-core}
\end{equation}
This core is exact for the dominant model and again requires the negative Lambert branch.

To give a precise meaning to finite arithmetic corrections off the integers, define the real trigonometric filter
\begin{equation}
 D_d(x)=\frac1d\sum_{r=0}^{d-1}\cos\left(\frac{2\pi r x}{d}\right).
 \label{t2:eq:divisibility-filter}
\end{equation}
At integer $x=n$, the finite geometric sum proves $D_d(n)=\mathbf1_{d\mid n}$. For fixed $D$, set
\begin{equation}
 F_{q,D}^{(w)}(x)=\frac{q^x}{x}
       \left(1+\sum_{d=2}^{D}w_dD_d(x)\ee^{-\omega_dx}\right),
 \label{t2:eq:necklace-finite-model}
\end{equation}
where $w_d=\varphi(d)$ or $w_d=\mu(d)$. This is an explicit, eventually positive and increasing function: the finite correction and its derivative tend exponentially to zero, while $(\log f)'=a-1/x$ tends to $a>0$.

Every inverse sector of this \emph{fixed finite model} follows from Theorem~\ref{ct:rmk:transport}, with $g$ the exact Lambert inverse of the dominant phase. Its first displacement is
\begin{equation}
 (F_{q,D}^{(w)})^{-1}(y)
 =X-\frac{w_2D_2(X)}{a-1/X}q^{-X/2}
          +\bigO(q^{-2X/3}),
 \label{t2:eq:necklace-finite-inverse}
\end{equation}
where $D_2(X)=(1+\cos\pi X)/2$. For $D=2$, the actual next nonlinear term is smaller, of order $q^{-X}$; the displayed weaker bound remains valid. For $D\ge3$, a divisor-$3$ sector may occur at order $q^{-2X/3}$.

At the true sequence values, Proposition~\ref{t2:prop:necklace-cutoff} and Lemma~\ref{t1:thm:certificate} imply the rigorous range approximations
\begin{align}
 (F_{q,D}^{(\varphi)})^{-1}(N_q(n))
 &=n+\bigO\!\left(nq^{-n(1-1/(D+1))}\right),\nonumber\\
 (F_{q,D}^{(\mu)})^{-1}(L_q(n))
 &=n+\bigO\!\left(nq^{-n(1-1/(D+1))}\right).
 \label{t2:eq:necklace-range}
\end{align}
Indeed, the relative forward error has exactly that size, and the logarithmic derivative of the finite model stays bounded away from zero. These estimates justify finite-cutoff inverse computations without inventing an infinite real interpolation.

It would be unjustified to replace $D$ by infinity in \eqref{t2:eq:necklace-finite-model} without a convergence analysis: for real noninteger $x$, the filters no longer restrict the sum to divisors of a single integer. The finite divisor identity has become an unrelated infinite series. This is a concrete place where a formal manipulation of a combinatorial formula can silently change the problem.
% ---- hand chunk 07-necklace-t3.tex ----

\section{Irreducible polynomials and the normalized arithmetic sectors}
\label{ct:sec:irreducibles}

For a prime power $q$ the primitive-necklace count $L_q(n)$ of
\cref{t2:eq:lyndon} is also the number $I_q(n)$ of monic irreducible
polynomials of degree $n$ over $\F_q$: the polynomial $X^{q^n}-X$ is the
product of all monic irreducibles whose degrees divide $n$, with no repeated
roots, so $q^n=\sum_{d\mid n}dI_q(d)$ and M\"obius inversion gives the
formula.  Throughout this section $h=\log q$ and $I_q(n)=L_q(n)$.

% ---- T3 lines 1264--1315 ----
Normalize by the leading term:
\begin{equation}\label{t3:eq:arithmetic-sectors}
 I_q(n)=\frac{q^n}{n}
 \left(1+\sum_{d\ge2}\mu(d)\one_{d\mid n}
       \ee^{-h(1-1/d)n}\right).
\end{equation}
At each integer $n$ the sum is finite. The exponential actions are
$h(1-1/d)$, with coefficients modulated by divisibility. They accumulate at $h$ as $d\to\infty$.

For a fixed cutoff $D_0\ge2$, the omitted relative tail satisfies
\begin{equation}\label{t3:eq:arithmetic-cutoff}
 \left|\sum_{\substack{d\mid n\\d>D_0}}
       \mu(d)\ee^{-h(1-1/d)n}\right|
 \le n\ee^{-h(1-1/(D_0+1))n}.
\end{equation}
Indeed, there are at most $n$ possible divisors and each exponent is at least the displayed action. This gives a rigorous finite-sector asymptotic decomposition uniformly over all integers $n$.

It is not, however, a single analytic power series in $\ee^{-\lambda n}$ with fixed constant coefficients. Along even $n$ the first relative correction is
\[
 -q^{-n/2}+O(nq^{-2n/3}),
\]
whereas along prime $n$ the exact correction is $-q^{1-n}$. The lowest nonzero action changes with the arithmetic subsequence. On a fixed congruence class modulo $\operatorname{lcm}(1,\ldots,D_0)$, the divisibility coefficients up to $D_0$ are fixed, but this does not fix the later coefficients.

There is a distinction here between an arithmetic sector decomposition and a particular formal transseries field. The accumulation of actions does not by itself exclude every possible Hahn-series construction. The concrete point is that no fixed finite exponential grid and no fixed analytic coefficient germ describes all the divisibility data in \eqref{t3:eq:arithmetic-sectors}.

\section{A finite-radical interpolation removes the moving divisors}

There is nevertheless a natural class of subsequences on which a full analytic inverse transseries can be constructed. Fix a squarefree integer $R\ge2$. Whenever $\rad(n)=R$, all nonzero M\"obius terms in \eqref{t2:eq:lyndon} correspond to divisors of $R$. Define
\begin{equation}\label{t3:eq:ER}
 E_R(t)=\sum_{d\mid R}\mu(d)t^{R-R/d}
       =1+\sum_{\substack{d\mid R\\d>1}}\mu(d)t^{R-R/d},
 \qquad \kappa=\frac hR.
\end{equation}
This is a finite polynomial with constant term one. The analytic function
\begin{equation}\label{t3:eq:fixedradical-function}
 J_{q,R}(x)=\frac{\ee^{hx}}x E_R(\ee^{-\kappa x})
          =\frac1x\sum_{d\mid R}\mu(d)\ee^{hx/d}
\end{equation}
agrees with $I_q(n)$ at every $n$ with radical $R$. It is positive and increasing for sufficiently large $x$.

The function is now explicitly specified, so its inverse has a well-defined exponentially small structure. We use
\begin{equation}\label{t3:eq:logphase}
 P(x)=hx-\log x,\qquad B(t)=\log E_R(t),\qquad
 D=P'(s)=h-\frac1s.
\end{equation}
The large real center is
\begin{equation}\label{t3:eq:Wcenter}
 \boxed{\displaystyle
 s=-\frac1h W_{-1}\!\left(-\frac hY\right).}
\end{equation}
The $W_{-1}$ branch is required because $s\to\infty$. The principal branch would give the small solution of $\ee^{hx}/x=Y$. The branch convention agrees with DLMF and \texttt{ProductLog[-1,z]} \cite{dlmfW,wlproductlog}.

% ---- T3 lines 1363--1483 ----
\begin{theorem}[Inverse transseries on a fixed-radical sheet]\label{t3:thm:radical-inverse}
For fixed $q>1$ and squarefree $R\ge2$, the large real inverse of \eqref{t3:eq:fixedradical-function} has a convergent expansion
\begin{equation}\label{t3:eq:radical-inverse}
 J_{q,R}^{-1}(Y)=s+\sum_{m\ge1}d_m(s)\ee^{-mhs/R},
\end{equation}
where $s$ is \eqref{t3:eq:Wcenter} and the coefficients are given by
\eqref{t3:eq:general-inverse-extraction} or \eqref{t3:eq:general-derivative}. Each coefficient is a rational function of $s,h,\kappa$ and the finitely many coefficients of $B$ needed at that weight.
\end{theorem}
\begin{proof}
For $s\ge2/h$, $P'(s)\ge h/2$. On a fixed disk $|\delta|\le\rho$ with $s$ large,
\[
 P(s+\delta)-P(s)-P'(s)\delta=O(\delta^2/s^2),
\]
uniformly in $s$. This follows by expanding $-\log(1+\delta/s)$, or from its second derivative. Since $B$ is analytic near zero, the same Rouch\'e argument used in Theorem~\ref{t3:thm:inverse} gives a holomorphic solution in a fixed disk of the auxiliary small parameter $t$, uniformly for sufficiently large $s$. Substituting $t=\ee^{-\kappa s}$ proves convergence. Formula \eqref{t3:eq:general-derivative}, together with $P'(s)=h-1/s$, proves rationality of the coefficients. Eventual monotonicity follows from $P'(x)\to h>0$ and the exponentially small tail derivative.
\end{proof}

\section{Prime-power indices: an explicit first example}

If $n=p^a$ with a fixed prime $p$, then
\[
 I_q(n)=\frac{q^n-q^{n/p}}n.
\]
It is convenient to use the primitive action
$\kappa=h(1-1/p)$ rather than $h/p$. The selected interpolation is
\begin{equation}\label{t3:eq:primepower-interp}
 J_{q,p}(x)=\frac{\ee^{hx}}x(1-\ee^{-\kappa x}),
 \qquad B(t)=\log(1-t).
\end{equation}
With $D=h-1/s$ and $t_0=\ee^{-\kappa s}$, the first two inverse sectors are
\begin{equation}\label{t3:eq:primepower-inverse}
 \boxed{\displaystyle
 J_{q,p}^{-1}(Y)=s+\frac{t_0}{D}
 +\left(\frac1{2D}-\frac{\kappa}{D^2}
              -\frac1{2s^2D^3}\right)t_0^2+O(t_0^3).}
\end{equation}
To check the second coefficient directly, expand
\[
 P(s+\delta)-P(s)=D\delta+\frac{\delta^2}{2s^2}+O(\delta^3),
 \qquad B(t)=-t-\frac{t^2}{2}+O(t^3),
\]
and substitute into the inverse equation. The coefficient of $t_0$ is $1/D$; the coefficient of $t_0^2$ is exactly the bracket in \eqref{t3:eq:primepower-inverse}.

The exponential parameter itself can be written in terms of $Y$ and the exact Lambert center:
\begin{equation}\label{t3:eq:primepower-scale}
 t_0=(Ys)^{-(1-1/p)},\qquad \ee^{hs}=Ys.
\end{equation}
Thus the inverse combines a logarithmic Lambert-$W$ scale with fractional inverse powers of $Y$ and corresponding powers of $s$. Writing
$\ell=\log(Y/h)$ and $v=\log\ell$, the beginning of the ordinary logarithmic expansion is
\begin{equation}\label{t3:eq:W-logexpansion}
 hs=\ell+v+\frac v\ell+\frac{v-v^2/2}{\ell^2}
       +O(v^3/\ell^3).
\end{equation}
It follows by substitution into $hs-\log(hs)=\ell$; the full logarithmic calculus is already developed in the source article \cite{tai}. Retaining $s$ exactly in \eqref{t3:eq:primepower-inverse} keeps all of those logarithmic corrections organized inside a single center.

For a second example, $R=6$ gives
\begin{equation}\label{t3:eq:R6}
 E_6(t)=1-t^3-t^4+t^5,
 \qquad
 J_{q,6}(x)=\frac{q^x-q^{x/2}-q^{x/3}+q^{x/6}}x.
\end{equation}
The inverse weights in $t_0=\ee^{-hs/6}$ start at $3,4,5$ and include their additive combinations. Weights $1$ and $2$ vanish. The first three nonzero terms are
\begin{equation}\label{t3:eq:R6-first}
 J_{q,6}^{-1}(Y)=s+\frac{t_0^3+t_0^4-t_0^5}{D}+O(t_0^6).
\end{equation}
Nonlinear interactions first enter at weight six.

\section{A uniform two-candidate theorem for the integer inverse}

The fixed-radical sheets do not define a single interpolation for all $n$. The integer threshold can nevertheless be controlled uniformly without choosing any interpolation of the arithmetic corrections.

\begin{lemma}[Uniform deficit bound]\label{t3:lem:irreducible-bound}
For every integer $q\ge2$ and $n\ge1$,
\begin{equation}\label{t3:eq:uniform-deficit}
 0\le q^n-nI_q(n)
 \le\sum_{j=1}^{\lfloor n/2\rfloor}q^j
 =\frac{q^{\lfloor n/2\rfloor+1}-q}{q-1}.
\end{equation}
In particular,
\begin{equation}\label{t3:eq:I-lower}
 \frac{q^n}{n}\left(1-\frac q{q-1}q^{-n/2}\right)
 \le I_q(n)\le\frac{q^n}{n}.
\end{equation}
\end{lemma}
\begin{proof}
The deficit counts nonprimitive words. Every nonprimitive word has a proper period dividing $n$, hence a period at most $\lfloor n/2\rfloor$. There are at most $q^j$ words obtained by repeating a block of length $j$. Taking the union bound over all $1\le j\le\lfloor n/2\rfloor$ proves \eqref{t3:eq:uniform-deficit}; allowing lengths that do not divide $n$ merely enlarges the upper bound. Dividing by $q^n$ gives \eqref{t3:eq:I-lower}.
\end{proof}

\begin{theorem}[An exact two-candidate threshold rule]\label{t3:thm:two-candidate}
Let $q\ge2$ be an integer, and let
\[
 X(Y)=-\frac1hW_{-1}(-h/Y),\qquad n_0=\lceil X(Y)\rceil.
\]
If $n_0\ge4$, then the integer threshold
$N_q(Y)=\min\{n\ge1:I_q(n)\ge Y\}$ satisfies
\begin{equation}\label{t3:eq:two-candidate}
 \boxed{\displaystyle
 N_q(Y)=
 \begin{cases}
 n_0,&I_q(n_0)\ge Y,\\
 n_0+1,&I_q(n_0)<Y.
 \end{cases}}
\end{equation}
Thus one exact divisor-sum comparison resolves the threshold after computing the Lambert center.
\end{theorem}
\begin{proof}
Put $B_q(n)=q^n/n$. Its values are nondecreasing for positive integers, and strictly increasing from $n\ge2$, because
$B_q(n+1)/B_q(n)=qn/(n+1)$. The choice of $n_0$ and the large increasing branch of $q^x/x$ imply
$I_q(n)\le B_q(n)<Y$ for every $n<n_0$.

It remains to prove $I_q(n_0+1)\ge Y$. By \eqref{t3:eq:I-lower}, for $n\ge4$,
\begin{align*}
 \frac{I_q(n+1)}{B_q(n)}
 &\ge\frac{qn}{n+1}
   \left(1-\frac q{q-1}q^{-(n+1)/2}\right)\\
 &\ge\frac{2n}{n+1}\left(1-2\cdot2^{-(n+1)/2}\right)>1.
\end{align*}
The last expression increases with $n$ and already exceeds one at $n=4$.
Since $Y\le B_q(n_0)$, this proves $I_q(n_0+1)>Y$. No separate monotonicity theorem for $I_q(n)$ was required.
\end{proof}

The extra candidate cannot be discarded. For $Y=q^n/n$ with $n\ge4$, nonprimitive words exist, so $I_q(n)<Y$ and the answer is $n+1$. On the other hand, at $Y=I_q(n)$ with $n$ sufficiently large the answer is $n$. The switch occurs inside a relative strip of width $O(q^{-n/2})$. This is precisely where uncontrolled rounding of a smooth approximation loses information.
% ---- hand chunk 07-necklace-all.tex ----

\section{All necklaces on the same scale}
\label{ct:sec:all-necklaces}

For the count $N_q(n)=M_q(n)$ of all necklaces, \cref{t2:eq:necklace}, the
same considerations apply with the totient in place of the M\"obius function.

% ---- T3 lines 1494--1499 ----
Its relative arithmetic sectors have the same actions as \eqref{t3:eq:arithmetic-sectors}, now with coefficients $\varphi(d)\one_{d\mid n}$, all positive. A uniform elementary estimate is
\[
 0\le\frac{nM_q(n)}{q^n}-1
 \le n^2q^{-n/2}.
\]
There are at most $n$ divisor terms, each $\varphi(d)\le n$, and $n/d\le n/2$ for $d>1$. Thus the Lambert center still gives the leading inverse scale, but its correction is on the opposite side from that for primitive necklaces. Unlike the M\"obius sum, fixing the radical alone does not fix all nonzero totient terms; powers of its primes continue to appear among the divisors. This explains why the finite polynomial $E_R$ construction is specific to the primitive-word problem.
% ---- hand chunk 08-part7-head.tex ----

\part{$q$-products and finite-field enumeration}
\label{ct:part:qproducts}

The families of this part are exact finite $q$-products: orders of the
classical groups over $\F_q$, flag counts, Gaussian binomial and multinomial
coefficients, $q$-Catalan numbers, and Galois numbers.  Three features make
them especially tractable.  Their normalized tails are analytic in $q^{-x}$,
so their transseries genuinely converge.  Inversion produces coefficients
with an explicit finite combinatorial structure, by the convergent theorem
of \cref{ct:sec:quadratic}.  And two different kinds of discrete modulation
occur, a finite number of residue-class theta constants and an unbounded
collection of divisibility indicators, which must not be conflated.

Throughout the fixed-$q$ chapters,

% ---- T3 lines 90--115 ----
Throughout the fixed-$q$ sections,
\begin{equation}\label{t3:eq:qnotation}
 q>1,\qquad Q=q^{-1}\in(0,1),\qquad h=\log q=-\log Q.
\end{equation}
A finite-field interpretation requires $q$ to be a prime power. The analytic formulas themselves are meaningful for every real $q>1$. A separate subsection treats $q$-factorials with base in $(0,1)$.

\begin{longtable}{@{}>{\raggedright\arraybackslash}p{0.25\textwidth}>{\raggedright\arraybackslash}p{0.32\textwidth}>{\raggedright\arraybackslash}p{0.34\textwidth}@{}}
\caption{The principal examples and their inverse scales. Constants and residue labels are retained in the exact centers used later.}\label{t3:tab:map}\\
\toprule
Family & Leading size & Structure of the inverse correction\\
\midrule\endfirsthead
\toprule Family & Leading size & Structure of the inverse correction\\\midrule\endhead
$\lvert\GL(n,q)\rvert$ & $\Pinf q^{n^2}$ & Powers of $q^{-s}$, with $s\asymp\sqrt{\log Y}$\\
Complete flags, $[n]_q!$ & $\Pinf q^{n(n+1)/2}(q-1)^{-n}$ & Powers of $q^{-s}$ around an exact quadratic center\\
Fixed-rank $\qbinq{n}{k}$ & $q^{k(n-k)}/\pp{k}$ & A convergent Puiseux series in $Y^{-1/k}$\\
Scaled Gaussian coefficients and partial flags & $Cq^{a_0n^2}$ & An additive semigroup of exponential weights\\
Binomial-version $q$-Fuss--Catalan & $(1-Q)\Pinf^{-1}q^{r n(n-1)}$ & Powers of $q^{-s}$; explicit cancellation rules\\
$\lvert\Sp(2n,q)\rvert$ & $\mathcal P_{Q^2}q^{2n^2+n}$ & Only even weights in $q^{-s}$\\
Unitary isometry groups & $\mathcal P_{-Q}q^{n^2}$ & Two parity sheets with different first corrections\\
Generalized Galois numbers, $r$ blocks & $\T^{(r)}_{n\bmod r}\Pinf^{1-r}q^{(r-1)n^2/(2r)}$ & Root-lattice theta coefficients; powers of $q^{-s/r}$\\
Normalized product limits & A nonzero finite limit & A logarithm plus a convergent power or Puiseux series in the deficit\\
Primitive necklaces, irreducible polynomials & $q^n/n$ & $W_{-1}$ center; arithmetic sectors or fixed-radical analytic sheets\\
\bottomrule
\end{longtable}

The finite product definitions and $q$-binomial identities used below are standard; DLMF and the Wolfram documentation specify their conventions \cite{dlmfq,wlqbinomial,wlqpochhammer}. OEIS is used to identify representative integer specializations, not as a substitute for proof. The leading theta asymptotics of generalized Galois numbers have an established literature \cite{kousidis,bliem}; our development gives a direct coefficient-level construction and then carries it through inversion.
% ---- hand chunk 08-tail-head.tex ----

\chapter{The analytic tail of a finite $q$-product}
\label{ct:ch:qtail}

% ---- T3 lines 180--259 ----
\label{t3:sec:qtail}

\section{The common building block}

Define
\begin{equation}\label{t3:eq:pinf}
 \Pinf=(Q;Q)_\infty=\prod_{j=1}^{\infty}(1-Q^j)>0,
 \qquad
 \A(t)=-\log(Qt;Q)_\infty.
\end{equation}
The logarithm is the analytic branch equal to zero at $t=0$. The first possible singularity is at $t=Q^{-1}$.

\begin{lemma}[Tail coefficients and an explicit remainder]\label{t3:lem:Atail}
For $|t|<Q^{-1}$,
\begin{equation}\label{t3:eq:Aseries}
 \A(t)=\sum_{m\ge1}\alpha_m t^m,
 \qquad \alpha_m=\frac{Q^m}{m(1-Q^m)}.
\end{equation}
For $M\ge0$ and $Q|t|<1$,
\begin{equation}\label{t3:eq:Atailbound}
 \left|\A(t)-\sum_{m=1}^M\alpha_m t^m\right|
 \le \frac{(Q|t|)^{M+1}}
 {(M+1)(1-Q)(1-Q|t|)}.
\end{equation}
\end{lemma}
\begin{proof}
Expand $-\log(1-Q^jt)=\sum_{m\ge1}Q^{jm}t^m/m$ and sum over $j\ge1$. Absolute convergence on compact subsets of $|t|<Q^{-1}$ permits interchanging the sums and gives \eqref{t3:eq:Aseries}. For the remainder use $1-Q^m\ge1-Q$, $m\ge M+1$, and then sum a geometric series.
\end{proof}

The interpolation of a finite product that will be used throughout is
\begin{equation}\label{t3:eq:product-interpolation}
 \mathcal P_Q(x)=\frac{(Q;Q)_\infty}{(Q^{x+1};Q)_\infty}
 =\Pinf\exp\A(Q^x),\qquad x>-1.
\end{equation}
At $x=n\in\N$ it is precisely $\pp{n}$. It decreases to $\Pinf$ as $x\to\infty$, since $Q^x$ decreases and $\A$ has positive coefficients. The notation $\mathcal P_Q(x)$ is distinguished from the limiting constant $\Pinf$.

\section{Forward coefficients without logarithms}

Two Euler expansions, in the convention of \cite{dlmfq}, are
\begin{align}
 (Qt;Q)_\infty
 &=\sum_{j\ge0}\frac{(-1)^jQ^{j(j+1)/2}}{\pp{j}}t^j,
 \label{t3:eq:euler-entire}\\
 \frac1{(Qt;Q)_\infty}
 &=\sum_{j\ge0}\frac{Q^j}{\pp{j}}t^j,
 \qquad |Qt|<1.
 \label{t3:eq:euler-recip}
\end{align}
For completeness, write $E(z)=(z;Q)_\infty$. Its functional equation
$E(z)=(1-z)E(Qz)$ gives
$c_j=-Q^{j-1}c_{j-1}/(1-Q^j)$ for its coefficients. This proves the first expansion, including entire convergence. The reciprocal equation
$(1-z)/E(z)=1/E(Qz)$ gives $c_j=c_{j-1}/(1-Q^j)$ for $1/E(z)$ and proves the second. In particular,
\begin{equation}\label{t3:eq:P-forward}
 \mathcal P_Q(x)=\Pinf\sum_{m\ge0}\frac{Q^{m(x+1)}}{\pp{m}}.
\end{equation}

More generally, if $B(t)=\sum b_m t^m$, then
\begin{equation}\label{t3:eq:forward-Bell}
 \coefB{t}{m}\ee^{B(t)}
 =\sum_{\substack{k_1,\ldots,k_m\ge0\\\sum j k_j=m}}
       \prod_{j=1}^m\frac{b_j^{k_j}}{k_j!}.
\end{equation}
This is a finite formula at every weight. It follows by multiplying
$\prod_j\sum_{k\ge0}b_j^k t^{jk}/k!$; only $j\le m$ can contribute to weight $m$.

\section{A finite-product coefficient rule}

Many families can be reduced immediately to this format. Suppose
\begin{equation}\label{t3:eq:product-rule}
 \ee^{B(t)}=\prod_{\nu=1}^J
       (Q^{\beta_\nu}t^{r_\nu};Q^{e_\nu})_\infty^{\eta_\nu},
 \qquad r_\nu\in\N_{>0},\quad e_\nu>0,
\end{equation}
where the powers use the branch equal to one at $t=0$. Real $\beta_\nu$ and real or complex $\eta_\nu$ are permitted. Taking the logarithm gives the explicit formula
\begin{equation}\label{t3:eq:product-log-coeff}
 b_m=-\sum_{\substack{1\le\nu\le J\\r_\nu\mid m}}
 \eta_\nu\frac{Q^{\beta_\nu m/r_\nu}}
 {(m/r_\nu)(1-Q^{e_\nu m/r_\nu})}.
\end{equation}
The identity is valid in a neighborhood of zero where all the factors are nonzero. Combining \eqref{t3:eq:product-log-coeff}, \eqref{t3:eq:forward-Bell}, and the inversion theorem below gives an explicit procedure for the forward and inverse coefficients of any such finite product.
% ---- hand chunk 08-GL-head.tex ----

\chapter{General linear groups, complete flags, and $q$-factorials}
\label{ct:ch:GL}

% ---- T3 lines 427--540 ----
\label{t3:sec:GL}

\section{Ordered bases and general linear groups}

An invertible $n\times n$ matrix over $\F_q$ is an ordered basis. Its first column can be any nonzero vector, its second can be any vector outside the first column's span, and so on. Therefore
\begin{equation}\label{t3:eq:GLcount}
 g_n(q)=|\GL(n,q)|=\prod_{j=0}^{n-1}(q^n-q^j)
       =q^{n^2}\pp{n}.
\end{equation}
The specialization $q=2$ is OEIS A002884 \cite{oeisGL}. Our chosen interpolation is
\begin{equation}\label{t3:eq:GLinterp}
 g_q(x)=\Pinf\exp\bigl(hx^2+\A(Q^x)\bigr).
\end{equation}
It agrees with \eqref{t3:eq:GLcount} at every nonnegative integer. Its exact forward transseries is especially simple:
\begin{equation}\label{t3:eq:GLforward}
 g_q(x)=\Pinf q^{x^2}
       \sum_{m\ge0}\frac{Q^{m(x+1)}}{\pp{m}}.
\end{equation}
For $q\ge2$, the interpolation is already increasing for $x\ge1$. Indeed,
\[
 \frac{(\log g_q)'(x)}h
 =2x-\sum_{m\ge1}\frac{Q^{m(x+1)}}{1-Q^m},
\]
and the sum is bounded by
$Q^{x+1}/((1-Q)(1-Q^{x+1}))\le2/3$ in this range.
For arbitrary fixed $q>1$, eventual monotonicity is sufficient and follows from Theorem~\ref{t3:thm:inverse}.

Put
\begin{equation}\label{t3:eq:GLcenter}
 s=\sqrt{\frac{\log(Y/\Pinf)}h},\qquad D=2hs,\qquad t_0=Q^s.
\end{equation}
Using $a=\lambda=h$ and $b_m=\alpha_m$ gives
\begin{align}\label{t3:eq:GLinverse}
 g_q^{-1}(Y)=s
 &-\frac{\alpha_1}{D}t_0\\[-0.4em]
 &-\left(\frac{\alpha_2}{D}
       +\frac{h\alpha_1^2}{D^2}
       +\frac{h\alpha_1^2}{D^3}\right)t_0^2
 +O(t_0^3/D).\notag
\end{align}
Every later coefficient is given by \eqref{t3:eq:finite-dm}. All corrections are nonpositive. This agrees with the elementary inequality
$g_q(x)>\Pinf q^{x^2}$, which puts the true inverse below the leading center.

For $q=2$,
\[
 \alpha_1=1,\qquad\alpha_2=\frac16,\qquad
 \alpha_3=\frac1{21},\qquad\alpha_4=\frac1{60}.
\]
Thus even the first two inverse sectors can be evaluated with no special functions beyond a logarithm and a square root, once $\Pinf$ is known.

\section{Complete flags and the inversion statistic on permutations}

A complete flag in $\F_q^n$ is a chain
$0=V_0\subset V_1\subset\cdots\subset V_n=\F_q^n$ with $\dim V_j=j$.
Each flag has exactly
$(q-1)^n q^{n(n-1)/2}$ ordered bases adapted to it: the $j$th vector may be any element of $V_j\setminus V_{j-1}$.
Dividing \eqref{t3:eq:GLcount} by this number gives
\begin{equation}\label{t3:eq:qfactorial}
 [n]_q!=\prod_{j=1}^n\frac{q^j-1}{q-1}
 =q^{n(n+1)/2}(q-1)^{-n}\pp{n}.
\end{equation}
At $q=2$ these are the products $\prod_{j=1}^n(2^j-1)$ in A005329 \cite{oeisflags}.

There is also a weighted permutation interpretation. If $\operatorname{inv}(\pi)$ denotes the inversion number of a permutation, insertion of the largest element creates any one of $0,1,\ldots,n-1$ new inversions. Induction therefore proves
\begin{equation}\label{t3:eq:Mahonian}
 \sum_{\pi\in S_n}q^{\operatorname{inv}(\pi)}=[n]_q!.
\end{equation}
This relates a finite-field count to an ordinary combinatorial generating polynomial.

For the analytic interpolation, use
\begin{equation}\label{t3:eq:flagphase}
 a=\frac h2,\quad b=\frac h2-\log(q-1),\quad
 c=\log\Pinf,\quad B(t)=\A(t),\quad\lambda=h.
\end{equation}
The center is the quadratic root \eqref{t3:eq:center}, and all inverse coefficients follow from the same $\alpha_m$ used for general linear groups. Only $a$ and $D$ change. In particular,
\begin{equation}\label{t3:eq:flaginverse}
 ([x]_q!)^{-1}(Y)
 =s-\frac{\alpha_1}{D}Q^s
 -\left(\frac{\alpha_2}{D}
   +\frac{h\alpha_1^2}{D^2}
   +\frac{h\alpha_1^2}{2D^3}\right)Q^{2s}
 +O(Q^{3s}/D).
\end{equation}
Here the notation on the left means the inverse of the interpolation just defined, not the reciprocal of a factorial.

\section{The other side of \texorpdfstring{$q=1$}{q = 1}: an exponential rather than quadratic phase}
\label{t3:sec:qfactorial-smallbase}

Use $Q\in(0,1)$ as the base of the usual $q$-factorial. Then
\begin{equation}\label{t3:eq:smallbase-factorial}
 [n]_Q!=\frac{\pp{n}}{(1-Q)^n},\qquad
 f_Q(x)=\Pinf\exp\bigl(\beta x+\A(Q^x)\bigr),
 \quad \beta=-\log(1-Q)>0.
\end{equation}
The logarithmic phase is now linear. Set
\begin{equation}\label{t3:eq:smallbase-center}
 s=\frac{\log(Y/\Pinf)}\beta,\qquad
 \rho=\frac h\beta,\qquad
 t_0=Q^s=(Y/\Pinf)^{-\rho}.
\end{equation}
The inverse is a convergent expansion in a generally nonintegral power of $Y$:
\begin{align}\label{t3:eq:smallbase-inverse}
 f_Q^{-1}(Y)=s
 -\frac{\alpha_1}{\beta}t_0
 -\left(\frac{\alpha_2}{\beta}
       +\frac{h\alpha_1^2}{\beta^2}\right)t_0^2
 +O(t_0^3).
\end{align}
The exact all-order formula simplifies to
\begin{equation}\label{t3:eq:linear-extraction}
 d_m=-\frac1{mh}\coefB{u}{m}\exp\left(\frac{mh}{\beta}\A(u)\right).
\end{equation}

Consequently the same polynomial sequence $[n]_q!$, viewed at a fixed real base, has three different large-index regimes. For $q>1$ its logarithm is quadratic in $n$; for $0<q<1$ it is linear; and at $q=1$ it is $\log(n!)$, with dominant term $n\log n$. This is a genuine change of asymptotic scaling, not something obtained by substituting $q=1$ into a fixed-$q$ inverse transseries. Section~\ref{t3:sec:double} studies a related simultaneous limit in detail.
% ---- hand chunk 08-gauss-head.tex ----

\chapter{Gaussian binomial coefficients and prescribed partial flags}
\label{ct:ch:gaussian}

% ---- T3 lines 543--675 ----
\label{t3:sec:gaussian}

\section{Counting subspaces}

For integers $0\le k\le n$, the number of $k$-dimensional subspaces of $\F_q^n$ is
\begin{equation}\label{t3:eq:qbin-count}
 \qbinq{n}{k}
 =\frac{\prod_{j=0}^{k-1}(q^n-q^j)}
        {\prod_{j=0}^{k-1}(q^k-q^j)}
 =q^{k(n-k)}\frac{\pp{n}}{\pp{k}\pp{n-k}}.
\end{equation}
The first quotient counts linearly independent ordered $k$-tuples and divides by the number of bases of each subspace. This proves both the enumerative interpretation and the product identity. The definition agrees with \texttt{QBinomial[n,k,q]} \cite{wlqbinomial}.

\section{Fixed rank: algebraic inverse series}
\label{t3:sec:fixedrank}

Fix $k\ge1$. The finite-product interpolation is
\begin{equation}\label{t3:eq:fixedrank-interp}
 H_{k,q}(x)=\frac{q^{k(x-k)}}{\pp{k}}
              \prod_{j=0}^{k-1}(1-q^jQ^x),\qquad x>k-1.
\end{equation}
It is positive and increasing in this range, as is evident from the first quotient in \eqref{t3:eq:qbin-count} with $n$ replaced by $x$. In the notation of the inversion theorem,
\begin{align}
 P(x)&=khx-hk^2-\log\pp{k},\label{t3:eq:fixedrank-phase}\\
 B(t)&=\sum_{j=0}^{k-1}\log(1-q^jt),\qquad
 b_m=-\frac{q^{km}-1}{m(q^m-1)}.\label{t3:eq:fixedrank-tail}
\end{align}
The center and small parameter are
\begin{equation}\label{t3:eq:fixedrank-center}
 s=\frac{\log Y+hk^2+\log\pp{k}}{kh},
 \qquad t_0=q^{-k}\bigl(\pp{k}Y\bigr)^{-1/k}.
\end{equation}
Thus a convergent inverse Puiseux expansion follows:
\begin{equation}\label{t3:eq:fixedrank-inverse}
 H_{k,q}^{-1}(Y)
 =s+\sum_{m\ge1}d_m t_0^m,
 \qquad
 d_m=-\frac1{mh}\coefB{u}{m}
       \prod_{j=0}^{k-1}(1-q^ju)^{m/k}.
\end{equation}
All powers are expanded at $u=0$ with constant term one. In particular,
\begin{equation}\label{t3:eq:fixedrank-first}
 H_{k,q}^{-1}(Y)=s+\frac{[k]_q}{kh}\,t_0+O(t_0^2).
\end{equation}
Although the leading inverse contains a logarithm, its correction coefficients are constants: $D=kh$ is fixed.

There is an independent algebraic explanation of convergence. Writing $z=q^x$, the equation is
\begin{equation}\label{t3:eq:fixedrank-algebraic}
 \prod_{j=0}^{k-1}(z-q^j)=Y\prod_{j=0}^{k-1}(q^k-q^j).
\end{equation}
After scaling $z$ by $Y^{1/k}$, the desired large positive root is analytic in $Y^{-1/k}$ near zero, because the limiting polynomial has a simple positive root. Taking a logarithm gives \eqref{t3:eq:fixedrank-inverse}.

For $k=1$ this reduces to
\[
 H_{1,q}^{-1}(Y)=\frac{\log(1+(q-1)Y)}h.
\]
For $k=2$, put $K_q=(q^2-1)(q^2-q)$. The inverse is exactly
\begin{equation}\label{t3:eq:k2-inverse}
 H_{2,q}^{-1}(Y)=\frac1h\log\left(
 \frac{q+1+\sqrt{(q-1)^2+4K_qY}}2\right).
\end{equation}
These elementary cases provide useful checks on the coefficient formula.

\section{Proportionally scaled rank: a quadratic inverse center}

Let $u,v$ be fixed positive integers. Define
\begin{equation}\label{t3:eq:scaled-gaussian}
 H_{u,v}(x)=\frac{q^{uvx^2}}{\Pinf}
 \exp\bigl(\A(t^{u+v})-\A(t^u)-\A(t^v)\bigr),
 \qquad t=Q^x.
\end{equation}
At integers $x=n$ it is $\qbinq{(u+v)n}{un}$. The logarithmic coefficients are
\begin{equation}\label{t3:eq:scaled-gaussian-b}
 b_m=\one_{u+v\mid m}\alpha_{m/(u+v)}
     -\one_{u\mid m}\alpha_{m/u}
     -\one_{v\mid m}\alpha_{m/v}.
\end{equation}
Use $a=uvh$, $b=0$, $c=-\log\Pinf$, and $\lambda=h$ in the inversion theorem. The first exponential weight is $\min(u,v)$; if $u=v$, its coefficient is doubled. The available inverse weights belong to the additive semigroup generated by the nonzero degrees in \eqref{t3:eq:scaled-gaussian-b}.

For the central case $u=v=1$,
\begin{equation}\label{t3:eq:central-B}
 B(t)=\A(t^2)-2\A(t),\qquad
 b_m=-2\alpha_m+\one_{2\mid m}\alpha_{m/2}.
\end{equation}
The corresponding $q=2$ sequence is A006098 \cite{oeiscentral}. Its center is
\begin{equation}\label{t3:eq:central-center}
 s=\sqrt{\frac{\log(Y\Pinf)}h},\qquad D=2hs.
\end{equation}
The first two inverse sectors are
\begin{align}\label{t3:eq:central-inverse}
 H_{1,1}^{-1}(Y)=s
 &+\frac{2\alpha_1}{D}Q^s\\[-0.3em]
 &-\left(\frac{\alpha_1-2\alpha_2}{D}
       +\frac{4h\alpha_1^2}{D^2}
       +\frac{4h\alpha_1^2}{D^3}\right)Q^{2s}
 +O(Q^{3s}/D).\notag
\end{align}
The leading correction has the opposite sign from that for $\GL(n,q)$.

An explicit forward coefficient formula follows directly from the Euler expansions:
\begin{equation}\label{t3:eq:central-forward}
 \ee^{B(t)}=\frac{(Qt;Q)_\infty^2}{(Qt^2;Q)_\infty},\qquad
 [t^m]\ee^{B(t)}=
 \sum_{r+s+2j=m}
 \frac{(-1)^{r+s}Q^{r(r+1)/2+s(s+1)/2+j}}
 {\pp{r}\pp{s}\pp{j}}.
\end{equation}
The sum ranges over nonnegative integers and is finite. Equations \eqref{t3:eq:central-B} and \eqref{t3:eq:finite-dm} give a similarly finite formula for every inverse coefficient.

\section{Prescribed partial flags}

Fix positive integers $r_1,\ldots,r_\ell$ and write
\[
 R=r_1+\cdots+r_\ell,\qquad
 A_0=\sum_{1\le i<j\le\ell}r_i r_j.
\]
A flag with successive quotient dimensions $r_1n,\ldots,r_\ell n$ is counted by the Gaussian multinomial coefficient
\begin{equation}\label{t3:eq:multinomial}
 \genfrac{[}{]}{0pt}{}{Rn}{r_1n,\ldots,r_\ell n}_q
 =q^{A_0n^2}\frac{\pp{Rn}}{\prod_i\pp{r_i n}}.
\end{equation}
It follows either by successive subspace choices or by dividing adapted bases by the relevant block-upper-triangular stabilizer. The interpolation has
\begin{equation}\label{t3:eq:partialflag-phase}
 P(x)=A_0hx^2+(1-\ell)\log\Pinf,
 \qquad
 B(t)=\A(t^R)-\sum_{i=1}^{\ell}\A(t^{r_i}).
\end{equation}
Thus
\begin{equation}\label{t3:eq:partialflag-b}
 b_m=\one_{R\mid m}\alpha_{m/R}
     -\sum_{i:r_i\mid m}\alpha_{m/r_i}.
\end{equation}
All proportional partial-flag families are consequently covered by one inversion theorem and one explicit arithmetic coefficient sequence. Repeated block sizes automatically produce their correct multiplicities; they should not be reduced to distinct sizes before applying this formula.
% ---- hand chunk 08-gauss-t1.tex ----

\section{The central coefficient revisited: a convergent forward transseries and the inverse as a differential operator}
\label{ct:sec:gauss-t1}

The first source treated the central case $u=v=1$ independently, writing
$\mathcal B_q(x)=\qbin{2x}{x}{q}$ for the interpolation
\cref{t3:eq:scaled-gaussian} and $a=h$, $t=Q^x$, $P=\Pinf$.  Its forward
transseries is the Euler-expanded form of \cref{t3:eq:central-forward}, its
logarithmic coefficients $h_r$ are the $b_m$ of \cref{t3:eq:central-B}, and
its inverse theorem is the operator form \cref{t3:eq:general-derivative} of
the convergent theorem, with $P(x)=ax^2$; its first two coefficients
$d_1,d_2$ coincide with \cref{t3:eq:central-inverse} after $D=2aX$ is
substituted.  The section is kept because the operator form, the
convergence-radius argument, and the warning about the limit $q\downarrow1$
are stated there most directly.

% ---- T1 lines 1080--1176 ----
\subsection{A convergent forward transseries}
Write $a=\log q$, $t=Q^x$. Expansion of $\log(1-z)$ in the infinite products gives the exact identity
\begin{equation}
 \mathcal B_q(x)=\frac{\ee^{ax^2}}P\exp H(t),\qquad
 H(t)=-\sum_{m\ge1}\frac{2t^m-t^{2m}}{m(q^m-1)}.
 \label{t1:eq:qbinH}
\end{equation}
The exchanges of the two sums are justified by absolute convergence for $|t|<\sqrt q$, in particular for the physical interval $0<t<1$. The coefficient form is
\begin{equation}
 H(t)=\sum_{r\ge1}h_rt^r,\qquad
 h_r=-\frac2{r(q^r-1)}+
 \mathbf1_{2\mid r}\frac2{r(q^{r/2}-1)}.
 \label{t1:eq:qbinh}
\end{equation}
Thus $h_1=-2/(q-1)$ and $h_2=q/(q^2-1)$. The complete forward transseries is the genuinely convergent expansion
\begin{equation}
 \mathcal B_q(x)=\frac{q^{x^2}}P
 \sum_{r\ge0}\ExpC_r(h_1,h_2,\ldots)q^{-rx}.
 \label{t1:eq:qbinforward}
\end{equation}
Unlike a Stirling series, this is a Taylor series of an exact analytic amplitude. The first singularity of $H$ in modulus lies at $t=\sqrt q$, coming from a denominator factor; the numerator factors first vanish at modulus $q$. In particular the radius stated above is not just a formal estimate.

The derivative of the logarithm is
\[
 (\log\mathcal B_q)'(x)=2ax-atH'(t)=2ax+\OO(q^{-x}),
\]
which is positive for large $x$. Hence a unique large real inverse is defined.

\subsection{An all-order inverse with a quadratic phase}
Set
\begin{equation}
 L=\log(yP),\qquad X=\sqrt{L/a},\qquad T=\ee^{-aX},
 \label{t1:eq:qbincore}
\end{equation}
for sufficiently large $y$. The carrier is exact for $q^{x^2}/P$. All remaining inverse corrections are exponentially small in $X$; there is no separate algebraic tail at this stage.

For $1\le m\le n$, define the finite coefficient
\begin{equation}
 H_{n,m}=\coefA{t^n}H(t)^m
 =\sum_{\substack{r_1+\cdots+r_m=n\\r_j\ge1}}h_{r_1}\cdots h_{r_m},
 \label{t1:eq:qHnm}
\end{equation}
and the differential operator, acting on functions of $X$,
\begin{equation}
 \mathfrak D_n=\frac1{2aX}(\partial_X-an).
 \label{t1:eq:qoperator}
\end{equation}

\begin{theorem}[Convergent quadratic-core inverse]
\label{t1:thm:qbininverse}
For sufficiently large $y$,
\begin{equation}
 \mathcal B_q^{-1}(y)=X+\sum_{n\ge1}d_n(X)T^n,
 \label{t1:eq:qbininverse}
\end{equation}
where
\begin{equation}
 d_n(X)=\sum_{m=1}^n\frac{(-1)^m}{m!}H_{n,m}
 \mathfrak D_n^{m-1}\left(\frac1{2aX}\right).
 \label{t1:eq:qbininversecoeff}
\end{equation}
This is a convergent expansion at the physical $T=\ee^{-aX}$, not merely a formal asymptotic assertion. The coefficients are finite rational expressions in $X$, $a$, and $q$.
\end{theorem}
\begin{proof}
Use \cref{t1:thm:phase} with $\phi(x)=ax^2$ and $R(x)=H(\ee^{-ax})$. Then $g=\sqrt{L/a}=X$, $g'=1/(2aX)$, and $D=(2aX)^{-1}\partial_X$. Expanding $H(T)^m$ using \eqref{t1:eq:qHnm}, observe that
\[
 D\bigl(T^n f(X)\bigr)=T^n\mathfrak D_n f(X).
\]
Collecting $T^n$ yields \eqref{t1:eq:qbininversecoeff}.

For a direct convergence argument, put $x=X+\delta$ and temporarily regard $T$ as independent of $X$. The equation to solve is
\[
 2aX\delta+a\delta^2+H(T\ee^{-a\delta})=0.
\]
At $(T,\delta)=(0,0)$ the slope is $2aX$. On fixed small disks in $T$ and $\delta$, analyticity of $H$ and this growing slope give a contraction for all sufficiently large $X$. Thus the solution is analytic in $T$ with a radius bounded below independently of large $X$. The physical $T$ eventually lies inside it. Uniqueness identifies its Taylor coefficients with the formal coefficients already derived.
\end{proof}

The first two coefficient functions simplify to
\begin{align}
 d_1(X)&=\frac1{a(q-1)X},\nonumber\\
 d_2(X)&=-\frac{q}{2a(q^2-1)X}
 -\frac1{a(q-1)^2X^2}
 -\frac1{2a^2(q-1)^2X^3}.
 \label{t1:eq:qbind12}
\end{align}
It is essential to apply \eqref{t1:eq:qoperator} as a variable-coefficient differential operator. Treating $1/(2aX)$ as a constant at successive differentiations loses the last terms in \eqref{t1:eq:qbind12} and all of their higher-order analogues.

In terms of the target $y$, the basic inverse scale is
\begin{equation}
 T=\exp\!\left(-\sqrt{(\log q)\log(yP)}\right).
 \label{t1:eq:qbinscale}
\end{equation}
This is slower than any fixed power $y^{-c}$ but faster than every inverse power of $\log y$. Therefore it is not captured by a power--logarithmic expansion alone.

\begin{warning}[The limit $q\downarrow1$ is not uniform]
Every theorem in this section holds for fixed $q>1$. The coefficients $h_r$ and $d_n$ contain powers of $(q^r-1)^{-1}$, while $P=(q^{-1};q^{-1})_\infty$ changes rapidly as $q\downarrow1$. Substituting $q=1$ in these formulas does not recover the ordinary central binomial asymptotics. A joint limit must specify how $\log q$ scales with $x$ and rederive the leading phase from the product. This is a different asymptotic problem, not a removable singularity in the fixed-$q$ expansion.
\end{warning}
% ---- hand chunk 08-catalan-head.tex ----

\chapter{$q$-Catalan numbers and the symplectic and unitary orders}
\label{ct:ch:qcatalan}

% ---- T3 lines 678--867 ----
\label{t3:sec:catalan}

\section{Definition and polynomiality}

There are several inequivalent $q$-Catalan conventions. Here we use only the binomial quotient
\begin{equation}\label{t3:eq:qfuss-definition}
 C_{r,n}(q)=\frac1{[rn+1]_q}\qbinq{(r+1)n}{n},
 \qquad r\in\N_{>0},\qquad
 [m]_q=\frac{q^m-1}{q-1}.
\end{equation}
At $r=1$ and $q=2$ this is A015030 \cite{oeisqcatalan}. At $q=1$, by continuity, it is the usual Fuss--Catalan number
$\binom{(r+1)n}{n}/(rn+1)$. No identification with a different area-weighted or recursively defined $q$-Catalan polynomial is intended.

\begin{proposition}\label{t3:prop:qfuss-polynomial}
For positive integers $r,n$, the expression \eqref{t3:eq:qfuss-definition} is a polynomial in $q$ with integer coefficients.
\end{proposition}
\begin{proof}
The exponent of the cyclotomic polynomial $\Phi_d(q)$ in
$\qbinq{(r+1)n}{n}$ is
\[
 e_d=\left\lfloor\frac{(r+1)n}{d}\right\rfloor
     -\left\lfloor\frac n d\right\rfloor
     -\left\lfloor\frac{rn}{d}\right\rfloor\in\{0,1\}.
\]
This follows by factoring each $q^j-1$ into cyclotomic factors. A factor of $[rn+1]_q$ has $d>1$ and $d\mid rn+1$. Write $rn=ad-1$. Then
$e_d=1+\lfloor(n-1)/d\rfloor-\lfloor n/d\rfloor=1$, because $d\nmid n$ follows from $\gcd(n,rn+1)=1$. Hence every denominator factor cancels, and all remaining cyclotomic exponents are nonnegative integers.
\end{proof}
This proof establishes integer coefficients. No coefficient-positivity assertion is needed in the analysis below; positivity of the quotient for real $q>1$ follows directly from its factors.

\section{An exact logarithmic tail}

From \eqref{t3:eq:scaled-gaussian}, and
$[rx+1]_q=q^{rx+1}(1-Qt^r)/(q-1)$, obtain the interpolation
\begin{equation}\label{t3:eq:qfuss-interp}
 C_{r,q}(x)=\frac{1-Q}{\Pinf}\,q^{r x(x-1)}\ee^{B_r(t)},
 \qquad t=Q^x,
\end{equation}
where
\begin{equation}\label{t3:eq:qfuss-B}
 B_r(t)=\A(t^{r+1})-\A(t)-\A(t^r)-\log(1-Qt^r).
\end{equation}
Consequently
\begin{equation}\label{t3:eq:qfuss-b}
 b_m=\one_{r+1\mid m}\alpha_{m/(r+1)}-\alpha_m
 -\one_{r\mid m}\alpha_{m/r}
 +\one_{r\mid m}\frac{Q^{m/r}}{m/r}.
\end{equation}
This formula includes $r=1$ correctly; two copies of $-\A(t)$ then occur.

For $r=1$, the leading asymptotic specializes to
\begin{equation}\label{t3:eq:qCatalan-leading}
 C_{1,n}(2)\sim\frac1{\mathcal P_{1/2}}\,2^{n^2-n-1}.
\end{equation}
This is the asymptotic constant recorded by Vladimir Reshetnikov in the 2021 contribution to A015030 \cite{oeisqcatalan}. Formula \eqref{t3:eq:qfuss-b} supplies all of its exponential corrections, while Theorem~\ref{t3:thm:inverse} supplies all inverse corrections.

\section{Inverse formulas}

The exact quadratic center is
\begin{equation}\label{t3:eq:qfuss-center}
 s=\frac12+\sqrt{\frac14+
       \frac{\log\bigl(Y\Pinf/(1-Q)\bigr)}{rh}},
 \qquad D=rh(2s-1).
\end{equation}
Use $a=rh$ and $\lambda=h$ in \eqref{t3:eq:finite-dm}. For $r=1$,
\begin{equation}\label{t3:eq:qCatalan-b12}
 b_1=-\frac{Q(1+Q)}{1-Q},\qquad
 b_2=\alpha_1-2\alpha_2+\frac{Q^2}{2}.
\end{equation}
Hence, with $\gamma=Q(1+Q)/(1-Q)$,
\begin{align}\label{t3:eq:qCatalan-inverse}
 C_{1,q}^{-1}(Y)=s
 &+\frac{\gamma}{D}Q^s\\[-0.3em]
 &-\left(\frac{b_2}{D}
       +\frac{h\gamma^2}{D^2}
       +\frac{h\gamma^2}{D^3}\right)Q^{2s}
 +O(Q^{3s}/D).\notag
\end{align}
For $r>1$, $b_1=-\alpha_1$, so the first correction is
\begin{equation}\label{t3:eq:qFuss-first}
 C_{r,q}^{-1}(Y)=s+\frac{\alpha_1}{rh(2s-1)}Q^s
                  +O(Q^{2s}/s).
\end{equation}
The case $r=1$ is therefore not obtained by treating all terms in \eqref{t3:eq:qfuss-B} as having distinct weights: a collision occurs at the very first weight.

For example, at $q=2$, $r=1$, the coefficients are
\[
 b_1=-\frac32,\qquad b_2=\frac{19}{24},\qquad
 b_3=-\frac{3}{56}.
\]
The source and inverse coefficients are rational functions of $Q$, together with the powers of $h$ and $D$ dictated by the universal formula. There is no unresolved recurrence in these expressions.

\chapter{Symplectic and unitary orders: missing sectors and parity}
\label{t3:sec:classical}

\section{Symplectic groups}

Let $V$ be a $2n$-dimensional symplectic vector space over $\F_q$. Choose the first vector $e_1\ne0$ of a symplectic basis in $q^{2n}-1$ ways. The equation $\langle e_1,f_1\rangle=1$ then has $q^{2n-1}$ solutions. The symplectic complement of their span has dimension $2n-2$. Iterating this argument proves
\begin{equation}\label{t3:eq:Sp-count}
 s_n(q)=|\Sp(2n,q)|
 =q^{n^2}\prod_{j=1}^n(q^{2j}-1)
 =q^{2n^2+n}(Q^2;Q^2)_n.
\end{equation}
The $q=2$ specialization is A003923 \cite{oeisSp}.

Choose the interpolation
\begin{equation}\label{t3:eq:Sp-interp}
 s_q(x)=\mathcal P_{Q^2}
   \exp\bigl(h(2x^2+x)+\mathcal A_{Q^2}(Q^{2x})\bigr).
\end{equation}
With $t=Q^x$, its logarithmic tail is
\begin{equation}\label{t3:eq:Sp-B}
 B(t)=\mathcal A_{Q^2}(t^2),\qquad
 b_{2m}=\frac{Q^{2m}}{m(1-Q^{2m})},\qquad b_{2m-1}=0.
\end{equation}
The inverse therefore contains no odd weights in $Q^s$. Put
\begin{equation}\label{t3:eq:Sp-center}
 s=\frac{-1+\sqrt{1+8\log(Y/\mathcal P_{Q^2})/h}}4,
 \qquad D=h(4s+1),\qquad \beta_m=\frac{Q^{2m}}{m(1-Q^{2m})}.
\end{equation}
Using the primitive small parameter $Q^{2s}$, so that $\lambda=2h$, gives
\begin{align}\label{t3:eq:Sp-inverse}
 s_q^{-1}(Y)=s
 &-\frac{\beta_1}{D}Q^{2s}\\[-0.3em]
 &-\left(\frac{\beta_2}{D}
       +\frac{2h\beta_1^2}{D^2}
       +\frac{2h\beta_1^2}{D^3}\right)Q^{4s}
 +O(Q^{6s}/D).\notag
\end{align}
This illustrates why the exponential scale should be determined from the source coefficients, not guessed from a generic family resemblance.

\section{Unitary isometry groups and the order formula}

Let $U_n(q)$ denote the isometry group of a nondegenerate Hermitian form on $\F_{q^2}^n$. This is the full unitary isometry group, not the determinant-one subgroup and not a group of similitudes. The corresponding OEIS convention uses $\mathrm{GU}(n,q)$ \cite{oeisU}.

We include a count to fix all powers of $q$ and all constants. For the standard Hermitian form, write $N(v)=\sum_i v_i^{q+1}\in\F_q$. The norm map $\F_{q^2}^{\times}\to\F_q^{\times}$ has $q+1$ preimages for each nonzero value. If $\psi$ is a nontrivial additive character of $\F_q$, then
\[
 \sum_{z\in\F_{q^2}}\psi(z^{q+1})
 =1+(q+1)\sum_{a\in\F_q^{\times}}\psi(a)=-q.
\]
By character orthogonality, the number of vectors of any prescribed nonzero norm is
\begin{equation}\label{t3:eq:unitnorm}
 \frac{q^{2n}-(-q)^n}{q}=q^{n-1}(q^n-(-1)^n).
\end{equation}
The orthogonal complement of a unit vector is again a nondegenerate Hermitian space of one smaller dimension. Counting orthonormal bases therefore yields
\begin{equation}\label{t3:eq:U-count}
 u_n(q)=|U_n(q)|
 =q^{n(n-1)/2}\prod_{j=1}^n(q^j-(-1)^j)
 =q^{n^2}\prod_{j=1}^n(1-(-Q)^j).
\end{equation}
The reduction to the standard Hermitian form follows by the usual orthogonal basis construction; every nonzero diagonal norm can be rescaled to one because the norm map is onto.

\section{Two real analytic sheets}

Define the positive constant
\begin{equation}\label{t3:eq:Pminus}
 \mathcal P_{-Q}=\prod_{j=1}^{\infty}(1-(-Q)^j)
 =(Q^2;Q^2)_\infty(-Q;Q^2)_\infty.
\end{equation}
For $\eps\in\{0,1\}$ define
\begin{align}
 B_{\eps}(t)&=\sum_{m\ge1}
 \frac{(-1)^{m(\eps+1)}Q^m}{m(1-(-Q)^m)}t^m,
 \label{t3:eq:U-B}\\
 u_{q,\eps}(x)&=\mathcal P_{-Q}
                  \exp\bigl(hx^2+B_{\eps}(Q^x)\bigr).
 \label{t3:eq:U-sheet}
\end{align}
The series in \eqref{t3:eq:U-B} is analytic for $|t|<Q^{-1}$. At an integer $n\equiv\eps\pmod2$, the omitted tail of the finite product in \eqref{t3:eq:U-count} starts with $(-Q)^{n+1}$; expansion of its logarithm gives exactly \eqref{t3:eq:U-B}. Thus each sheet matches its specified parity.

Both inverse centers are
\begin{equation}\label{t3:eq:U-center}
 s=\sqrt{\frac{\log(Y/\mathcal P_{-Q})}h},\qquad D=2hs.
\end{equation}
Since $b_1=(-1)^{\eps+1}Q/(1+Q)$,
\begin{equation}\label{t3:eq:U-first}
 u_{q,\eps}^{-1}(Y)
 =s+(-1)^{\eps}\frac{Q}{(1+Q)D}Q^s+O(Q^{2s}/D).
\end{equation}
The even sheet lies above its center and the odd sheet below. Averaging the two sheets would erase the leading correction, although that correction is present on each integer subsequence.

At $q=2$,
\begin{equation}\label{t3:eq:U-constant-check}
 \mathcal P_{-1/2}=1.210724130301059180136172856\ldots.
\end{equation}
Thus $u_n(2)\sim\mathcal P_{-1/2}2^{n^2}$. If this is written as
$c\,2^{n^2-1}$, the equivalent constant is
$c=2\mathcal P_{-1/2}=2.42144826060211836027\ldots$.
The product \eqref{t3:eq:U-count} fixes this normalization independently of the form in which an asymptotic comment is tabulated.

All higher parity-dependent corrections follow from \eqref{t3:eq:finite-dm} with the coefficients \eqref{t3:eq:U-B}. Integer thresholds are obtained from \eqref{t3:eq:residue-rounding} with $r=2$, followed by the minimum over the two parities.
% ---- hand chunk 08-galois-head.tex ----

\chapter{Galois numbers and root-lattice theta sectors}
\label{ct:ch:galois}

% ---- T3 lines 870--1138 ----
\label{t3:sec:galois}

\section{From one Grassmannian to all subspaces}

The ordinary Galois number is the number of all subspaces of $\F_q^n$:
\begin{equation}\label{t3:eq:Galois-def}
 G_n(q)=\sum_{k=0}^n\qbinq{n}{k}.
\end{equation}
At $q=2$ this is A006116 \cite{oeisGalois}:
\[
 1,\ 2,\ 5,\ 16,\ 67,\ 374,\ 2825,\ 29212,\ldots.
\]
The terms near $k=n/2$ have the largest exponential size, but there is no growing Gaussian width: the ratio of the dominant quadratic factors at displacement $j$ is $q^{-j^2}$. Summing these displacements therefore gives a theta constant. The parity of $n$ determines whether the displacements are integral or half-integral.

Define
\begin{equation}\label{t3:eq:theta2}
 \T_0(Q)=\sum_{j\in\Z}Q^{j^2},\qquad
 \T_1(Q)=\sum_{j\in\Z+1/2}Q^{j^2}.
\end{equation}
The leading form is
\begin{equation}\label{t3:eq:Galois-leading}
 G_n(q)\sim\frac{\T_{n\bmod2}(Q)}{\Pinf}\,q^{n^2/4}
 \qquad(n\to\infty\text{ within a fixed parity}).
\end{equation}
Theta leading constants in this enumeration are known \cite{kousidis}. We now prove a stronger exact representation with explicit coefficients at every exponential weight, and extend it at once to flags of arbitrary fixed length.

\section{Generalized Galois numbers}

For $r\ge2$, define
\begin{equation}\label{t3:eq:generalGalois-def}
 G_n^{(r)}(q)=\sum_{\substack{k_1,\ldots,k_r\ge0\\k_1+\cdots+k_r=n}}
 \genfrac{[}{]}{0pt}{}{n}{k_1,\ldots,k_r}_q.
\end{equation}
This counts weak flags
\[
 0\subseteq V_1\subseteq\cdots\subseteq V_{r-1}\subseteq\F_q^n;
\]
the successive quotient dimensions are the $k_i$, so repeated subspaces are permitted. For $r=2$ it reduces to \eqref{t3:eq:Galois-def}. For $r=3$, $q=2$, its first terms are
\[
 1,\ 3,\ 12,\ 66,\ 513,\ 5769,\ 95706,\ 2379348,\ldots.
\]
The parameter $r$ is fixed in all asymptotic statements below.

For an integer $\eps$ define the lattice coset
\begin{equation}\label{t3:eq:Lattice}
 \mathcal L_{\eps}^{(r)}=
 \left\{(z_1-\eps/r,\ldots,z_r-\eps/r):
 z_i\in\Z,\ \sum_i z_i=\eps\right\}.
\end{equation}
It lies in the hyperplane $\sum_i j_i=0$. Increasing $\eps$ by $r$ leaves the set unchanged, so the subscript is understood modulo $r$. For $\eps=0$, this is the $A_{r-1}$ root lattice in its usual Euclidean realization. Put
\begin{equation}\label{t3:eq:root-theta}
 \T_{\eps}^{(r)}(Q)=
 \sum_{\boldsymbol j\in\mathcal L_{\eps}^{(r)}}
 Q^{\|\boldsymbol j\|^2/2}>0.
\end{equation}
These series converge absolutely, since a positive definite quadratic form dominates the number of lattice points in successive shells. Reflection of the lattice gives
$\T_{\eps}^{(r)}=\T_{r-\eps}^{(r)}$. At $r=2$, definition \eqref{t3:eq:root-theta} is exactly \eqref{t3:eq:theta2}.

\section{An entire function containing all correction coefficients}

Let
\begin{equation}\label{t3:eq:comp-coeff}
 A_{r,M}(Q)=
 \sum_{\substack{m_1,\ldots,m_r\ge0\\m_1+\cdots+m_r=M}}
 \prod_{i=1}^r\frac1{\pp{m_i}},\qquad A_{r,0}=1.
\end{equation}
This is a finite sum. Define
\begin{equation}\label{t3:eq:h-coeff}
 h_{\eps,M}^{(r)}=
 (-1)^M Q^{M^2/(2r)+M/2}\T_{\eps+M}^{(r)}(Q)\,A_{r,M}(Q),
\end{equation}
and
\begin{equation}\label{t3:eq:H-entire}
 H_{\eps}^{(r)}(t)=\sum_{M\ge0}h_{\eps,M}^{(r)}t^M.
\end{equation}

\begin{theorem}[Exact theta-tail representation]\label{t3:thm:theta-tail}
The function $H_{\eps}^{(r)}$ is entire. For every nonnegative integer $n\equiv\eps\pmod r$,
\begin{equation}\label{t3:eq:Galois-exact}
 \boxed{\displaystyle
 G_n^{(r)}(q)=\Pinf^{1-r}
 q^{(r-1)n^2/(2r)}
 \exp\!\bigl(\A(t^r)\bigr)H_{\eps}^{(r)}(t),
 \qquad t=Q^{n/r}.}
\end{equation}
An equivalent entire representation of the tail is
\begin{equation}\label{t3:eq:H-bilateral}
 H_{\eps}^{(r)}(t)=
 \sum_{\boldsymbol j\in\mathcal L_{\eps}^{(r)}}
 Q^{\|\boldsymbol j\|^2/2}
 \prod_{i=1}^r(QtQ^{j_i};Q)_\infty.
\end{equation}
The series in \eqref{t3:eq:H-bilateral} converges absolutely and locally uniformly for all complex $t$.
\end{theorem}

\begin{proof}
Start with the multinomial product
\[
 \genfrac{[}{]}{0pt}{}{n}{k_1,\ldots,k_r}_q
 =q^{(n^2-\sum_i k_i^2)/2}
     \frac{\pp{n}}{\prod_i\pp{k_i}}.
\]
Write $j_i=k_i-n/r$. Then $\boldsymbol j\in\mathcal L_{\eps}^{(r)}$ and
\begin{equation}\label{t3:eq:quadratic-completion}
 \frac{n^2-\sum_i k_i^2}{2}
 =\frac{r-1}{2r}n^2-\frac12\|\boldsymbol j\|^2.
\end{equation}
Using \eqref{t3:eq:product-interpolation}, each denominator contributes
$\Pinf^{-1}(Q^{k_i+1};Q)_\infty$, and the numerator contributes
$\Pinf\exp\A(Q^n)$. This gives \eqref{t3:eq:Galois-exact} with the sum in \eqref{t3:eq:H-bilateral} initially restricted to $k_i\ge0$.

At $t=Q^{n/r}$, a lattice point with any $k_i<0$ contributes zero: the product $(Q^{k_i+1};Q)_\infty$ contains the factor $1-Q^0$. Thus the restricted sum may be extended to the whole lattice coset, once convergence is established.

Expand each product in \eqref{t3:eq:H-bilateral} by \eqref{t3:eq:euler-entire}. For a multi-index $\boldsymbol m=(m_1,\ldots,m_r)$ with $M=\sum_i m_i$, its contribution contains the exponent
\[
 \frac12\|\boldsymbol j\|^2
 +\sum_i m_i j_i+\frac12\sum_i m_i(m_i+1).
\]
Complete the square:
\begin{equation}\label{t3:eq:theta-complete}
 \frac12\left\|\boldsymbol j+\boldsymbol m-\frac Mr\boldsymbol1\right\|^2
 +\frac{M^2}{2r}+\frac M2.
\end{equation}
Translation by $\boldsymbol m-(M/r)\boldsymbol1$ carries
$\mathcal L_{\eps}^{(r)}$ bijectively to $\mathcal L_{\eps+M}^{(r)}$.
The lattice sum at fixed $\boldsymbol m$ is therefore exactly
$Q^{M^2/(2r)+M/2}\T_{\eps+M}^{(r)}$.
Summing over compositions of $M$ proves the coefficient formula \eqref{t3:eq:h-coeff}.

Here is a simultaneous justification of all rearrangements. Since
$\pp{m_i}\ge\Pinf$, and there are $\binom{M+r-1}{r-1}$ weak compositions of $M$,
\begin{equation}\label{t3:eq:H-bound}
 |h_{\eps,M}^{(r)}|
 \le\frac{\max_{0\le j<r}\T_j^{(r)}}{\Pinf^r}
 \binom{M+r-1}{r-1}Q^{M^2/(2r)+M/2}.
\end{equation}
The sum of the right-hand side times $R^M$ converges for every finite $R$. Applying the same completion of squares to the absolute values of the expanded terms proves absolute convergence of the full multiple series, uniformly for $|t|\le R$. Tonelli's theorem and uniform convergence now validate the expansions, the rearrangements, and the entire character of $H$. This also justifies the zero-extension argument at integer $n$.
\end{proof}

\section{Forward coefficients, logarithmic coefficients, and the first terms}

Normalize by the leading theta constant. From \eqref{t3:eq:Galois-exact},
\begin{equation}\label{t3:eq:Galois-forward-normalized}
 G_n^{(r)}(q)=\frac{\T_{\eps}^{(r)}}{\Pinf^{r-1}}
 q^{(r-1)n^2/(2r)}\sum_{m\ge0}v_{\eps,m}^{(r)}Q^{mn/r},
\end{equation}
where the series converges and
\begin{equation}\label{t3:eq:Galois-v}
 v_{\eps,m}^{(r)}=
 \frac1{\T_{\eps}^{(r)}}
 \sum_{j=0}^{\lfloor m/r\rfloor}
 \frac{Q^j}{\pp{j}}\,h_{\eps,m-rj}^{(r)}.
\end{equation}
This is a finite formula using only $r$ theta constants and finite composition sums.

For inversion, define
\begin{equation}\label{t3:eq:Galois-B}
 B_{\eps}^{(r)}(t)=\A(t^r)+
       \log\frac{H_{\eps}^{(r)}(t)}{\T_{\eps}^{(r)}}.
\end{equation}
Since $H(0)=\T_{\eps}^{(r)}>0$, the logarithm is analytic in a nonzero disk about zero. If
$c_m=h_{\eps,m}^{(r)}/\T_{\eps}^{(r)}$ for $m\ge1$, then its coefficients are
\begin{equation}\label{t3:eq:Galois-b}
 b_m=\one_{r\mid m}\alpha_{m/r}
 +\sum_{k=1}^m\frac{(-1)^{k+1}}k
  \sum_{\substack{j_1,\ldots,j_k\ge1\\\sum j_i=m}}
       c_{j_1}\cdots c_{j_k}.
\end{equation}
Thus the entire-function representation is not hiding any recurrence in the logarithmic or inverse coefficients.

The first two normalized entire coefficients are
\begin{align}
 c_1&=-\frac{rQ^{(r+1)/(2r)}}{1-Q}
               \frac{\T_{\eps+1}^{(r)}}{\T_{\eps}^{(r)}},
 \label{t3:eq:Galois-c1}\\
 c_2&=Q^{1+2/r}\frac{\T_{\eps+2}^{(r)}}{\T_{\eps}^{(r)}}
 \left(\frac r{(1-Q)(1-Q^2)}
       +\frac{\binom r2}{(1-Q)^2}\right).
 \label{t3:eq:Galois-c2}
\end{align}
Consequently
\begin{equation}\label{t3:eq:Galois-b12}
 b_1=c_1,\qquad b_2=c_2-\frac{c_1^2}2+\one_{r=2}\alpha_1.
\end{equation}
The special contribution at $r=2$ comes from the numerator product $\exp\A(t^r)$.

For ordinary Galois numbers, \eqref{t3:eq:h-coeff} becomes the particularly compact formula
\begin{equation}\label{t3:eq:H-r2}
 h_{\eps,M}^{(2)}=(-1)^M Q^{M(M+2)/4}\T_{\eps+M}(Q)
     \sum_{j=0}^M\frac1{\pp{j}\pp{M-j}}.
\end{equation}
All higher exponentially small corrections to \eqref{t3:eq:Galois-leading} are therefore explicit.

\section{Inverse theta transseries}

Define the residue-sheet interpolation by replacing $n$ by a real $x$ on the right-hand side of \eqref{t3:eq:Galois-exact}, with $\eps$ held fixed. It is positive and increasing for all sufficiently large $x$. Its phase data are
\begin{equation}\label{t3:eq:Galois-phase}
 a=\frac{(r-1)h}{2r},\quad b=0,\quad
 c=\log\frac{\T_{\eps}^{(r)}}{\Pinf^{r-1}},\quad
 \lambda=\frac hr.
\end{equation}
In particular,
\begin{equation}\label{t3:eq:Galois-center}
 s_{\eps}=\sqrt{\frac{2r}{(r-1)h}
       \log\left(\frac{Y\Pinf^{r-1}}{\T_{\eps}^{(r)}}\right)},
 \qquad D_{\eps}=\frac{(r-1)h}{r}s_{\eps}.
\end{equation}
The complete inverse is
\begin{equation}\label{t3:eq:Galois-inverse-full}
 \boxed{\displaystyle
 (G_{q,\eps}^{(r)})^{-1}(Y)
 =s_{\eps}+\sum_{m\ge1}d_{\eps,m}(D_{\eps})Q^{ms_{\eps}/r},}
\end{equation}
with \eqref{t3:eq:Galois-b} substituted into \eqref{t3:eq:finite-dm}. The expansion converges for all sufficiently large $Y$.

Its first term is
\begin{align}\label{t3:eq:Galois-first}
 (G_{q,\eps}^{(r)})^{-1}(Y)=s_{\eps}
 &+\frac{rQ^{(r+1)/(2r)}}{(1-Q)D_{\eps}}
     \frac{\T_{\eps+1}^{(r)}}{\T_{\eps}^{(r)}}
      Q^{s_{\eps}/r}
 +O(Q^{2s_{\eps}/r}/D_{\eps}).
\end{align}
This correction is positive. The next coefficient is obtained from \eqref{t3:eq:d2} and \eqref{t3:eq:Galois-b12}, and every further coefficient is still a finite expression.

At $Q=1/2$ the constants are numerically
\begin{align*}
 \T_0^{(2)}&=2.1289368272118771586694585485\ldots,\\
 \T_1^{(2)}&=2.1289312505130275585916134026\ldots,\\
 \T_0^{(3)}&=5.2335188744705752675068634419\ldots,\\
 \T_1^{(3)}=\T_2^{(3)}&=5.2335186066094946295319364209\ldots.
\end{align*}
Their small numerical differences must not be discarded in an all-order inverse calculation. For fixed $Q$, a nonzero change of $c$ shifts the center by order $1/s$, which eventually dominates every exponential sector.

\section{Weighted flags and Rogers--Szeg\H{o} polynomials}
\label{t3:sec:weighted-Galois}

A multivariate extension is almost automatic. Fix positive real weights $z_1,\ldots,z_r$ and define
\begin{equation}\label{t3:eq:weightedGalois}
 G_n^{(r)}(q;\boldsymbol z)=
 \sum_{k_1+\cdots+k_r=n}
 \genfrac{[}{]}{0pt}{}{n}{k_1,\ldots,k_r}_q
 \prod_i z_i^{k_i}.
\end{equation}
Replace the theta constant by
\begin{equation}\label{t3:eq:weightedtheta}
 \T_{\eps}^{(r)}(Q;\boldsymbol z)=
 \sum_{\boldsymbol j\in\mathcal L_{\eps}^{(r)}}
 Q^{\|\boldsymbol j\|^2/2}\prod_i z_i^{j_i}.
\end{equation}
The Gaussian factor guarantees convergence for every fixed positive weight vector. The proof of Theorem~\ref{t3:thm:theta-tail}, including its absolute-convergence argument, gives the entire coefficients
\begin{align}\label{t3:eq:weightedh}
 h_{\eps,M}^{(r)}(\boldsymbol z)
 ={}&(-1)^M Q^{M^2/(2r)+M/2}
       \T_{\eps+M}^{(r)}(Q;\boldsymbol z)\\[-0.3em]
 &\times\sum_{m_1+\cdots+m_r=M}
       \prod_{i=1}^r\frac{z_i^{M/r-m_i}}{\pp{m_i}}.\notag
\end{align}
The extra factor is exactly what results from translating the linear weight under the completion of squares \eqref{t3:eq:theta-complete}. The leading prefactor acquires
$\prod_i z_i^{n/r}$, and the inverse phase has
\begin{equation}\label{t3:eq:weightedphase}
 a=\frac{(r-1)h}{2r},\qquad
 b=\frac1r\sum_i\log z_i,\qquad
 c=\log\frac{\T_{\eps}^{(r)}(Q;\boldsymbol z)}{\Pinf^{r-1}}.
\end{equation}
Nothing else changes in the inversion theorem.

At $r=2$, $z_1=z$, $z_2=1$, the finite sum is the Rogers--Szeg\H{o} polynomial
$\sum_k\qbinq{n}{k}z^k$. The relation between generalized Galois numbers, flags, and generalized Rogers--Szeg\H{o} polynomials is also studied in \cite{bliem}. Here fixed positive weights preserve an explicit convergent inverse transseries with the same small parameter $Q^{s/r}$.
% ---- hand chunk 08-double-head.tex ----

\chapter{The singular transition $q\to1$}
\label{ct:ch:double}

% ---- T3 lines 1502--1671 ----
\label{t3:sec:double}

\section{Why a fixed-\texorpdfstring{$q$}{q} transseries cannot simply be specialized}

The coefficients $\alpha_m=Q^m/[m(1-Q^m)]$ grow like $1/(hm^2)$ as $h\to0$. Simultaneously $\Pinf$ tends to zero exponentially in $1/h$. A limit that lets $q\to1$ while $x\to\infty$ therefore reorganizes both the nominal constant and the finite-product tail.

Consider the scaling
\begin{equation}\label{t3:eq:double-scaling}
 Q=\ee^{-h},\qquad h\downarrow0,\qquad x=\frac\tau h,
 \qquad \tau>0.
\end{equation}
The analytic parameter $q$ is continuous here; there is no family of finite fields whose orders tend to one. The finite-field interpretation motivates the formulas, but this limiting problem concerns their real-$q$ special-function extensions.

\section{An all-order dilogarithmic expansion of the tail}

Write $z=\ee^{-\tau}$. The exact tail is
\begin{equation}\label{t3:eq:double-Aexact}
 \mathcal A_{\ee^{-h}}(z)
 =\sum_{m\ge1}\frac{\ee^{-m\tau}}{m(\ee^{hm}-1)}.
\end{equation}
Let the Bernoulli numbers be defined by
$v/(\ee^v-1)=\sum_{j\ge0}B_jv^j/j!$ near zero.

\begin{proposition}[Uniform finite-order tail expansion]\label{t3:prop:double-A}
For every fixed integer $K\ge1$ and every $\tau_0>0$, uniformly for $\tau\ge\tau_0$ as $h\downarrow0$,
\begin{align}\label{t3:eq:double-A}
 \mathcal A_{\ee^{-h}}(\ee^{-\tau})
 ={}&\frac{\Li_2(\ee^{-\tau})}{h}
     +\frac12\log(1-\ee^{-\tau})\\[-0.3em]
 &+\sum_{k=1}^K\frac{B_{2k}}{(2k)!}h^{2k-1}
          \Li_{2-2k}(\ee^{-\tau})
 +O_{K,\tau_0}(h^{2K+1}).\notag
\end{align}
The assertion remains valid after any fixed number of derivatives with respect to $\tau$, with a corresponding change of the implied constant.
\end{proposition}
\begin{proof}
The remainder in the local expansion
\[
 \frac1{\ee^v-1}=\frac1v-\frac12+
   \sum_{k=1}^K\frac{B_{2k}}{(2k)!}v^{2k-1}+r_K(v)
\]
satisfies $|r_K(v)|\le C_Kv^{2K+1}$ for all real $v>0$. Near zero this follows from analyticity after removal of $1/v$. On a compact positive interval it follows from continuity, and for large $v$ the displayed polynomial and $1/v$ grow no faster than a constant times $v^{2K+1}$.

Insert $v=hm$ into \eqref{t3:eq:double-Aexact}. The sum of the absolute remainders is at most
\[
 C_Kh^{2K+1}\sum_{m\ge1}\ee^{-m\tau}m^{2K},
\]
which is uniformly bounded by a constant times $h^{2K+1}$ for $\tau\ge\tau_0$. Summing the explicit powers of $m$ gives the polylogarithms in \eqref{t3:eq:double-A}. Differentiation in $\tau$ adds only a fixed power of $m$, leaving the same argument valid.
\end{proof}

\section{The modular factor and its separate exponential scale}

The modular transformation of the Euler product gives the exact identity \cite{dlmfqmod}
\begin{equation}\label{t3:eq:modularP}
 \mathcal P_{\ee^{-h}}
 =\sqrt{\frac{2\pi}{h}}
   \exp\left(-\frac{\pi^2}{6h}+\frac h{24}\right)
   \mathcal P_{\widehat Q},
 \qquad \widehat Q=\ee^{-4\pi^2/h}.
\end{equation}
For example, it follows from the Dedekind eta transformation
$\eta(-1/\tau)=\sqrt{-i\tau}\,\eta(\tau)$ at $\tau=ih/(2\pi)$, together with
$\eta(ih/(2\pi))=\ee^{-h/24}\mathcal P_{\ee^{-h}}$.
The exponentially small product left over is explicit:
\begin{equation}\label{t3:eq:modularsectors}
 -\log\mathcal P_{\widehat Q}
 =\sum_{N\ge1}\sigma_{-1}(N)\widehat Q^N,
 \qquad\sigma_{-1}(N)=\sum_{d\mid N}\frac1d.
\end{equation}
To prove \eqref{t3:eq:modularsectors}, expand each logarithm of the Euler product and group the terms with the same product of indices. This series converges for $0<\widehat Q<1$.

The new variable $\widehat Q$ is not a power of $Q^x=\ee^{-\tau}$. The simultaneous limit therefore contains two different mechanisms: a finite-size tail described by \eqref{t3:eq:double-Aexact}, and the modular exponential factor \eqref{t3:eq:modularsectors}. The latter is an exact separated contribution. It is not being asserted to describe the entire beyond-all-orders structure of the finite-size tail's Bernoulli expansion.

\section{Central Gaussian coefficients and the crossover phase}

For the central Gaussian interpolation,
\begin{equation}\label{t3:eq:central-double-exact}
 \log H_{1,1}(\tau/h)
 =\frac{\tau^2}{h}-\log\mathcal P_{\ee^{-h}}
 +\mathcal A_{\ee^{-h}}(\ee^{-2\tau})
 -2\mathcal A_{\ee^{-h}}(\ee^{-\tau}).
\end{equation}
Introduce
\begin{align}
 S(\tau)&=\tau^2+\frac{\pi^2}{6}
               +\Li_2(\ee^{-2\tau})-2\Li_2(\ee^{-\tau}),
 \label{t3:eq:Sphase}\\
 J(\tau)&=\frac12\log(1-\ee^{-2\tau})
                  -\log(1-\ee^{-\tau}),\label{t3:eq:Jamp}\\
 E_1(\tau)&=-\frac1{24}
       +\frac{\Li_0(\ee^{-2\tau})-2\Li_0(\ee^{-\tau})}{12},
 \label{t3:eq:E1}\\
 E_k(\tau)&=\frac{B_{2k}}{(2k)!}
       \left(\Li_{2-2k}(\ee^{-2\tau})
                 -2\Li_{2-2k}(\ee^{-\tau})\right),\quad k\ge2.
 \label{t3:eq:Ek}
\end{align}
Combining \eqref{t3:eq:double-A} and \eqref{t3:eq:modularP} gives
\begin{align}\label{t3:eq:central-double-asympt}
 \log H_{1,1}(\tau/h)
 ={}&\frac{S(\tau)}h-\frac12\log\frac{2\pi}{h}+J(\tau)
     +\sum_{k=1}^K h^{2k-1}E_k(\tau)\\[-0.3em]
 &-\log\mathcal P_{\widehat Q}
     +O_{K,\tau_0}(h^{2K+1}).\notag
\end{align}
The modular term is retained exactly, although a finite algebraic remainder is of course much larger than that term. Equation \eqref{t3:eq:central-double-exact} remains available when the finite-size tail itself must be evaluated beyond that algebraic accuracy.

The derivative of the new dominant phase has a useful elementary form:
\begin{equation}\label{t3:eq:Sderivative}
 \boxed{\displaystyle
 S'(\tau)=2\log(1+\ee^{\tau})>0,\qquad
 S''(\tau)=\frac{2\ee^\tau}{1+\ee^\tau}>0.}
\end{equation}
Indeed, differentiating the dilogarithms gives
$2\tau+2\log(1-\ee^{-2\tau})-2\log(1-\ee^{-\tau})$
which reduces to \eqref{t3:eq:Sderivative}. Also $S(0)=0$ by continuity. Thus $S$ is increasing from zero to infinity and has an unambiguous positive inverse.

For orientation, its endpoint expansions are
\begin{align}
 S(\tau)&=2(\log2)\tau+\frac{\tau^2}{2}
                  +\frac{\tau^3}{12}+O(\tau^5),\quad\tau\to0,
 \label{t3:eq:Ssmall}\\
 S(\tau)&=\tau^2+\frac{\pi^2}{6}+O(\ee^{-\tau}),\quad\tau\to\infty.
 \label{t3:eq:Slarge}
\end{align}
They follow from \eqref{t3:eq:Sderivative} and the defining dilogarithms. The small-$\tau$ phase and $J(\tau)\sim\tfrac12\log(2/\tau)$ display the classical central-binomial phase and amplitude. The uniform remainder theorem above is stated for $\tau\ge\tau_0>0$; matching it uniformly all the way to $\tau=0$ requires a different error analysis.

\section{Inverting the crossover phase}

Let $\ell=h\log Y$ be fixed and positive, and define
\begin{equation}\label{t3:eq:crossover-center}
 \tau_*=S^{-1}(\ell),\qquad
 L_h=\frac12\log\frac{2\pi}{h}.
\end{equation}
On compact positive ranges of $\ell$, the derivative in \eqref{t3:eq:Sderivative} is bounded away from zero. The inverse index therefore has the expansion
\begin{equation}\label{t3:eq:crossover-inverse}
 x=\frac{\tau_*}{h}+c_1(\tau_*,L_h)
                         +h c_2(\tau_*,L_h)
                         +O\bigl(h^2(1+L_h)^3\bigr),
\end{equation}
where
\begin{align}
 c_1&=\frac{L_h-J(\tau_*)}{S'(\tau_*)},\label{t3:eq:crossover-c1}\\
 c_2&=-\frac{\tfrac12 S''(\tau_*)c_1^2
                  +J'(\tau_*)c_1+E_1(\tau_*)}{S'(\tau_*)}.
 \label{t3:eq:crossover-c2}
\end{align}
To prove these formulas, write $\tau=\tau_*+hc_1+h^2c_2+\cdots$, multiply \eqref{t3:eq:central-double-asympt} by $h$, and compare coefficients in
\[
 S(\tau)+h(J(\tau)-L_h)+h^2E_1(\tau)+\cdots=\ell.
\]
The error estimate follows from Proposition~\ref{t3:prop:double-A}, Taylor's theorem, and the positive derivative. The modular term is smaller than every fixed power of $h$ and so does not affect this finite algebraic estimate.

There is a finite coefficient formula at every higher order. Treat $L_h$ as an independent symbol during coefficient extraction, and put
\begin{equation}\label{t3:eq:crossoverR}
 \mathcal R(h,\tau)=h(J(\tau)-L_h)
                  +\sum_{j\ge1}h^{2j}E_j(\tau).
\end{equation}
The formal series $\tau=\tau_*+\sum_{m\ge1}c_mh^m$ has
\begin{equation}\label{t3:eq:crossover-all}
 c_m=\sum_{k=1}^m\frac{(-1)^k}{k!}
 \left[
 \left(\frac1{S'(\tau)}\frac{\ddn}{\ddn\tau}\right)^{k-1}
 \frac{[h^m]\mathcal R(h,\tau)^k}{S'(\tau)}
 \right]_{\tau=\tau_*}.
\end{equation}
The proof is the same inverse-perturbation identity used for
\eqref{t3:eq:general-derivative}. At each fixed $m$, only finitely many $E_j$ occur. Each $c_m$ is a polynomial in $L_h$ of degree at most $m$, with smooth coefficients in $\tau_*$. The finite-order expansions are asymptotic with remainders controlled by the preceding propositions; no convergence assertion is made for this Bernoulli-based expansion in $h$.

The main lesson is quantitative. At fixed positive $\ell=h\log Y$, the correct leading center is $S^{-1}(\ell)/h$, not the fixed-$q$ square-root center with its finite-size tail dropped. The tail contributes at order $1/h$ in this regime and therefore belongs to the dominant phase itself.
% ---- hand chunk 085-part8-head.tex ----

\part{Slow growth and finite limits}
\label{ct:part:slow}

\chapter{Harmonic functions: exponential and endpoint-Puiseux inversion}
\label{ct:ch:harmonic}

The families so far grow at least exponentially.  Slow growth reverses the
scales: an inverse-power expansion of the harmonic function in a centred
variable becomes an \emph{exponential} inverse series, and a generalized
harmonic function that converges to $\zeta(s)$ has an inverse that is a
generalized Puiseux series in the distance to that endpoint.

% ---- T2 lines 1678--1790 ----
\label{t2:sec:harmonic}

\section{A slowly increasing combinatorial function}

Harmonic numbers occur throughout permutation enumeration, including the $k=2$ fixed-cycle Stirling formula. Their standard real continuation is
\begin{equation}
 H(x)=\psi(x+1)+\gamma,\qquad x\ge0,
 \label{t2:eq:harmonic-cont}
\end{equation}
which agrees with $H_n=\sum_{j=1}^n1/j$ at integers. The digamma and generalized-harmonic conventions agree with \texttt{HarmonicNumber} \cite{wlharmonic}. Since
$H'(x)=\sum_{j\ge0}(x+1+j)^{-2}>0$, this function has a unique increasing inverse on $[0,\infty)$.

The centered variable $w=x+1/2$ eliminates odd inverse powers from the logarithmic correction. The shifted digamma expansion gives
\begin{equation}
 H(x)-\gamma\sim\log w+C(w^{-2}),\qquad
 C(z)=-\sum_{m\ge1}\frac{B_{2m}(1/2)}{2m}z^m.
 \label{t2:eq:harmonic-centered}
\end{equation}
For example, $C(z)=z/24-7z^2/960+\cdots$. This follows by differentiating the shifted gamma expansion, or by Euler--Maclaurin summation with the half-integer shift \cite{dlmfgamma,dlmfbernoulli}.

\begin{theorem}[All coefficients of the harmonic inverse]
\label{t2:thm:harmonic-inverse}
Let $X=\ee^{y-\gamma}$. Then
\begin{equation}
 H^{-1}(y)\sim X-\frac12+\sum_{m\ge1}h_mX^{1-2m},
 \qquad
 h_m=-\frac1{2m-1}[z^m]\exp\bigl((2m-1)C(z)\bigr).
 \label{t2:eq:harmonic-inverse}
\end{equation}
The first five coefficients are
\begin{align}
 h_1&=-\frac1{24},&h_2&=\frac3{640},&
 h_3&=-\frac{1525}{580608},\nonumber\\
 h_4&=\frac{615881}{199065600},&
 h_5&=-\frac{3058641}{504627200}.
 \label{t2:eq:harmonic-coeff}
\end{align}
A truncation through $h_MX^{1-2M}$ has error $\bigO(X^{-2M-1})$.
\end{theorem}
\begin{proof}
Put $v=1/w$ and $t=1/X$. Equation~\eqref{t2:eq:harmonic-centered} becomes formally
\[
 v=t\exp(C(v^2)).
\]
The Laurent form of Lagrange inversion gives, for $m\ge1$,
\[
 [t^{2m-1}]\frac1{v(t)}
 =-\frac1{2m-1}[u^{2m}]
             \exp\bigl((2m-1)C(u^2)\bigr),
\]
which is the coefficient in \eqref{t2:eq:harmonic-inverse}. One may derive this Laurent version by applying the usual residue proof of Lagrange inversion to $v^{-1}$; the leading $t^{-1}$ term is separated before extracting positive powers. Oddness of $v(t)$ forces the displayed parity of powers in $w=1/v$.

Substitution of a finite truncation into a sufficiently long centered digamma expansion gives a forward residual of order $X^{-2M-2}$. Since $H'(x)\asymp X^{-1}$ near the proposed inverse, Lemma~\ref{t1:thm:certificate} gives an inverse error $\bigO(X^{-2M-1})$. Formula~\eqref{t2:eq:E} computes the displayed rational coefficients exactly.
\end{proof}

In the target variable, the correction sequence is
\[
 \ee^{y-\gamma},\quad 1,\quad
 \ee^{-(y-\gamma)},\quad\ee^{-3(y-\gamma)},\quad\ldots.
\]
Thus even an inverse-power expansion in the centered forward variable becomes an exponential inverse series. The Bernoulli coefficient series remains generally divergent; no convergence claim is made by writing all its coefficients.

\section{Generalized harmonic functions near their finite limit}

Fix a real $s>1$ and define
\begin{equation}
 H_s(x)=\zeta(s)-\zeta(s,x+1),\qquad x\ge0.
 \label{t2:eq:general-harmonic}
\end{equation}
The Hurwitz-zeta series shows
$H_s'(x)=s\zeta(s+1,x+1)>0$. Its range approaches the finite endpoint $\zeta(s)$, so the inverse asymptotic problem is now $y\uparrow\zeta(s)$, not $y\to\infty$.

Put
\begin{equation}
 p=s-1,\qquad \delta=\zeta(s)-y,\qquad
 T=(p\delta)^{-1/p},\qquad w=x+\frac12.
 \label{t2:eq:general-harmonic-core}
\end{equation}
The centered Hurwitz-zeta expansion \cite{dlmfhurwitz} gives
\begin{align}
 \zeta(s,x+1)&\sim\frac{w^{-p}}pD_s(w^{-2}),\nonumber\\
 D_s(z)&=1+\sum_{m\ge1}
  \frac{pB_{2m}(1/2)\rising{s}{2m-1}}{(2m)!}z^m.
 \label{t2:eq:general-harmonic-D}
\end{align}
In particular $D_s(z)=1-ps\,z/24+\cdots$.

\begin{theorem}[Endpoint inverse coefficients]
\label{t2:thm:general-harmonic-inverse}
For fixed $s>1$,
\begin{equation}
 H_s^{-1}(y)\sim T-\frac12+\sum_{m\ge1}h_{s,m}T^{1-2m},
 \quad
 h_{s,m}=-\frac1{2m-1}[z^m]
                    D_s(z)^{-(2m-1)/p}.
 \label{t2:eq:general-harmonic-inverse}
\end{equation}
In particular,
\begin{equation}
 H_s^{-1}(y)=T-\frac12-\frac{s}{24T}+\bigO(T^{-3}).
 \label{t2:eq:general-harmonic-first}
\end{equation}
The remainder after order $M$ is $\bigO(T^{-2M-1})$.
\end{theorem}
\begin{proof}
The forward equation reads $T=wD_s(w^{-2})^{-1/p}$. Thus with $v=1/w$ and $t=1/T$,
\[
 v=tD_s(v^2)^{-1/p}.
\]
The same Laurent Lagrange formula used for the harmonic inverse gives \eqref{t2:eq:general-harmonic-inverse}. The first coefficient is $-s/24$ by direct expansion. For a finite truncation, the relative forward error in $\delta$ is $\bigO(T^{-2M-2})$, hence the absolute error is $\bigO(T^{-p-2M-2})$. The derivative magnitude of the tail is $\asymp T^{-p-1}$, so the inverse error is $\bigO(T^{-2M-1})$.
\end{proof}

Because $T$ is a fractional power of the distance $\delta$ to the endpoint, this is a Puiseux-type expansion in $\delta^{1/p}$. When $p$ is not rational, the powers need not be rational; ``generalized Puiseux'' is the appropriate description. The finite coefficient formula still makes sense for every fixed real $s>1$.
% ---- hand chunk 085-endpoint-head.tex ----

\chapter{Inverting convergence to a finite limit}
\label{ct:ch:endpoint}

The normalized $q$-products of \cref{ct:part:qproducts} converge to positive
constants, and their inverse problem asks how large an index is needed to
reach a small prescribed deficit.  The correct leading inverse is
logarithmic, and the convergent tail makes every correction explicit.

% ---- T3 lines 1141--1248 ----
\label{t3:sec:endpoint}

The previous examples grow without bound, but their normalized versions often converge to positive constants. Their inverse problem asks how large an index is needed to reach a small prescribed deficit. The correct leading inverse is logarithmic.

\section{A general endpoint inversion formula}

\begin{theorem}[Logarithmic endpoint inverse]\label{t3:thm:endpoint}
Let $B(t)=b_1t+b_2t^2+\cdots$ be analytic at zero with $b_1>0$, and set
$\eps=B(\ee^{-\lambda x})$. For positive sufficiently small $\eps$, the large real solution has a convergent expansion
\begin{equation}\label{t3:eq:endpoint-all}
 \boxed{\displaystyle
 x=\frac1\lambda\log\frac{b_1}{\eps}
 -\frac1\lambda\sum_{m\ge1}\frac{\eps^m}{m}
     \coefB{u}{m}\left(\frac{u}{B(u)}\right)^m.}
\end{equation}
In particular,
\begin{equation}\label{t3:eq:endpoint-first}
 x=\frac1\lambda\log\frac{b_1}{\eps}
  +\frac{b_2}{\lambda b_1^2}\eps
  +\frac1\lambda\left(\frac{b_3}{b_1^3}
             -\frac{3b_2^2}{2b_1^4}\right)\eps^2
  +O(\eps^3).
\end{equation}
\end{theorem}
\begin{proof}
The analytic inverse theorem gives $t=t(\eps)$ with $B(t)=\eps$ and
$t=\eps/b_1+O(\eps^2)$. Let $R(t)=t/B(t)$, analytic at zero. Since
$t/\eps=R(t)$, we have
$x=-\lambda^{-1}\log\eps-\lambda^{-1}\log R(t(\eps))$.
Lagrange--B\"urmann gives, for $m\ge1$,
\[
 [\eps^m]\log R(t(\eps))
 =\frac1m[u^{m-1}]\frac{R'(u)}{R(u)}R(u)^m
 =\frac1m[u^m]R(u)^m.
\]
The constant term is $\log R(0)=-\log b_1$. This proves \eqref{t3:eq:endpoint-all}; expanding the first two coefficients proves \eqref{t3:eq:endpoint-first}. All series used are analytic near zero, establishing convergence.
\end{proof}

\section{Invertibility probabilities of random matrices}

The probability that a uniformly random $n\times n$ matrix over $\F_q$ is invertible is
\begin{equation}\label{t3:eq:probability}
 p_n(q)=\frac{|\GL(n,q)|}{q^{n^2}}=\pp{n}.
\end{equation}
It decreases to $\Pinf$. With the canonical interpolation and the logarithmic relative excess
\begin{equation}\label{t3:eq:probability-eps}
 \eps=\log\frac{\mathcal P_Q(x)}{\Pinf},
\end{equation}
we have exactly $\eps=\A(Q^x)$. Therefore
\begin{equation}\label{t3:eq:probability-inverse}
 x=\frac1h\log\frac{\alpha_1}{\eps}
 +\frac{\alpha_2}{h\alpha_1^2}\eps
 +\frac1h\left(\frac{\alpha_3}{\alpha_1^3}
       -\frac{3\alpha_2^2}{2\alpha_1^4}\right)\eps^2
 +O(\eps^3).
\end{equation}
At $q=2$, this becomes
\begin{equation}\label{t3:eq:probability-q2}
 x=\frac1{\log2}\left(
       \log\frac1\eps+\frac\eps6+\frac{\eps^2}{168}
       +O(\eps^3)\right).
\end{equation}
For a relative excess $\Delta=\mathcal P_Q(x)/\Pinf-1$, substitute
$\eps=\log(1+\Delta)$, not $\eps=\Delta$. Confusing these two small variables changes every coefficient after the leading approximation.

Since the interpolation decreases, the least integer $n$ satisfying
$p_n(q)\le\Pinf\ee^{\eps}$ is the ceiling of this real inverse, provided it lies in the nonnegative range. Strict inequalities require the corresponding strict-rounding convention.

\section{Normalized central Gaussian coefficients}

Define
\begin{equation}\label{t3:eq:central-normalized}
 R_q(x)=\Pinf q^{-x^2}H_{1,1}(x)
       =\exp\bigl(\A(t^2)-2\A(t)\bigr),\qquad t=Q^x.
\end{equation}
For sufficiently large $x$, this approaches one from below. Indeed,
\[
 -\log R_q(x)=2\A(t)-\A(t^2)=2\alpha_1t+O(t^2)>0.
\]
Its derivative with respect to $t$ is positive near zero, so $R_q$ increases to one there. Theorem~\ref{t3:thm:endpoint} applies to
$\eps=-\log R_q(x)$ with
\begin{equation}\label{t3:eq:central-endpoint-b}
 b_m=2\alpha_m-\one_{2\mid m}\alpha_{m/2}.
\end{equation}
It gives an explicit logarithmic inverse for attaining a prescribed normalized accuracy, independently of the growing-count inverse in Section~\ref{t3:sec:gaussian}.

\section{Higher-order leading tails and Puiseux corrections}

Suppose instead
\[
 B(t)=b_rt^r(1+O(t)),\qquad b_r>0.
\]
Set $\eta=(\eps/b_r)^{1/r}$ and choose the analytic root
\[
 C(t)=\left(B(t)/b_r\right)^{1/r}=t(1+O(t)).
\]
The ordinary analytic inverse of $C$ gives $t=\eta(1+O(\eta))$. Hence
\begin{equation}\label{t3:eq:puiseux-endpoint}
 x=\frac1{r\lambda}\log\frac{b_r}{\eps}
       +\sum_{m\ge1}\gamma_m\eps^{m/r},
\end{equation}
with coefficients obtained by the same Lagrange formula applied to $C$. The series in $\eps^{1/r}$ converges near zero on the chosen positive branch.

There is an important simplification when $B(t)=\widehat B(t^r)$: put $u=t^r$ and invert $\eps=\widehat B(u)$ directly. Only integer powers of $\eps$ remain after the logarithm. In particular the normalized symplectic count has this simpler structure, because its tail is $\mathcal A_{Q^2}(Q^{2x})$. Its leading index is
\[
 x\sim\frac1{2h}\log\frac{Q^2}{(1-Q^2)\eps},
\]
followed by a convergent series in integer powers of $\eps$.
% ---- hand chunk 09-part9-head.tex ----

\part{Certification, computation, and synthesis}
\label{ct:part:certification}

\chapter{Error certification, threshold decisions, and computation}
\label{ct:ch:certification}

The certificate of \cref{ct:sec:residual} is the bridge from a symbolic
inverse to a number and from a number to an integer index.  This chapter
records how it is applied to the families of the volume, which analytic
tail bounds feed it, what a safe ceiling criterion is, and how the
coefficients are computed without symbolic explosion.

% ---- T3 lines 1674--1719 ----
\label{t3:sec:certification}

\section{Analytic error bounds versus numerical evidence}

Theorem~\ref{t3:thm:inverse} supplies an analytic, geometrically decreasing outer-sector remainder. For a concrete parameter value, it is often more economical to certify an approximation by its residual. Let
\begin{equation}\label{t3:eq:residual}
 \mathcal E(x)=P(x)+B(\ee^{-\lambda x})-\log Y.
\end{equation}
Suppose a known interval $[u,v]$ contains the desired root and
$\mathcal E'(x)\ge\mu>0$ throughout it. Then for any $x_M\in[u,v]$,
\begin{equation}\label{t3:eq:residualbound}
 |x_M-F^{-1}(Y)|\le\frac{|\mathcal E(x_M)|}{\mu}.
\end{equation}
This is the mean value theorem. Alternatively, checking opposite signs at
$x_M-\eta$ and $x_M+\eta$, together with monotonicity, directly certifies a root enclosure of radius $\eta$.

For $q$-product tails, Lemma~\ref{t3:lem:Atail} bounds the omitted logarithmic terms using only elementary quantities. Derivatives can be bounded in the same way after differentiating the absolutely convergent series. For a theta-tail function, the majorant \eqref{t3:eq:H-bound} bounds the omitted entire coefficients; a fixed polynomial growth factor is dominated by the Gaussian decrease in their weight.

The supplied program uses high-precision floating-point arithmetic for its analytic tests. Those tests are reproducibility checks, not interval-arithmetic certificates. The proofs and inequalities above establish the analytical assertions independently of floating-point evidence. For a certified implementation one would evaluate the residual, derivative bounds, and product or theta remainders with outward-rounded interval arithmetic.

\section{A safe ceiling criterion}

If a real inverse has a certified enclosure $[u,v]$ and
$\lceil u\rceil=\lceil v\rceil$, its integer threshold is determined. If the enclosure crosses an integer, evaluate the original exact sequence at that candidate integer. The same rule applies to a residue-class lattice, using \eqref{t3:eq:residue-rounding}.

It is essential to retain equality cases. For a threshold defined using $a_n\ge Y$, the correct rounding is a ceiling of the real inverse. For a condition $a_n>Y$, an exact integer inverse instead requires the next lattice point. A floating-point value such as $n-10^{-50}$ can therefore produce the wrong answer even when it appears extraordinarily accurate.

For primitive necklaces, Theorem~\ref{t3:thm:two-candidate} provides a particularly simple implementation: bracket the leading solution using exact comparisons of $q^n/n$ with $Y$, then perform one exact comparison of $I_q(n_0)$ with $Y$. This avoids trying to choose one fixed-radical interpolation before the integer index is known.

\section{Computing coefficients without symbolic explosion}

The finite formulas make arbitrary weights well defined, but a direct symbolic expansion of all compositions is not always the fastest implementation. The following operations suffice for a stable truncated-power-series implementation:
\begin{equation}\label{t3:eq:seriesexpalgorithm}
 f(t)=\ee^{g(t)}\quad\Longrightarrow\quad
 f_0=\ee^{g_0},\qquad
 f_n=\frac1n\sum_{j=1}^n jg_jf_{n-j},
\end{equation}
and, if $f_0\ne0$,
\begin{equation}\label{t3:eq:serieslogalgorithm}
 g(t)=\log f(t)\quad\Longrightarrow\quad
 g_n=\frac{nf_n-\sum_{j=1}^{n-1}jg_jf_{n-j}}{nf_0}.
\end{equation}
Both follow by comparing coefficients in $f'=g'f$. Substituting these into
\eqref{t3:eq:dm-recurrence} computes the inverse through any chosen finite weight without differentiating a large expression. The code also checks the independent nonrecursive formula \eqref{t3:eq:finite-dm} against this recurrence through weight eight.

For generalized Galois numbers, compute the finite products $\pp{j}$ once, form the composition coefficients $A_{r,M}$ by truncated convolution, multiply by the appropriate theta constant in \eqref{t3:eq:h-coeff}, take a truncated logarithm, and apply the inverse-coefficient routine. Only $r$ theta constants are required for any number of sectors at fixed rank $r$.
% ---- T1 lines 1285--1288 ----
\section{From inverse enclosures to exact integer indices}
For any eventual monotone interpolation, \eqref{t2:eq:ceil} converts the real inverse to the threshold index. A certified interval $[u,v]$ for the inverse determines that index immediately when $\lceil u\rceil=\lceil v\rceil$. If the interval straddles an integer, evaluate the original exact sequence at the few candidate indices and compare it with $y$. At the special targets $y=a_n$, a nearest-integer recovery succeeds whenever the inverse approximation error is less than $1/2$, but that nearest-integer rule is not valid for a general threshold target.

Parity-resolved families require \eqref{t1:eq:zigzagthreshold} rather than an ordinary ceiling of either branch. Integer recovery is consequently a final arithmetic step, not an exponentially small correction to a smooth function.
% ---- hand chunk 09-numerics-head.tex ----

\chapter{Reproducible numerical checks}
\label{ct:ch:numerics}

Three verification programs accompany this volume, one per absorbed source,
retained with their recorded outputs under \path{verification/source1/},
\path{verification/source2/}, and \path{verification/source3/}
(\cref{ct:app:reproduce}).  They use exact rational arithmetic (SymPy or
Python integers) for coefficient identities and \texttt{mpmath} at 90 to 110
decimal digits for the numerical comparisons; none is an interval-arithmetic
certificate.  The tables of this chapter are their recorded results, with
the source stated for each.  Where two sources tabulated the same quantity
--- the Catalan inverse at $n=100$, the involution inverse at $n=100$ ---
their numbers agree to the displayed digits.

% ---- T1 lines 1290--1359 ----
\section{Forward numerical checks}
The first source's program \texttt{verify.py}, retained under \path{verification/source1/}, computes all rational coefficient formulas with SymPy and evaluates numerical checks with mpmath at $90$ decimal digits. It compares the endpoint blocks with their finite combinatorial sums for $n=0,\ldots,20$, and the Gaussian-binomial product interpolation with its finite product for $n=1,\ldots,12$. The same source's \texttt{audit.py} checks the zigzag parity formulas for $n=1,\ldots,20$. These are $74$ high-precision identity tests, in addition to the separate Gaussian block check for $I_{100}$. They passed at the tolerances stated in the scripts.

At $n=100$, the measured forward relative errors are:
\begin{center}
\small
\begin{tabularx}{\textwidth}{@{}Xcl@{}}
\toprule
Family & Last retained index & Relative error\\
\midrule
Catalan & $j=5$, powers $n^{-j}$ & $1.130441935\cdot10^{-12}$\\
Involutions, both exact-block amplitudes & $j=6$, powers $n^{-j/2}$ & $2.760730722\cdot10^{-9}$\\
Motzkin, both endpoint amplitudes & $j=5$, powers $n^{-j}$ & $8.849657729\cdot10^{-11}$\\
Central trinomial, both endpoint amplitudes & $j=5$, powers $n^{-j}$ & $2.844260581\cdot10^{-14}$\\
\bottomrule
\end{tabularx}
\end{center}
Here ``both amplitudes'' means both exact blocks are replaced by their stated finite amplitude truncations. It does not mean that those truncated amplitudes have been evaluated as exact functions. At this $n$, the negative endpoint blocks are far below the errors caused by truncating the dominant algebraic amplitude, as the exact-block distinction predicts.

\section{Inverse numerical checks}
The targets below are exact sequence values $y=a_n$, so the true integer index is known independently of root finding. All errors shown are absolute errors of the real inverse approximation.
\begin{center}
\small
\begin{tabular}{@{}llrl@{}}
\toprule
Family & Truncation & $n$ & $|x_{\rm approx}-n|$\\
\midrule
Catalan & Lambert core only & $100$ & $8.167707070\cdot10^{-3}$\\
 & through $X^{-1}$ & $100$ & $5.188458858\cdot10^{-5}$\\
 & through $X^{-2}$ & $100$ & $1.357576716\cdot10^{-7}$\\
 & through $X^{-4}$ & $100$ & $3.304373394\cdot10^{-11}$\\
\addlinespace
Stirling, $k=3$ & logarithmic core only & $30$ & $1.424106365\cdot10^{-5}$\\
 & total marker degree $\le1$ & $30$ & $2.917192261\cdot10^{-11}$\\
 & total marker degree $\le2$ & $30$ & $5.561043091\cdot10^{-17}$\\
 & total marker degree $\le4$ & $30$ & $2.312654479\cdot10^{-27}$\\
\addlinespace
Eulerian, $m=1$ & logarithmic core only & $30$ & $4.165204886\cdot10^{-8}$\\
 & through first sector & $30$ & $5.453038651\cdot10^{-16}$\\
 & through second sector & $30$ & $9.224385539\cdot10^{-24}$\\
\addlinespace
Involutions & Lambert core only & $100$ & $4.334874595$\\
 & through $v_0$ & $100$ & $6.069814763\cdot10^{-2}$\\
 & through $v_1$ & $100$ & $1.197099872\cdot10^{-2}$\\
 & through $v_2$ & $100$ & $4.933553106\cdot10^{-4}$\\
 & through $v_3$ & $100$ & $1.029749159\cdot10^{-5}$\\
\addlinespace
Zigzag, even branch & centered Lambert core & $100$ & $9.969553297\cdot10^{-5}$\\
 & through $Z^{-1}$ & $100$ & $6.865593829\cdot10^{-10}$\\
 & through $Z^{-3}$ & $100$ & $2.069198748\cdot10^{-14}$\\
\addlinespace
Gaussian binomial, $q=2$ & quadratic core only & $20$ & $6.879303859\cdot10^{-8}$\\
 & through $T$ & $20$ & $2.526734188\cdot10^{-14}$\\
 & through $T^2$ & $20$ & $6.434897235\cdot10^{-21}$\\
\bottomrule
\end{tabular}
\end{center}
The Eulerian second-sector coefficient used here, with $a=\log2$ and $X=\log_2 y$, is
$-(X+1)^2/(2a)+(X+1)/a^2$, multiplying $2^{-2X}$.
For the involution rows, only the dominant-block algebraic inverse is used; the target is the full exact $I_{100}$. The omitted inverse oscillatory contribution is much smaller than the displayed algebraic errors. The same separation applies to the zigzag rows, where the odd-integer sectors are negligible at the displayed precision.

An additional calculation isolates the involution small block directly from the exact integrals:
\begin{align}
 \frac{A_-(100)}{A_+(100)}
 &=1.944286460629523489640756\cdot10^{-9},\nonumber\\
 \frac{A_-(100)/A_+(100)}{(\log A_+)'(100)}
 &=8.264930508117383796848111\cdot10^{-10}.
 \label{t1:eq:ivsmallnumeric}
\end{align}
The second number is the first-order displacement of the positive-block inverse at the target $I_{100}$, with positive sign because $100$ is even and $I_{100}>A_+(100)$. Applying the full inverse sector subtracts that displacement. It is not itself an independently certified root difference, and no such claim is needed for the check.
% ---- hand chunk 09-numerics-t2.tex ----

\section{Checks of the second source}
\label{ct:sec:numerics-t2}

% ---- T2 lines 1963--2052 ----
\label{t2:app:verification}

The second source's program \texttt{verify.py}, retained under \path{verification/source2/}, uses exact SymPy rational arithmetic for the coefficient identities and \texttt{mpmath} with 100 decimal digits for numerical evaluations. The companion \texttt{audit\_symbolic.py} checks several formulas independently. Both scripts are included with their actual output reports. These finite checks complement, rather than replace, the proofs.

\Needspace{12\baselineskip}
\section{Catalan inverse}

For each listed integer $n$, the target is the exact Catalan number $y=C_n$. The first column compares the Lambert core \eqref{t2:eq:catalan-core} with $n$; the second uses all three corrections in \eqref{t2:eq:catalan-inverse}.
\begin{center}
\begin{tabular}{rcc}
\toprule
$n$ & $|X-n|$ & Absolute error after three corrections\\\midrule
20 & $4.1976\times10^{-2}$ & $1.5033\times10^{-6}$\\
100 & $8.1677\times10^{-3}$ & $2.2520\times10^{-9}$\\
1000 & $8.1203\times10^{-4}$ & $2.2216\times10^{-13}$\\\bottomrule
\end{tabular}
\end{center}
The observed improvement is consistent with the proved $\bigO(X^{-4})$ inverse remainder. The exact gamma interpolation is specified; the experiment is not testing an arbitrary interpolation of the same integers.

\section{Involution forward and inverse checks}

The forward approximation uses the positive-saddle leading factor and the six amplitude coefficients through $x^{-5/2}$ in \eqref{t2:eq:involution-B-first}. It is compared with the exact integer sum \eqref{t2:eq:involution-enumeration}; hence its error includes the omitted negative saddle as well as the algebraic truncation. The inverse uses $z_0,z_1,z_2$ in \eqref{t1:eq:ivvfirst} at $y=I_n$.
\begin{center}
\begin{tabular}{rcc}
\toprule
$n$ & Relative forward error & Absolute inverse error\\\midrule
100 & $1.3669\times10^{-8}$ & $4.9336\times10^{-4}$\\
400 & $1.6562\times10^{-10}$ & $1.0608\times10^{-4}$\\
1600 & $2.4275\times10^{-12}$ & $2.2684\times10^{-5}$\\\bottomrule
\end{tabular}
\end{center}
The independent symbolic audit substitutes the displayed inverse into \eqref{t1:eq:ivVeq} and verifies exact cancellation through $t^2$. Thus the half-power coefficients are checked algebraically, not merely fitted to numerical values.

\section{The regenerated graph sector}

The connected graph counts are computed independently by the exact component-of-vertex-$1$ recurrence
\begin{equation}
 c_n=2^{\binom n2}-\sum_{j=1}^{n-1}
         \binom{n-1}{j-1}c_j2^{\binom{n-j}{2}}.
 \label{t2:eq:graph-check-recurrence}
\end{equation}
This follows by choosing the vertex set of the connected component containing label $1$, selecting its connected graph, and placing an arbitrary graph on the remaining labels. It is used only as a verification algorithm; the volume's asymptotic coefficients have the nonrecursive formula \eqref{t2:eq:graph-b-finite}.

At $y=c_n$, the successive inverse errors are
\begin{center}
\begin{tabular}{rccc}
\toprule
$n$ & Quadratic core & One sector & Two sectors\\\midrule
12 & $7.3731\times10^{-4}$ & $1.8669\times10^{-6}$ & $6.4675\times10^{-8}$\\
20 & $2.8223\times10^{-6}$ & $4.8514\times10^{-11}$ & $1.0455\times10^{-14}$\\
40 & $2.6575\times10^{-12}$ & $9.1870\times10^{-23}$ & $3.8930\times10^{-32}$\\\bottomrule
\end{tabular}
\end{center}
The dramatic second improvement illustrates why a vanished forward sector must not simply be deleted from the inverse calculation.

The audit also evaluates graph forward remainders as exact rational numbers, avoiding cancellation at insufficient numerical precision. For example, after retaining sectors through $K=2$, the ratio of the remainder to $n^3 2^{-3n}$ is approximately $-18.2542$, $-19.7600$, and $-20.5400$ at $n=20,40,80$. These values approach the predicted leading coefficient $-128/6=-64/3$.

\section{A convergent inverse and an exponential inverse}

For $S_3^{-1}$, the complete multi-index formula is truncated at total degree six. For $H^{-1}$, the terms through $X^{-5}$ are used.
\begin{center}
\begin{tabular}{rcc}
\toprule
$n$ & Stirling-II logarithmic core error & Degree-six inverse error\\\midrule
10 & $4.8582\times10^{-2}$ & $9.6875\times10^{-13}$\\
20 & $8.2157\times10^{-4}$ & $3.9637\times10^{-25}$\\
50 & $4.2827\times10^{-9}$ & $4.1486\times10^{-62}$\\\bottomrule
\end{tabular}
\end{center}
\begin{center}
\begin{tabular}{rc}
\toprule
$n$ & Harmonic inverse absolute error at $y=H_n$\\\midrule
10 & $2.1550\times10^{-10}$\\
100 & $2.9870\times10^{-17}$\\
1000 & $3.0830\times10^{-24}$\\\bottomrule
\end{tabular}
\end{center}
The first experiment uses a genuinely convergent inverse series. The second uses a finite truncation of a generally divergent asymptotic series. Similar numerical success does not erase that analytic distinction.

\section{Additional independent checks and reproducibility}

The symbolic audit verifies the first two Motzkin logarithmic coefficients, the Stirling gamma-polynomial identity through $n=8$, and the generalized-harmonic inverse coefficient $-s/24$. Direct high-precision integration checks the Motzkin, Delannoy, Schr\"oder, and two-Gaussian involution moment identities for $n=0,\ldots,5$, with observed absolute discrepancies below $10^{-85}$. These are independent checks of the exact moment normalizations, which strongly affect both leading constants and inverse cores.

To reproduce the calculations, use Python 3.10 or later with SymPy and mpmath installed, then run
\begin{lstlisting}
python verify.py
python audit_symbolic.py
\end{lstlisting}
The scripts write \nolinkurl{verification_results.json}, \nolinkurl{verification_report.txt}, and \nolinkurl{symbolic_audit_report.txt} in their own directory. Exact assertions raise an error on failure. No internet access, OEIS download, or proprietary software is needed to run them.
% ---- hand chunk 09-numerics-t3.tex ----

\section{Checks of the third source: the $q$-product families}
\label{ct:sec:numerics-t3}

% ---- T3 lines 1722--1807 ----
\label{t3:sec:numerics}

The third source's program \texttt{verify.py}, retained under \path{verification/source3/}, and its generated
\texttt{verification\_report.json}. The supplied run passed \textbf{1,090 checks} using 110 decimal digits for analytic calculations and exact Python integers or rational numbers for the enumeration and threshold tests. The checks include the small integer values of the group, Gaussian, and $q$-Catalan families; generalized Galois identities for $r=2,3$; a weighted Rogers--Szeg\H{o} identity; inverse coefficient substitutions through weight eight; the finite nonrecursive coefficient formula; endpoint inversion; fixed-radical logarithmic-phase inverses; and the two-candidate threshold rule at rational test targets.

For the inverse tables below, a known real value $x_*$ is chosen and
$Y=F(x_*)$ is computed from the specified interpolation. The reported error is
\begin{equation}\label{t3:eq:tableerror}
 E_M=\left|s+\sum_{m=1}^M d_m(s)t_0^m-x_*\right|.
\end{equation}
The displayed values are rounded; the JSON report retains additional digits. The truncation index $M$ always refers to the chosen small variable, so a zero weight still counts toward $M$. In particular the symplectic row uses $t_0=Q^s$, not the reduced variable $Q^{2s}$.

\begin{table}[htbp]
\centering\small
\setlength{\tabcolsep}{4pt}
\begin{tabular}{@{}lrrrrr@{}}
\toprule
Family & $x_*$ & $E_0$ & $E_2$ & $E_4$ & $E_8$\\
\midrule
$\GL$ & $8.25$ & $2.874\times10^{-4}$ & $2.674\times10^{-10}$ & $6.375\times10^{-16}$ & $6.908\times10^{-27}$ \\
Central Gaussian & $8.25$ & $5.738\times10^{-4}$ & $1.25\times10^{-9}$ & $7.171\times10^{-15}$ & $4.475\times10^{-25}$ \\
$q$-Catalan & $8.25$ & $4.578\times10^{-4}$ & $1.063\times10^{-9}$ & $6.282\times10^{-15}$ & $3.912\times10^{-25}$ \\
Symplectic & $8.25$ & $1.526\times10^{-7}$ & $1.983\times10^{-13}$ & $4.006\times10^{-19}$ & $2.455\times10^{-30}$ \\
Unitary, even sheet & $8.25$ & $9.558\times10^{-5}$ & $1.484\times10^{-10}$ & $3.491\times10^{-16}$ & $2.874\times10^{-27}$ \\
Unitary, odd sheet & $8.25$ & $9.589\times10^{-5}$ & $1.489\times10^{-10}$ & $3.5\times10^{-16}$ & $2.881\times10^{-27}$ \\
Galois $r=2$, $\eps=0$ & $16.25$ & $1.512\times10^{-3}$ & $4.131\times10^{-9}$ & $2.824\times10^{-14}$ & $2.62\times10^{-24}$ \\
Galois $r=2$, $\eps=1$ & $17.25$ & $1.007\times10^{-3}$ & $1.307\times10^{-9}$ & $4.276\times10^{-15}$ & $9.105\times10^{-26}$ \\
Galois $r=3$, $\eps=0$ & $18.25$ & $6.633\times10^{-3}$ & $3.316\times10^{-7}$ & $1.318\times10^{-12}$ & $3.45\times10^{-19}$ \\
\bottomrule
\end{tabular}
\caption{Quadratic-phase examples at $q=2$. Absolute inverse errors after selected exponential weights.}
\end{table}

\begin{table}[htbp]
\centering\small
\setlength{\tabcolsep}{4pt}
\begin{tabular}{@{}lrrrrr@{}}
\toprule
Family & $x_*$ & $E_0$ & $E_2$ & $E_4$ & $E_8$\\
\midrule
Fixed rank $k=3$ & $12.25$ & $6.913\times10^{-4}$ & $3.333\times10^{-11}$ & $3.334\times10^{-18}$ & $5.338\times10^{-32}$ \\
$q$-factorial, base $1/2$ & $12.25$ & $2.962\times10^{-4}$ & $2.556\times10^{-11}$ & $5.128\times10^{-18}$ & $3.908\times10^{-31}$ \\
Prime-power index $p=2$ & $16.25$ & $5.682\times10^{-3}$ & $2.537\times10^{-9}$ & $2.81\times10^{-15}$ & $2.754\times10^{-27}$ \\
Fixed radical $R=6$ & $36.25$ & $5.334\times10^{-6}$ & $5.334\times10^{-6}$ & $1.212\times10^{-9}$ & $3.321\times10^{-17}$ \\
\bottomrule
\end{tabular}
\caption{Linear and logarithmic phases. The last two rows use the exact $W_{-1}$ center.}
\end{table}

For the finite-limit inverse at $q=2$, take $x_*=16.25$ and
\[
 \eps=\mathcal A_{1/2}(2^{-16.25})
      =1.28310884632239305154\ldots\times10^{-5}.
\]
Truncating the convergent series \eqref{t3:eq:endpoint-all} gives:
\begin{center}
\begin{tabular}{@{}rrrrr@{}}
\toprule
$E_0$ & $E_1$ & $E_2$ & $E_4$ & $E_8$\\
\midrule
$3.085\times10^{-6}$ & $1.414\times10^{-12}$ & $1.075\times10^{-18}$ & $7.137\times10^{-30}$ & $2.592\times10^{-52}$ \\
\bottomrule
\end{tabular}
\end{center}

As a separate check of the singular transition, set $\tau=1$ and retain
terms through $h^3$ in \eqref{t3:eq:central-double-asympt}, omitting its much
smaller modular factor. The observed errors are consistent with the proved
$O(h^5)$ remainder:
\begin{center}
\begin{tabular}{@{}rrr@{}}
\toprule
$h$ & Exact $\log H_{1,1}(1/h)$ & Absolute error through $h^3$\\
\midrule
$.2$ & $8.475418514412$ & $4.856\times10^{-7}$ \\
$.1$ & $17.97956130094$ & $1.551\times10^{-8}$ \\
$.05$ & $37.31564489494$ & $4.874\times10^{-10}$ \\
\bottomrule
\end{tabular}
\end{center}

The sparse symplectic and fixed-radical rows also verify the predicted missing
weights. Such tests are more informative than a single agreement of the
leading approximation: a wrong small parameter, a lost parity, or a missing
curvature term usually changes the next nonzero error by several orders of
magnitude.
% ---- hand chunk 09-synthesis-head.tex ----

\chapter{Synthesis: scales, closure, and what has been proved}
\label{ct:ch:synthesis}

The families of this volume extend the canonical volume's inversion
architecture in several directions at once.  Balanced gamma products show
how Bernoulli-polynomial data turn a factorial formula into all inverse
coefficients.  Stirling and Eulerian columns show that finite forward
exponential sums have infinite but convergent inverse expansions, with
generalized powers and polynomial-logarithmic amplitudes.  Endpoint moment
sequences show how an oscillatory sector is fixed by an explicit
negative-support block.  Involutions add a half-power saddle calculus and an
unbounded inverse drift before the first exponentially small correction.
Zigzag numbers turn inverse action collisions into finite
ordered-factorization sums.  Connected graphs give every exponential layer
of a superexponential class and an inverse rigorous on the range.  Necklaces
show where a finite exponential grid fails and how fixed-radical sheets
restore convergence.  The $q$-products give convergent inverse transseries
with quadratic, linear, and theta-modulated phases, and a singular
transition at $q\to1$ in which the tail joins the dominant phase.

% ---- T2 lines 1793--1905 ----
\label{t2:sec:synthesis}

\section{A comparative map}

Table~\ref{t2:tab:classes} summarizes the different dominant balances. A symbol such as $X$ in an exponential action always denotes the appropriate exact or formally specified dominant inverse, not an interchangeable rough approximation.

\begin{table}[htbp]
\centering
\small
\renewcommand{\arraystretch}{1.25}
\begin{tabular}{P{34mm}P{46mm}P{64mm}}
\toprule
Family & Forward structure & Inverse structure\\\midrule
Catalan, Fuss--Catalan, fixed-height rectangles
 & $C\rho^x x^\beta\sum a_jx^{-j}$
 & $W_{-1}$ core and inverse powers of that core\\
Fixed-column Stirling II
 & Finite sum of real exponentials
 & Convergent multivariate generalized powers of $1/y$\\
Fixed-column Stirling I
 & $\Gamma(x)$ times polynomials in $\log x$
 & $W_0$ core; $\log\log X/\log X$ and rational-logarithmic corrections\\
Motzkin
 & Two exact endpoint moments; relative $3^{-x}\cos\pi x$
 & $W_{-1}$ core; exact dominant inverse with oscillatory exponential corrections\\
Involutions
 & $(x/\ee)^{x/2}\ee^{\sqrt x}$; relative $\ee^{-2\sqrt x}\cos\pi x$
 & $W_0$ core; half-powers with rational-log coefficients, then stretched-exponential sectors\\
Alternating permutations
 & $\Gamma(x+1)(2/\pi)^x$ times odd Dirichlet sums
 & Factorial core and prime-generated exponential sectors\\
Connected graphs
 & $2^{x(x-1)/2}\sum P_k(x)2^{-kx}$ on integers
 & Quadratic core; rational-function sector coefficients; rigorous range inverse\\
Necklaces, Lyndon words
 & $q^n/n$ times divisor-dependent sectors
 & $W_{-1}$ core; arithmetic finite-cutoff inverses with accumulating actions\\
Harmonic functions
 & Logarithmic growth, or convergence to $\zeta(s)$
 & Exponential inverse, or endpoint generalized Puiseux expansion\\
\bottomrule
\end{tabular}
\caption{Dominant structures and their inverse scales. Delannoy and Schr\"oder inverses fall in the first row's balanced class, with the moment interpolation specified in Section~\ref{ct:ch:moments}.}
\label{t2:tab:classes}
\end{table}

The negative Lambert branch is not a rare special case. Any balanced exponentially growing sequence with an algebraic prefactor $x^\beta$, $\beta<0$, leads to it. By contrast, factorial growth has a positive Lambert core because the $x\log x$ term changes the dominant balance. Connected graph enumeration is faster still in logarithmic phase, but its inverse core is simpler: a quadratic root.

\section{Why nonlinear inversion creates sectors but not arbitrary constants}

\begin{proposition}[Additive support closure under inverse transport]
\label{t2:prop:support}
Suppose an exact small correction, expressed at the exact dominant inverse $x_0$, has sectors
\[
 \varepsilon(x_0)=\sum_{j=1}^r
       a_j(x_0)\exp(-\lambda_j\Omega(x_0)),\qquad\lambda_j>0.
\]
Assume the coefficient class is closed under products and under
$\mathcal D=((\log F_0)'(x_0))^{-1}\dd/\dd x_0$, including the factors introduced by differentiating $\Omega$. Then the formal inverse produced by Theorem~\ref{ct:rmk:transport} has exponential support contained in
\[
 \left\{\sum_{j=1}^rm_j\lambda_j:
              m_j\in\mathbb Z_{\ge0},\quad\sum_jm_j>0\right\}.
\]
At every finite action cutoff, the coefficients are finite sums. The same conclusion holds for any locally finite positive action set.
\end{proposition}
\begin{proof}
Expanding $\log(1+\varepsilon)$ produces only products of input sectors, so its actions are sums of the $\lambda_j$. Differentiation multiplies a sector by coefficient factors and does not change its exponential rate in the stated coefficient class. Formula~\eqref{ct:eq:transport} uses only these two operations. Positive local finiteness makes extraction below a fixed action cutoff finite. No new independent constant appears in any operation.
\end{proof}

The proposition explains the nonzero second inverse graph sector, the composite odd-integer sectors of the alternating-permutation inverse, and the higher involution actions $2m\sqrt{x_+}$. It does not manufacture Stokes constants for a dominant asymptotic expansion that never specified a small sector in the first place. Those constants belong to the choice and analysis of the original function, not to formal reversion alone.

For oscillatory coefficients, multiplication adds frequencies and differentiation retains them. Thus nonlinear inversion can generate $\cos(2\pi x_0)$ and $\sin(2\pi x_0)$ even if the forward interpolation had only $\cos(\pi x)$. Treating an oscillatory prefactor as a constant during differentiation gives incorrect second and higher layers.

\section{Status of the results}

The following distinctions summarize the logical scope of the volume.

\paragraph{Convergent identities.}
The finite-exponential inverse of Theorem~\ref{t2:thm:finite-exp} is locally convergent in its small variables. The Motzkin beta representations converge absolutely. The alternating-permutation Dirichlet sums and Euler products converge absolutely in the stated half-plane. The finite enumeration and divisor formulas are exact.

\paragraph{All-orders asymptotic expansions with error control.}
The gamma-quotient, fixed-column Stirling I, endpoint-moment, moving-saddle, and harmonic results have a remainder at every prescribed finite algebraic order. Their inverse errors follow from a quantified slope and a forward residual. These statements do not assert convergence of the corresponding infinite Bernoulli or saddle series.

\paragraph{Exactly normalized exponentially small sectors.}
Motzkin and involution sectors are defined by exact separate integrals. Alternating-permutation sectors are defined by exact Dirichlet sums. Their inverse transport formulas refer to the exact dominant inverse. This gives a meaningful all-sector formal expansion without confusing the error of a fixed algebraic truncation with the size of a subdominant sector.

\paragraph{Lattice transseries and range inverses.}
Theorem~\ref{t2:thm:graphs} proves every fixed exponential layer for connected graph counts, and Corollary~\ref{t2:cor:graph-range} proves the corresponding inverse on the range. The necklace results prove finite-cutoff arithmetic expansions and finite-model range inverses. Neither argument asserts a unique global real interpolation of the sequence.

\paragraph{Not claimed.}
No global resurgence or Borel summation theorem is established for all the examples. No uniform expansion is asserted as a fixed combinatorial parameter grows with $n$. No lower-endpoint Stokes coefficient is supplied for the Delannoy or Schr\"oder moment functions. No infinite noninteger divisor interpolation is assumed for necklaces. These are separate questions rather than consequences of the formulas already proved.

\section{A reproducible route from an enumeration formula to its inverse}

The derivations suggest a practical workflow. First choose an exact representation: a gamma quotient, finite exponential sum, positive moment, formal component generating function, or divisor sum. Then decide whether the desired output is a continuous inverse, a range inverse, or a staircase threshold. This prevents interpolation conventions from being inserted accidentally in the middle of an asymptotic calculation.

Next extract the logarithmic dominant phase and invert it as accurately as possible in closed form. The useful cores here are $ax+\beta\log x$, $\kappa x\log x+\lambda x$, and a quadratic polynomial. Keeping the core exact avoids repeated expansion of its internal logarithms. Normalize the remaining correction so that it has positive order, using a relative displacement if the absolute correction grows.

Then compute only the coefficients required by the chosen precision. Formula~\eqref{t2:eq:E} handles exponentiation, \eqref{t2:eq:L} handles logarithms, and the relevant finite Lagrange formula handles reversion. For multiple sectors, specify an action cutoff before truncating algebraic amplitudes. For arithmetic actions with a finite accumulation point, specify a finite divisor cutoff instead of claiming a complete bounded-action truncation.

Finally substitute the inverse approximation into the \emph{forward} equation. The residual is generally easier to estimate than the inverse error directly. Divide it by a justified slope bound, and use a rounding interval rather than an unqualified ceiling. This last step turns formal coefficient algebra into a trustworthy numerical or discrete conclusion.

\section{Further research directions}

Several concrete extensions are suggested by the results, but are not used as premises above.

One direction is to obtain exponentially improved remainders for the exact positive moment functions. For Motzkin and involution numbers, the small sector is already normalized by a real integral, so a useful next theorem would compare optimal truncation of the dominant amplitude directly with the known second sector. For Delannoy and Schr\"oder functions, the first task is instead to determine a contour- or summation-specific subdominant decomposition before proposing an inverse Stokes analysis.

A second direction is a uniform two-parameter theory. The fixed-$k$ Stirling expansions change when $k$ grows, and fixed-height rectangular tableaux change when the height is comparable with the width. The coefficient engines remain formally available, but the present remainder arguments are not uniform in those regimes. Identifying the transition scales and the appropriate dominant inverse would yield genuinely different asymptotic problems.

A third direction is an arithmetic transseries algebra for divisor sums. It should retain divisibility indicators, permit the action accumulation at $\log q$, and state exactly which products and inverse operations remain summable. The finite-cutoff estimates in Proposition~\ref{t2:prop:necklace-cutoff} provide a concrete benchmark for such a theory. Merely replacing indicators by analytic expressions and summing indefinitely is not an adequate definition.

Finally, the connected-component proof suggests a general theorem for labeled structures whose total counts have a strongly convex superexponential logarithm. A unique giant component can again dominate formal logarithmic extraction. The appropriate hypotheses should control both the smaller-component reciprocal coefficients and the contributions of partitions with no giant block. Under such hypotheses one can seek analogues of \eqref{t2:eq:graph-expansion} and then apply the same inverse coefficient transport. The graph theorem supplies a complete worked case, not a proof for every superexponential labeled class.
% ---- hand chunk 09-synthesis-t3.tex ----

\section{The $q$-product families in the same view}
\label{ct:sec:synthesis-t3}

% ---- T3 lines 1810--1828 ----
\label{t3:sec:conclusion}

\section{A unified coefficient calculus}

The examples exhibit the same underlying structure at three levels. Finite $q$-products yield explicit logarithmic coefficients by \eqref{t3:eq:product-log-coeff}. Theta-weighted flag sums yield them by \eqref{t3:eq:h-coeff} and \eqref{t3:eq:Galois-b}. Fixed-radical divisor sums yield them from the logarithm of the finite polynomial $E_R$. In all three cases, inversion is a finite operation at every exponential weight, expressed by \eqref{t3:eq:finite-dm} or \eqref{t3:eq:general-derivative}.

The coefficients record two distinct nonlinear effects: differentiation of the dominant inverse phase, and multiplication or shifting of exponential sectors. The first effect introduces factors such as $D^{-1}$ and derivatives of $D$; the second produces the composition sums $C_{m,k}$. Keeping these mechanisms separate explains the universality of the formulas without suppressing family-specific cancellations.

\section{Exact discrete identities can improve analytic interpolation}

The Galois-number construction illustrates a useful general strategy. A finite sum is extended to a lattice sum because an infinite-product factor vanishes outside the original finite domain. After an Euler expansion, completing a quadratic form converts every coefficient into a theta constant on a translated lattice coset. The result is an entire coefficient-generating function and an explicit residue-sheet interpolation, not only a leading asymptotic estimate.

For a more general finite-field flag or partition-function problem, the analogous questions are concrete: does a product factor enforce the finite domain by zeros; does the dominant exponent have a positive definite quadratic form on the constraint lattice; and do the coefficient translations preserve finitely many lattice cosets? When all three conditions hold, the proof of Theorem~\ref{t3:thm:theta-tail} suggests a direct route to all sectors and then to their inverses.

\section{Problems left open by this organization}

Several extensions merit separate treatment. The first is a uniform theory in which the flag length $r$ grows with $n$: the lattice dimension and the exponential spacing $h/r$ then both vary, so the fixed-$r$ coefficient bounds are not uniform. The second is a complete resurgent analysis of the finite-size tail in the double-scaling limit; the exact modular product isolates one contribution but does not replace that analysis. The third is efficient interval evaluation of high-dimensional theta constants, allowing the explicit inverse series to serve as certified threshold algorithms rather than only high-precision approximations.

These questions do not affect the established fixed-parameter results. For fixed $q>1$ and fixed rank or flag length, the product and theta inverses above are convergent, their coefficients are explicit to every order, and their integer thresholds are controlled by the stated rounding or exact-comparison rules.
% ---- hand chunk 10-app-head.tex ----

\appendix
\chapter{Coefficient formulas collected in one place}
\label{ct:app:formulas}

% ---- T2 lines 1909--1960 ----
\label{t2:app:coefficients}

The formulas below show explicitly that the constructions are not dependent on unlisted auxiliary recurrences. Recurrences may be more efficient in a computer implementation, but a finite coefficient or partition formula determines every requested order.

\begin{longtable}{P{43mm}P{104mm}}
\toprule
Operation or family & Finite coefficient specification\\\midrule
\endfirsthead
\toprule
Operation or family & Finite coefficient specification\\\midrule
\endhead
Exponential amplitude &
$\displaystyle [t^n]\ee^{\sum q_jt^j}=\cE_n(q)$, with the partition sum \eqref{t2:eq:E}.\\[3pt]
Logarithmic amplitude &
$\displaystyle [t^n]\log(1+\sum a_jt^j)=\cL_n(a)$, with \eqref{t2:eq:L}.\\[3pt]
Gamma quotient &
Bernoulli-polynomial sum \eqref{t2:eq:gamma-q}; normalized amplitudes $\cE_n(q)$.\\[3pt]
Balanced inverse &
$\displaystyle d_n=\sum_{m=1}^n\frac1m[t^nu^{m-1}]\Psi(t,u)^m$, equation~\eqref{t2:eq:balanced-d}.\\[3pt]
Finite exponential inverse &
Multi-index generalized-binomial sum \eqref{t2:eq:finite-exp-inverse}.\\[3pt]
Fixed-column Stirling I &
$\displaystyle P_{k,m}(\ell)=[z^{k-1}]\ee^{\ell z}\cE_m(r(z))/\Gamma(1+z)$; inverse by scaled Lagrange extraction.\\[3pt]
Moment endpoint &
$\displaystyle a_m=\sum_{j=0}^m\binom\theta j(-\kappa)^j(\nu)_jR_{m-j}(\alpha,\alpha+\nu+j)$.\\[3pt]
Involution saddle &
Gaussian partition sum \eqref{t2:eq:involution-S}, with prefactor \eqref{t2:eq:involution-B}. Only weights $\sum(k-2)m_k\le N$ contribute through $t^N$.\\[3pt]
Involution inverse &
After removing $z_0=-2/\ell$, finite extraction \eqref{t1:eq:ivVclosed}.\\[3pt]
Exact sector transport &
$\displaystyle \sum_{m\ge1}\frac{(-1)^m}{m!}\partial_L^{m-1}(g'h^m)$, truncated at the specified action.\\[3pt]
Connected graphs &
Reciprocal partition sum \eqref{t2:eq:graph-b-finite}; polynomial $P_k$ in \eqref{t2:eq:graph-P}; inverse extraction \eqref{t2:eq:graph-d}.\\[3pt]
Harmonic inverse &
$\displaystyle h_m=-\frac{[z^m]\ee^{(2m-1)C(z)}}{2m-1}$.\\[3pt]
Generalized harmonic inverse &
$\displaystyle h_{s,m}=-\frac{[z^m]D_s(z)^{-(2m-1)/(s-1)}}{2m-1}$.\\
\bottomrule
\end{longtable}

For the moment sequences, a useful compact parameter dictionary for \eqref{t2:eq:endpoint-coeff} is
\begin{center}
\begin{tabular}{lcccc}
\toprule
Sector & $\alpha$ & $\nu$ & $\theta$ & $\kappa$\\\midrule
Motzkin positive & $1$ & $3/2$ & $1/2$ & $3/4$\\
Motzkin negative & $1$ & $3/2$ & $1/2$ & $1/4$\\
Delannoy & $1$ & $1/2$ & $-1/2$ & $(3+2\sqrt2)/(4\sqrt2)$\\
Large Schr\"oder & $0$ & $3/2$ & $1/2$ & $(3+2\sqrt2)/(4\sqrt2)$\\\bottomrule
\end{tabular}
\end{center}
The normalization constants are outside this dictionary and are specified in the relevant section. Mixing those constants with normalized amplitude coefficients is a common source of inverse-core errors.
% ---- T3 lines 1899--1899 ----
The finite nonrecursive inverse coefficient is \eqref{t3:eq:finite-dm}; the general analytic-phase version is \eqref{t3:eq:general-derivative}. Product families are specified by \eqref{t3:eq:product-log-coeff}. All generalized Galois theta-tail coefficients are given by \eqref{t3:eq:h-coeff}, their logarithmic coefficients by \eqref{t3:eq:Galois-b}, and their full inverse sheets by \eqref{t3:eq:Galois-inverse-full}. The endpoint inverse is \eqref{t3:eq:endpoint-all}. Arithmetic fixed-radical sheets are \eqref{t3:eq:fixedradical-function} and \eqref{t3:eq:radical-inverse}; the uniform integer threshold rule is \eqref{t3:eq:two-candidate}. The simultaneous $q\to1$ phase is \eqref{t3:eq:Sphase}, with inverse coefficients \eqref{t3:eq:crossover-all}.
% ---- hand chunk 10-app-audit-head.tex ----

\chapter{An independent involution coefficient audit}
\label{ct:app:audit}

% ---- T1 lines 1371--1400 ----
\label{t1:app:audit}
The finite Gaussian-moment formula \eqref{t1:eq:ivcoeff} can be checked without using numerical fitting. The exact recurrence
\[
 I_n=I_{n-1}+(n-1)I_{n-2}
\]
follows by considering whether the last element is fixed or paired. Let $t=n^{-1/2}$, $S(t)=\sum b_jt^j$, and
\begin{equation}
 R_k(t)=\exp\left\{
 \left(\frac1{2t^2}-\frac k2\right)\log(1-kt^2)+\frac k2
 +\frac{\sqrt{1-kt^2}-1}{t}\right\}.
 \label{t1:eq:auditR}
\end{equation}
Dividing the recurrence by the positive saddle carrier gives the formal identity
\begin{align}
 S(t)-tR_1(t)S\!\left(\frac{t}{\sqrt{1-t^2}}\right)
 -(1-t^2)R_2(t)S\!\left(\frac{t}{\sqrt{1-2t^2}}\right)=0.
 \label{t1:eq:auditrec}
\end{align}
The negative saddle is flat in $t$ on the positive real axis and therefore does not alter this formal algebraic identity. Substituting $b_0,\ldots,b_8$ from \eqref{t1:eq:ivcoeff} makes every coefficient through $t^{10}$ vanish exactly in rational arithmetic. The script \texttt{audit.py} performs this calculation using finite polynomial products.

This audit also identifies a version-specific discrepancy worth recording for reproducibility. The displayed theorem in the HTML rendering of arXiv:2410.16334v1 \cite{ekz} gives the coefficient of $n^{-3}$ as $-562799/47775744$. Formula \eqref{t1:eq:ivcoeff} instead gives
\[
 b_6=\frac{6340449533}{687970713600}.
\]
Keeping the preceding coefficients fixed and replacing this $b_6$ by the displayed HTML value leaves the nonzero recurrence residual
\begin{equation}
 \frac{14444755133}{114661785600}\,t^8+\OO(t^9).
 \label{t1:eq:auditdiscrepancy}
\end{equation}
Later coefficients cannot cancel this term: changing $b_j$ first affects the normalized residual at degree $j+2$. This observation concerns that particular displayed formula, not the paper's separately linked high-order data, software, or any subsequent version. The coefficients used in this volume come from the independently proved Gaussian-moment formula rather than from that display.
% ---- hand chunk 10-app-wolfram-head.tex ----

\chapter{Wolfram Language recipes}
\label{ct:app:wolfram}

The three sources each supplied Wolfram Language translations of their
coefficient formulas.  They are collected here unchanged, one section per
source, with their own naming; none was executed in a Wolfram kernel in
preparing the sources or this volume, and the executed checks are the
Python programs of \cref{ct:ch:numerics}.  \texttt{ProductLog[0,z]} and
\texttt{ProductLog[-1,z]} are the Lambert branches $W_0,W_{-1}$
\cite{wlproductlog}; \texttt{QBinomial}, \texttt{QPochhammer},
\texttt{StirlingS1}, \texttt{StirlingS2}, \texttt{DirichletBeta},
\texttt{HarmonicNumber} and \texttt{CatalanNumber} provide direct
comparisons with the standard functions
\cite{wlqbinomial,wlqpochhammer,wlstirling1,wlstirling2,wlbeta,wlharmonic,wlcatalan}.

\section{Recipes of the first source}
\label{ct:app:wolfram-t1}

% ---- T1 lines 1406--1454 ----
The following Wolfram Language definitions are retained as \texttt{coefficients.wl} under \path{verification/source1/}. They express the same finite coefficient formulas using symbolic series operations. They were not executed in a Wolfram kernel in preparing this volume; the executed checks are the Python ones. \texttt{ProductLog[0,z]} and \texttt{ProductLog[-1,z]} correspond to the Lambert branches used above \cite{wlproductlog}; \texttt{QBinomial}, \texttt{QPochhammer}, and \texttt{StirlingS2} provide direct comparisons with the standard functions \cite{wlqbinomial,wlqpochhammer,wlstirling2}.
\begin{lstlisting}[language=Mathematica]
(* Exact coefficient formulas. These definitions use no machine literals. *)
ClearAll[ExpCoefficient, GammaRatioCoefficient, EndpointCoefficient,
  NormalHalfMoment, InvolutionCoefficient, QLogCoefficient,
  QInverseCoefficient, StirlingColumnInverse];
ExpCoefficient[d_List, n_Integer] /; 0 <= n <= Length[d] :=
 Module[{t}, SeriesCoefficient[
   Exp[Sum[d[[j]] t^j, {j, 1, n}]], {t, 0, n}]];
GammaRatioCoefficient[aa_, bb_, n_Integer] /; n >= 0 :=
 ExpCoefficient[Table[
   (-1)^(j + 1) (BernoulliB[j + 1, aa] -
      BernoulliB[j + 1, bb])/(j (j + 1)), {j, 1, n}], n];
EndpointCoefficient[alpha_, c_, n_Integer] /; n >= 0 :=
 Total[Table[Binomial[alpha, j] (-c)^j Pochhammer[alpha + 1, j]
   GammaRatioCoefficient[1, alpha + j + 2, n - j], {j, 0, n}]];
NormalHalfMoment[n_Integer] /; n >= 0 :=
 Sum[n! (1/2)^(n - 2 j)/(4^j j! (n - 2 j)!),
   {j, 0, Floor[n/2]}];
InvolutionCoefficient[m_Integer] /; m >= 0 :=
 Module[{t, v, poly},
  poly = Expand[SeriesCoefficient[
    Exp[Sum[(-1)^(r + 1) v^(r + 2) t^r/(r + 2), {r, 1, m}]],
    {t, 0, m}]];
  Total[Table[Coefficient[poly, v, j] NormalHalfMoment[j],
    {j, 0, 3 m}]]];
QLogCoefficient[r_Integer, q_] /; r >= 1 :=
 -2/(r (q^r - 1)) + If[EvenQ[r], 2/(r (q^(r/2) - 1)), 0];
QInverseCoefficient[n_Integer, x_Symbol, a_, q_] /; n >= 1 :=
 Module[{t, hh, apply},
  hh = Sum[QLogCoefficient[r, q] t^r, {r, 1, n}];
  apply[f_] := (D[f, x] - a n f)/(2 a x);
  Simplify[Total[Table[(-1)^m/m! Coefficient[Expand[hh^m], t, n]
    Nest[apply, 1/(2 a x), m - 1], {m, 1, n}]]]];
(* Total-degree truncation, exponential cost in k and maxDegree. *)
StirlingColumnInverse[k_Integer, y_, maxDegree_Integer] /;
    k >= 2 && maxDegree >= 0 :=
 Module[{a = Log[k], xx, kap, terms, indices},
  xx = Log[k! y]/a;
  kap = Table[Log[k/j]/a, {j, 1, k - 1}];
  terms = Table[(-1)^(k - j) Binomial[k, j] (j/k)^xx,
    {j, 1, k - 1}];
  indices = Select[Tuples[Range[0, maxDegree], k - 1],
    1 <= Total[#] <= maxDegree &];
  xx - Total[Map[Function[nu,
    Product[(kap.nu) - h, {h, 1, Total[nu] - 1}]
    Product[terms[[j]]^nu[[j]]/nu[[j]]!, {j, 1, k - 1}]],
    indices]]/a];
\end{lstlisting}
% ---- hand chunk 10-app-wolfram-t2.tex ----

\section{Recipes of the second source}
\label{ct:app:wolfram-t2}

% ---- T2 lines 2056--2100 ----
\label{t2:app:wolfram}

The user-facing Wolfram functions \texttt{CatalanNumber}, \texttt{StirlingS1}, \texttt{StirlingS2}, \texttt{HarmonicNumber}, \texttt{DirichletBeta}, and \texttt{ProductLog} provide useful reference definitions \cite{wlcatalan,wlstirling1,wlstirling2,wlharmonic,wlbeta,wlproductlog}. The following symbolic recipes implement the coefficient formulas rather than asking a numerical inverse solver to infer an asymptotic expansion. They are supplied as a translation of the proved formulas; the executed checks accompanying this volume used Python, not a Wolfram kernel.

Each row of \texttt{data} below is $\{a_i,b_i,\epsilon_i\}$ for a gamma factor $\Gamma(a_ix+b_i)^{\epsilon_i}$.
\begin{lstlisting}
ClearAll[gammaLogCoefficient, balancedInverseCoefficient,
  harmonicInverseCoefficient, catalanCore];

gammaLogCoefficient[data_List, r_Integer] /; r >= 1 :=
  (-1)^(r + 1)/(r (r + 1)) Total[
    (#[[3]] BernoulliB[r + 1, #[[2]]]/#[[1]]^r) & /@ data
  ];

balancedInverseCoefficient[q_List, a_, beta_, n_Integer] /;
    n >= 1 && Length[q] >= n :=
  Module[{t, u, psi},
    psi = -(beta Log[1 + t u] +
       Sum[q[[j]] (t/(1 + t u))^j, {j, 1, n}])/a;
    Simplify[Total[Table[
      SeriesCoefficient[
        SeriesCoefficient[psi^m, {t, 0, n}],
        {u, 0, m - 1}]/m,
      {m, 1, n}]]]
  ];

harmonicInverseCoefficient[m_Integer] /; m >= 1 :=
  Module[{z, c},
    c = -Sum[BernoulliB[2 j, 1/2] z^j/(2 j), {j, 1, m}];
    -SeriesCoefficient[Exp[(2 m - 1) c], {z, 0, m}]/(2 m - 1)
  ];

catalanCore[y_] := -3/(2 Log[4]) ProductLog[-1,
  -(2 Log[4])/3 (1/(Sqrt[Pi] y))^(2/3)];

catData = {{2, 1, 1}, {1, 1, -1}, {1, 2, -1}};
catQ = Table[gammaLogCoefficient[catData, j], {j, 1, 6}];
Table[balancedInverseCoefficient[catQ, Log[4], -3/2, n],
  {n, 1, 3}]
Table[harmonicInverseCoefficient[m], {m, 1, 5}]
\end{lstlisting}

The intended first output is the three coefficients in \eqref{t2:eq:catalan-inverse}; the second is \eqref{t2:eq:harmonic-coeff}. Exact input constants should be kept exact until a final numerical evaluation. For very large targets, use high precision and work with logarithmic phases: constructing an enormous combinatorial number at machine precision and then taking its logarithm can lose the small corrections being studied.

A safe finite-order implementation has three separate precision controls: algebraic order within a sector, exponential action cutoff, and numerical working precision. They should not be represented by a single undocumented integer. For example, computing the involution amplitude through $x^{-5/2}$ does not resolve the $\ee^{-2\sqrt x}$ sector, and six total degrees in the Stirling-II inverse do not mean six successive dominance-ordered terms.
% ---- hand chunk 10-app-wolfram-t3.tex ----

\section{Recipes of the third source}
\label{ct:app:wolfram-t3}

% ---- T3 lines 1832--1882 ----
\label{t3:app:wolfram}

The following functions use the documented \texttt{QPochhammer},
\texttt{QBinomial}, and \texttt{ProductLog} conventions
\cite{wlqpochhammer,wlqbinomial,wlproductlog}. The list \texttt{bs} contains
$\{b_1,\ldots,b_M\}$, without a leading zero.

\begin{lstlisting}
ClearAll[alpha, pInf, aTail, inverseCoefficient];
alpha[m_Integer, Q_] := Q^m/(m (1 - Q^m));
pInf[Q_] := QPochhammer[Q, Q];
aTail[Q_, t_] := -Log[QPochhammer[Q t, Q]];

inverseCoefficient[m_Integer, a_, lam_, D_, bs_List] :=
 Module[{u, bpoly},
  bpoly = Sum[bs[[j]] u^j, {j, 1, m}];
  -Sum[
    Coefficient[Expand[bpoly^k], u, m]/k *
     Sum[Binomial[k + j - 1, j] If[j == 0, 1, a^j] *
       (m lam)^(k - 1 - j)/
       ((k - 1 - j)! D^(k + j)), {j, 0, k - 1}],
    {k, 1, m}]
 ];

(* General linear group interpolation: work with its logarithm. *)
ClearAll[glLog, glInverseTruncated];
glLog[x_, q_] := x^2 Log[q] + Log[pInf[1/q]] +
                 aTail[1/q, q^(-x)];
glInverseTruncated[logY_, q_, M_Integer] :=
 Module[{h = Log[q], Q = 1/q, s, D, bs},
  s = Sqrt[(logY - Log[pInf[Q]])/h];
  D = 2 h s;
  bs = Table[alpha[j, Q], {j, 1, M}];
  s + Sum[inverseCoefficient[m, h, h, D, bs] q^(-m s),
          {m, 1, M}]
 ];

(* Exact integer examples; q is a prime power for field counting. *)
ClearAll[galoisNumber, qCatalanBinomial, irreducibleCount];
galoisNumber[n_Integer, q_] := Sum[QBinomial[n, k, q], {k, 0, n}];
qCatalanBinomial[n_Integer, q_] :=
  (q - 1) QBinomial[2 n, n, q]/(q^(n + 1) - 1);
irreducibleCount[n_Integer, q_Integer] /; n >= 1 && q >= 2 :=
  Total[(MoebiusMu[#] q^(n/#)) & /@ Divisors[n]]/n;

(* The leading inverse of q^x/x, on the large real branch. *)
ClearAll[necklaceCenter];
necklaceCenter[y_, q_] := -ProductLog[-1, -Log[q]/y]/Log[q];
\end{lstlisting}

Sequence indices are nonnegative; \texttt{irreducibleCount} requires a positive index. Inversion targets are positive and sufficiently large, with $q>1$ except where a base $Q\in(0,1)$ is explicitly used. The binomial-version $q$-Catalan function at $q=1$ is obtained by a limit. For symbolic calculations, take $q$ exact, such as $q=2$, and request numerical precision only after forming the expression. Near an integer threshold, verify candidate indices by exact comparison rather than relying on a machine-precision ceiling of \texttt{ProductLog}.
% ---- hand chunk 10-app-reproduce.tex ----

\section{Reproduction}
\label{ct:app:reproduce}

The volume is built with three \texttt{pdflatex} passes from
\path{Combinatorial_Transseries_Inverses.tex}; it loads the corpus
notation file \path{docs/fabius-notation.tex} by relative path and has no
figure or bibliography-database dependency.  The retained verification
programs live under \path{verification/}: \path{source1/} holds the first
source's \texttt{verify.py}, \texttt{audit.py}, \texttt{coefficients.wl} and
their recorded \texttt{verification\_results.json} and
\texttt{audit\_results.json}; \path{source2/} holds the second source's
\texttt{verify.py}, \texttt{audit\_symbolic.py}, \texttt{coefficient\_tools.wl}
and the recorded \texttt{verification\_results.json},
\texttt{verification\_report.txt}, \texttt{symbolic\_audit\_report.txt};
\path{source3/} holds the third source's \texttt{verify.py},
\texttt{reference\_implementation.wl} and \texttt{verification\_report.json}.
Each needs Python 3.10 or later with SymPy and \texttt{mpmath}
(\texttt{requirements.txt} in each directory) and no network access; run
\texttt{python verify.py} (and \texttt{python audit.py} or
\texttt{python audit\_symbolic.py}) inside the directory.  The programs were
not modified by the consolidation; their recorded outputs are the runs the
sources shipped.

% ---- hand chunk 11-notation.tex ----

\chapter{Notation, and the conventions of the three sources}
\label{ct:app:notation}

The mathematical notation of the canonical corpus is loaded from
\path{docs/fabius-notation.tex}; the symbols below are document-local, and
where a symbol had different meanings in the three absorbed sources the
convention adopted here is stated together with the sources' spellings.

\begin{center}
\small
\renewcommand{\arraystretch}{1.15}
\begin{tabularx}{\textwidth}{@{}p{0.27\textwidth}X@{}}
\toprule
Symbol & Meaning (and the sources' variants)\\
\midrule
$N_a(y)$, $\lceil F^{-1}(y)\rceil$ & threshold inverse and the ceiling rule, \cref{t2:eq:staircase,t2:eq:ceil}; residue-class rule \cref{t3:eq:residue-rounding}\\
$\cE_n(q)=\ExpC_n(q)$, $\cL_n(a)=\LogC_n(a)$ & the finite exponential and logarithmic partition formulas \cref{t2:eq:E,t2:eq:L}\\
$\sone{r}{m}$ & signed Stirling numbers of the first kind, \cref{t1:eq:stirlingsign}; $\StirI nk=c(n,k)$ unsigned, $\StirII nk=\Sii nk$ second kind\\
$R_n(a,b)=G_n(a,b)$ & gamma-ratio coefficients, \cref{t2:eq:ratio}; the first source's $G_n$ is the second's $R_n$\\
$\fall{z}{r}$, $\rising{z}{r}$ & falling and rising factorials\\
$\coefA{t^n}$, $\coefB{t}{n}$, $[t^n]$ & coefficient extraction (three spellings, one meaning)\\
$X$, $s$ & the exact dominant inverse (core); the third source writes $s$ for the centre of a quadratic phase\\
$D=\dd/\dd L$ (engines); $D=P'(s)$ ($q$-products) & differentiation in the target logarithm; the slope of a quadratic phase at its centre\\
$\Psi(t,u)$, $d_n$ & the balanced-core fixed-point series and its coefficients, \cref{t2:eq:balanced-d}; the first source wrote $c_n$\\
$q_j$ (gamma), $b_m$ ($q$-products) & logarithmic correction coefficients; the first source wrote $d_m$ and $h_r$\\
$A_\pm$, $M_\pm$, $F_\alpha$ & exact positive and negative endpoint blocks and the cosine interpolation; the first source's $J_\pm$ for involutions are the $A_\pm$ here\\
$B(t)$, $b_m$, $H(t)=\log B$, $h_m$ & involution saddle amplitude and its logarithm; the first source wrote $A(t)$, $a_m$\\
$X$, $\ell=\lambda=\log X$, $v_j=z_j$ & involution core, its logarithm, and the half-power inverse coefficients\\
$\mathcal G(x)$, $Z_0,Z_1$, $E_{\mathrm e},E_{\mathrm o}$, $E(x)$ & zigzag gamma block, its parity sheets, and the full cosine interpolation; $a=\pi/2$ there, and $\bar a=\log(2/\pi)$, $\varsigma=\log X+\bar a$ in the uncentred inverse\\
$N_q$, $L_q=I_q$, $M_q$ & all necklaces, Lyndon words (equivalently irreducible polynomials), all necklaces again in the third source's letter\\
$q>1$, $Q=q^{-1}$, $h=\log q$, $\Pinf$, $\A(t)$, $\pp{n}$ & the fixed-base notation \cref{t3:eq:qnotation,t3:eq:pinf}; $\qbin{n}{k}{q}=\qbinq{n}{k}$\\
$\T^{(r)}_\eps$ & root-lattice theta constants \cref{t3:eq:root-theta}\\
\bottomrule
\end{tabularx}
\end{center}

The first source's Motzkin exponent $p=\alpha+1$ and endpoint amplitude
$U_n(\alpha,c)$ specialize to the second source's $\nu=\tfrac32$ and
$a^\pm_m$ at $\alpha=\tfrac12$ (\cref{ct:sec:motzkin-extra}).  The
involution cores of the two sources are the same function:
$L=2\log y+\tfrac12+\log2$ with $X(\log X-1)=L$ in the first, and
$L'=\log(y\sqrt2)+\tfrac14$ with $X(\log X-1)/2=L'$ in the second, so
$L=2L'$.  The zigzag gamma block is written $\mathcal G(x)=2\Gamma(x+1)(\pi/2)^{-x-1}
=(4/\pi)(2/\pi)^x\Gamma(x+1)$ in both.  The Stirling-column inverse
coefficient is written as a falling factorial by the first source and as a
generalized binomial by the second; \cref{ct:rmk:falling} identifies them.
The central Gaussian binomial logarithmic coefficients $h_r$ of the first
source are the $b_m$ of \cref{t3:eq:central-B}.

% ---- hand chunk 12-provenance.tex ----

\chapter{Provenance of the consolidation}
\label{ct:app:provenance}

\section{What was merged}

Three independently written articles arrived on 2026-09-04 and were filed
the same day, as separate members beside the canonical volume, in
\path{drafts/series-and-transseries/}.  Each stated that it extended the
canonical volume's methods to complementary combinatorial families rather
than restating them.  The first two overlapped on most of their families and
on all of their machinery; the third overlapped with them only on central
Gaussian binomial coefficients and on necklaces.  They were merged
editorially into this volume on 2026-09-04: the foundations and the three
inversion engines are stated once (\cref{ct:part:foundations}); each family
appears in one chapter; where two sources proved the same result their
formulas were compared symbol by symbol before one statement was kept and
the other's distinct contributions were absorbed; and the three
verification programs are retained with their recorded outputs.  The three
directories and their arrival PDFs were deleted after a residue audit of
every titled result; git history is the archive.

\section{The absorbed sources and what each contributed}

\begin{description}[leftmargin=1.5em,style=nextline]
\item[Source 1 (labels \texttt{t1:}): \emph{Combinatorial Transseries and
Their Inverses: Further Families}] (\path{combinatorial_transseries_and_inverses/},
1{,}519 lines, 29 pp., 9 theorems).  The phase-coordinate inversion theorem
and the all-multiplicative-sector formula with signed Stirling numbers
(\cref{t1:thm:phase,t1:thm:multi}); the marker-convergence remark; the
flattened linear--logarithmic inverse and the action-collision rule; the
Binet completion and its error bound; the Eulerian columns
(\cref{t1:prop:eulerian}); the general-$\alpha$ endpoint treatment of
Motzkin and trinomial numbers with the convergent beta representation, the
remainder justification and the warning about algebraic truncations; the
Gaussian-moment saddle formula for involutions (\cref{t1:thm:ivforward}),
the scaled inverse with the fixed-point extraction
(\cref{t1:thm:ivinverse}) and the drift-scale rule; the parity-resolved
zigzag sheets, the centred Lambert carrier with odd-only corrections, the
ordered-factorization sector theorem (\cref{t1:thm:zigzaginverse}) and the
two-parity threshold; the operator form of the central Gaussian-binomial
inverse (\cref{t1:thm:qbininverse}) and the $q\downarrow1$ warning; the
general signed-support principle; the certificate with an approximate
phase; the involution coefficient audit against the published $n^{-3}$
value; the first numerical tables and Wolfram recipes.
\item[Source 2 (labels \texttt{t2:}): \emph{Combinatorial Transseries and
Their Inverses}] (\path{combinatorial_transseries_extension/}, 2{,}241
lines, 42 pp., 10 theorems).  The four-object definition and the
interpolation gauge-freedom proposition; the partition formulas and the
gamma-ratio lemma; the scaled reversion; the general (unbalanced)
gamma-quotient master expansion and the balanced inverse theorem with its
analytic justification; Catalan, Fuss--Catalan, multinomial and tableaux
sections; the convergent multiexponential inverse
(\cref{t2:thm:finite-exp}); the fixed-cycle Stirling numbers of the first
kind and their nested-logarithmic inverse; the endpoint lemma; Delannoy and
Schr\"oder numbers and the boundary on the endpoint interpretation; the
phase-Taylor saddle formula for involutions
(\cref{t2:thm:involution-coeff}) and the remark on the action of the exact
inverse; the residue proof of the zigzag Dirichlet decomposition, the full
cosine interpolation, the uncentred core and the prime Euler products; the
connected-graph theorem and its range inverse; the necklace obstruction and
the trigonometric finite models; the harmonic and generalized-harmonic
inverses; the support-closure proposition, the status taxonomy, the
workflow and the research directions; the coefficient table; the second
numerical tables and Wolfram recipes.
\item[Source 3 (labels \texttt{t3:}): \emph{Combinatorial Transseries and
Their Inverses: $q$-Products, Finite Fields, Theta Sectors, and Arithmetic
Sheets}] (\path{combinatorial_transseries_extension-2/}, 1{,}970 lines,
39 pp., 6 theorems).  The multiplicative flat-ambiguity proposition and the
residue-class rounding rule; the analytic tail of a finite $q$-product; the
convergent inverse transseries theorem with the finite nonrecursive
coefficient formula, the computation recurrence and the general-phase
derivative formula (\cref{t3:thm:inverse,t3:thm:coefficient}); general
linear groups, flags, $q$-factorials at both bases, fixed-rank and scaled
Gaussian coefficients, partial flags, binomial $q$-Catalan numbers,
symplectic and unitary orders with their parity sheets; the exact theta-tail
representation of generalized Galois numbers and its inverse
(\cref{t3:thm:theta-tail}); the logarithmic endpoint inverse for finite
limits; the fixed-radical sheets for irreducible polynomials, the
prime-power and $R=6$ examples, the uniform deficit bound and the
two-candidate threshold theorem (\cref{t3:thm:two-candidate}); the singular
transition $q\to1$ with its dilogarithmic phase and modular factor; the
certification and computation sections; the third numerical tables and
Wolfram recipes.
\end{description}

\section{Conventions reconciled and results compared}

Where the sources proved the same theorem, the following were checked to
agree exactly before one statement was kept: the Catalan inverse
coefficients through $X^{-3}$; the Fuss--Catalan constants and $q_1$; the
Stirling-column inverse coefficients (falling-factorial versus
generalized-binomial form, \cref{ct:rmk:falling}) and the $k=3$ second-order
terms; the Motzkin block coefficients through $m=3$ and the first inverse
displacement; the involution amplitude coefficients through $t^5$ and the
inverse coefficients $v_0,v_1,v_2$ against $z_0,z_1,z_2$
(\cref{ct:rmk:zj}); the zigzag gamma block and its first parity sectors;
the central Gaussian-binomial logarithmic coefficients and inverse
coefficients $d_1,d_2$ (\cref{ct:sec:gauss-t1}).  No discrepancy was found.
The notation was reconciled as recorded in \cref{ct:app:notation}: the
involution blocks are $A_\pm$ and the amplitude $B(t)$; the zigzag gamma
block is $\mathcal G$; the first source's differential-operator letters and
the third's centre letters are kept in their own chapters.  The first
source's Taylor-integral proof of the phase-coordinate theorem is the one
printed; the second's transport theorem is stated as
\cref{ct:rmk:transport}.  No mathematical claim of any source was found to
be false.

\section{Status of the claims}

Every theorem, proposition, lemma and corollary carries a proof.  The
numerical tables are the recorded runs of the three retained programs,
which the consolidation did not rerun.  The results are human proofs: none
of the new theorems has a Lean counterpart, and the canonical volume's
register \cite{tai} lists which pieces of the shared apparatus --- the
residual bracket with a root-existence certificate, the Lagrange--B\"urmann
coefficient formulas, the Bell-polynomial calculus, the Wright omega and
Lambert expansions --- are formalized there.  Convergent identities,
Poincar\'e expansions with remainders, exactly normalized sectors, formal
range inverses and finite arithmetic cutoffs are distinguished as in
\cref{ct:ch:synthesis}; no complete complex Stokes classification is claimed.

% ---- hand chunk 13-bibliography.tex ----

\backmatter
\begin{thebibliography}{99}
\small
\setlength{\itemsep}{0.25em}

\bibitem{tai}
V. Reshetnikov, \emph{Transseries: the polynomial--logarithmic calculus,
series reversal at infinity, and the inversion of rapidly growing
functions}, canonical consolidated volume of
\path{drafts/series-and-transseries/Transseries_And_Inversion/} in the
\texttt{ProveIt} repository, consulted 4 September 2026.
\url{https://github.com/VladimirReshetnikov/ProveIt/tree/main/Analysis/FabiusFunction/docs/semi-formalized-research-frontiers/drafts/series-and-transseries/Transseries_And_Inversion}

\bibitem{dlmfgamma}
NIST Digital Library of Mathematical Functions, \S5.11, \emph{Asymptotic Expansions} (gamma and digamma functions, shifted Stirling expansions, gamma ratios, positive-axis error bounds). \url{https://dlmf.nist.gov/5.11}

\bibitem{dlmfbinet}
NIST Digital Library of Mathematical Functions, \S5.9, \emph{Integral Representations} (Binet formulas). \url{https://dlmf.nist.gov/5.9}

\bibitem{dlmfbernoulli}
NIST Digital Library of Mathematical Functions, \S24.2, \emph{Definitions and Generating Functions} (Bernoulli and Euler polynomials). \url{https://dlmf.nist.gov/24.2}

\bibitem{dlmflaplace}
NIST Digital Library of Mathematical Functions, \S2.3, \emph{Integrals of a Real Variable} (Watson's lemma, Laplace's method). \url{https://dlmf.nist.gov/2.3}

\bibitem{dlmfstirling}
NIST Digital Library of Mathematical Functions, \S26.8, \emph{Set Partitions: Stirling Numbers}. \url{https://dlmf.nist.gov/26.8}

\bibitem{dlmflegendre}
NIST Digital Library of Mathematical Functions, \S14.12, \emph{Integral Representations} (Legendre functions). \url{https://dlmf.nist.gov/14.12}

\bibitem{dlmfhurwitz}
NIST Digital Library of Mathematical Functions, \S25.11, \emph{Hurwitz Zeta Function}. \url{https://dlmf.nist.gov/25.11}

\bibitem{dlmfq}
NIST Digital Library of Mathematical Functions, \S17.2, \emph{Calculus} ($q$-Pochhammer symbols, Gaussian coefficients, the $q$-binomial theorem; equations 17.2.1, 17.2.3, 17.2.27, 17.2.35--17.2.38). \url{https://dlmf.nist.gov/17.2}

\bibitem{dlmfqmod}
NIST Digital Library of Mathematical Functions, equation 17.2.6\_1 (modular transformation of the Euler product) and \S23.18, \emph{Modular Transformations}. \url{https://dlmf.nist.gov/17.2.E6_1}

\bibitem{dlmfW}
NIST Digital Library of Mathematical Functions, \S4.13, \emph{Lambert $W$-Function}. \url{https://dlmf.nist.gov/4.13}

\bibitem{oeiscatalan}
OEIS Foundation, \emph{A000108: Catalan numbers}. \url{https://oeis.org/A000108}

\bibitem{oeisfuss}
OEIS Foundation, \emph{A001764: $\binom{3n}{n}/(2n+1)$, ternary trees}. \url{https://oeis.org/A001764}

\bibitem{oeistableaux}
OEIS Foundation, \emph{A005789: standard Young tableaux of shape $(n,n,n)$}. \url{https://oeis.org/A005789}

\bibitem{oeismotzkin}
OEIS Foundation, \emph{A001006: Motzkin numbers}. \url{https://oeis.org/A001006}

\bibitem{oeistrinomial}
OEIS Foundation, \emph{A002426: central trinomial coefficients}. \url{https://oeis.org/A002426}

\bibitem{oeisdelannoy}
OEIS Foundation, \emph{A001850: central Delannoy numbers}. \url{https://oeis.org/A001850}

\bibitem{oeisschroder}
OEIS Foundation, \emph{A006318: large Schr\"oder numbers}. \url{https://oeis.org/A006318}

\bibitem{oeisinvolution}
OEIS Foundation, \emph{A000085: number of involutions}. \url{https://oeis.org/A000085}

\bibitem{oeiszigzag}
OEIS Foundation, \emph{A000111: Euler or up/down numbers}. \url{https://oeis.org/A000111}

\bibitem{oeiseulerian}
OEIS Foundation, \emph{A008292: triangle of Eulerian numbers}. \url{https://oeis.org/A008292}

\bibitem{oeisallgraphs}
OEIS Foundation, \emph{A006125: $2^{\binom n2}$, labeled graphs}. \url{https://oeis.org/A006125}

\bibitem{oeisconnected}
OEIS Foundation, \emph{A001187: connected labeled graphs}. \url{https://oeis.org/A001187}

\bibitem{oeisnecklace}
OEIS Foundation, \emph{A000031: binary necklaces}. \url{https://oeis.org/A000031}

\bibitem{oeislyndon}
OEIS Foundation, \emph{A001037: primitive binary necklaces, binary Lyndon words, irreducible polynomials over $\F_2$}. \url{https://oeis.org/A001037}

\bibitem{oeisGL}
OEIS Foundation, \emph{A002884: order of the general linear group $\GL(n,2)$}. \url{https://oeis.org/A002884}

\bibitem{oeisflags}
OEIS Foundation, \emph{A005329: products $\prod_{j=1}^n(2^j-1)$}. \url{https://oeis.org/A005329}

\bibitem{oeiscentral}
OEIS Foundation, \emph{A006098: central Gaussian binomial coefficients at $q=2$}. \url{https://oeis.org/A006098}

\bibitem{oeisqcatalan}
OEIS Foundation, \emph{A015030: binomial-version $q$-Catalan numbers at $q=2$}; the asymptotic entry by Vladimir Reshetnikov is dated 26 September 2021. \url{https://oeis.org/A015030}

\bibitem{oeisSp}
OEIS Foundation, \emph{A003923: orders of $\Sp(2n,2)$}. \url{https://oeis.org/A003923}

\bibitem{oeisU}
OEIS Foundation, \emph{A070953: orders of the general unitary groups over $\F_4$, convention $\mathrm{GU}(n,2)$}. \url{https://oeis.org/A070953}

\bibitem{oeisGalois}
OEIS Foundation, \emph{A006116: numbers of subspaces of $\F_2^n$}. \url{https://oeis.org/A006116}

\bibitem{wlstirling2}
Wolfram Research, \emph{StirlingS2}, Wolfram Language Documentation. \url{https://reference.wolfram.com/language/ref/StirlingS2.html}

\bibitem{wlstirling1}
Wolfram Research, \emph{StirlingS1}, Wolfram Language Documentation (signed first-kind convention). \url{https://reference.wolfram.com/language/ref/StirlingS1.html}

\bibitem{wlqbinomial}
Wolfram Research, \emph{QBinomial}, Wolfram Language Documentation. \url{https://reference.wolfram.com/language/ref/QBinomial.html}

\bibitem{wlqpochhammer}
Wolfram Research, \emph{QPochhammer}, Wolfram Language Documentation. \url{https://reference.wolfram.com/language/ref/QPochhammer.html}

\bibitem{wlbeta}
Wolfram Research, \emph{DirichletBeta}, Wolfram Language Documentation. \url{https://reference.wolfram.com/language/ref/DirichletBeta.html}

\bibitem{wlproductlog}
Wolfram Research, \emph{ProductLog}, Wolfram Language Documentation. \url{https://reference.wolfram.com/language/ref/ProductLog.html}

\bibitem{wlcatalan}
Wolfram Research, \emph{CatalanNumber}, Wolfram Language Documentation. \url{https://reference.wolfram.com/language/ref/CatalanNumber.html}

\bibitem{wlharmonic}
Wolfram Research, \emph{HarmonicNumber}, Wolfram Language Documentation. \url{https://reference.wolfram.com/language/ref/HarmonicNumber.html}

\bibitem{ekz}
S. B. Ekhad, M. Kauers, and D. Zeilberger, \emph{The $O(1/n^{85})$ asymptotic expansion of OEIS sequence A85}, arXiv:2410.16334, 2024; the HTML rendering of version 1 was consulted for the audit of \cref{ct:app:audit}. \url{https://arxiv.org/abs/2410.16334}

\bibitem{kousidis}
S. Kousidis, \emph{Asymptotics of generalized Galois numbers via affine Kac--Moody algebras}, 2011, arXiv:1109.2546. \url{https://arxiv.org/abs/1109.2546}

\bibitem{bliem}
T. Bliem and S. Kousidis, \emph{The number of flags in finite vector spaces: asymptotic normality and Mahonian statistics}, 2011, arXiv:1109.4624. \url{https://arxiv.org/abs/1109.4624}

\end{thebibliography}
\end{document}

